% !TeX root = ../Book2.tex
% 纲要标题：# 第二编　多线性代数与形式
% 纲要位置：工作记录/计划纲要.md，第 1303 行。

\BookPart{多线性代数与形式}{b2:part:02}
\BookPartGuide
两个向量可以相加，也可以作为两个独立变量放入同一个公式。内积、行列式及向量的外积都与后一种运算有关，但它们满足的交换规则并不相同。本编以张量积为基础，分别把多线性、对称与交替这些条件写成代数构造，再研究双线性型和二次型的分类。

\ProseParagraphBreak
\chapref{b2:ch:08}将两个因子的张量积推广到多个因子。把所有次数放在一起便得到张量代数；加入交换关系则得到对称代数。有限自由模的对称代数可识别为多项式环，但这个识别用到了基。需要特别区分对称幂这一商模与张量幂中的不变部分：在次数 \(d\ge1\) 中，若 \(d!\) 在底环 \(R\) 中可逆，平均对称化给出两者的自然同构；它不可逆时，不能一般地使用除以 \(d!\) 的平均公式。特征零也不保证这一条件，例如 \(2\) 在整数环中不可逆。

\ProseParagraphBreak
\chapref{b2:ch:09}从“两个输入相同则值为零”的交替性出发构造外幂，并以楔积组成外代数。这项条件在特征二时仍有意义，不能只改写为交换输入后变号。在有限自由模上，最高外幂把行列式解释为一个秩一模上的作用，各阶外幂则记录子式，行列式乘法性和 Cauchy--Binet 公式也随之得到解释。随后讨论伴随矩阵与 Cayley--Hamilton 定理。章末的 Plücker 坐标\footnote{普吕克（Julius Plücker，1801-1868），德国数学家及物理学家，以齐次坐标研究射影几何。}和 Koszul 复形\footnote{科斯居尔（Jean-Louis Koszul，1921-2018），法国数学家，研究微分几何与 Lie 群；以其名字命名的复形也用于交换代数。}是外代数的两种应用：前者用中间外幂记录子空间，后者用收缩算子构造次数 \(-1\) 的微分，本章即给出 Koszul 复形的初步构造；一般的正合性问题将在第三卷结合正则序列继续研究。

\ProseParagraphBreak
\chapref{b2:ch:10}研究给定的配对，而不只是收集全部多线性映射。换基对双线性型产生合同变换，对线性算子则产生相似变换，这两类分类问题须分开。特征不为二时的对角化、实数域上的惯性及 Witt 理论\footnote{维特（Ernst Witt，1911-1991），德国数学家，研究二次型、域与代数的结构。}依次说明，哪些数据能够从改变坐标后的型中恢复。对域上对合的考虑进一步引出 Hermitian 型\footnote{埃尔米特（Charles Hermite，1822-1901），法国数学家，研究数论与代数，并证明 \(e\) 的超越性。}与酉群；章末的 Clifford 代数\footnote{克利福德（William Kingdon Clifford，1845-1879），英国数学家，研究非欧几何，并引入推广外代数的 Clifford 代数。}联系二次型与按关系定义的代数。

\begin{center}
\begin{tikzcd}[column sep=large,row sep=large]
\begin{gathered}\text{第6章}\\\text{张量积与双模}\end{gathered}\arrow[r]
&\begin{gathered}\text{第8章}\\\text{张量代数与对称代数}\end{gathered}\arrow[r]\arrow[d,dashed]
&\begin{gathered}\text{第9章}\\\text{外代数与行列式}\end{gathered}\\
\begin{gathered}\text{第1章}\\\text{向量空间与对偶}\end{gathered}\arrow[r]
&\begin{gathered}\text{第10章}\\\text{双线性型与二次型}\end{gathered}
&
\end{tikzcd}
\end{center}

实线表示主要的直接依赖，虚线专指 Clifford 代数所用的张量代数构造。第8、9章还使用第一编中自由模、商模和泛性质的基本语言；第10章关于型的分类以有限维向量空间为基础。

\ProseParagraphBreak
本编默认的系数范围随问题变化：有限自由模上的张量、对称与外代数可在交换环上讨论；非退化型、Witt 分解及经典群则主要在域上讨论。每到一个结论，先辨认底环或基域、特征与维数假设，再检查所用映射。对一条抽象恒等式，可以先在二维或三维基上计算；对一个分类结论，则应尝试改变特征或去掉非退化条件，看看证明在哪一步不再成立。第五卷第7.3节的 Schur--Weyl 理论\footnote{舒尔（Issai Schur，1875-1941），俄国出生的德国数学家，研究群表示和矩阵代数。外尔（Hermann Weyl，1885-1955），德国裔美国数学家，研究表示论、几何与数学物理。}和第三卷第15.2--15.3节的 Koszul 复形及其正合性将进一步使用这些构造。

% !TeX root = ../../Book2.tex
% 纲要标题：## 第8章　张量代数与对称代数

\chapter{张量代数与对称代数}
\label{b2:ch:08}
\BookChapterNumber{b2:ch:08}

\begin{chapterabstract}
多重张量积把多线性映射改写成线性映射，张量代数把所有次数的有序乘积放在同一分次代数中，对称代数则进一步施加交换关系。本章证明这些构造的泛性质，建立自由张量积的基与有限自由情形的多项式解释，并比较对称商与不变张量，说明阶乘可逆时的同构及这一条件缺失时可能出现的差别。最后把对偶、张量积与对称幂用于群表示，说明等变映射如何成为不变张量，以对偶空间的对称代数定义多项式不变环，并比较位置置换与一般线性群作用的对称算子。
\end{chapterabstract}

\begin{learninggoals}
\begin{itemize}
\item 用多重张量积表示多线性映射，说明改变括号与置换因子的指定同构及其相容性。
\item 构造分次张量代数，区分一般模的张量代数与选定自由生成元后的字代数，利用生成元和关系定义同态。
\item 证明对称代数的泛性质及多项式解释，计算对称幂的基和维数，并检查阶乘是否可逆。
\item 写出张量、对偶、对称幂表示的作用公式，区分位置置换不变与群作用不变，将有限维 Hom 识别为张量表示。
\end{itemize}
\end{learninggoals}

取自由模 \(V=Re_1\oplus Re_2\)。在张量代数中，\(e_1e_2\) 与 \(e_2e_1\) 记录两种不同次序；若希望它们相同，就必须再施加交换关系，得到对称代数中的同一个乘积。这不是改变记号而已：给两个元素指定像时，映到一般结合代数无需它们可交换，映到交换代数才自动满足新关系。不同的关系决定了构造可以接受哪些映射。

\ProseParagraphBreak

这里还要区分两种构造：一种选出在位置置换下已经固定的张量，另一种把不同排列识别为同一个类。它们将在\secref{b2:sec:0803}中分别定义；次数二的计算先说明为什么不能无条件混同。即使取 \(R=\mathbb Z\)，张量 \(e_1\otimes e_2+e_2\otimes e_1\) 在交换商中的像仍是 \(2e_1e_2\)，不能由它得到 \(e_1e_2\)。所以在二不可逆时，不能只凭对称性就断言不变张量与对称商相同；在一般次数中，阶乘可逆将提供统一的平均方法；不可逆时，则须直接检查从固定子模到商的映射。随后还要区分位置置换与群在各因子上同时作用，这两种作用所固定的张量也未必相同。

\ProseParagraphBreak

本章以第六章的张量积泛性质、剩余作用和结合同构为直接基础；自然性采用第五章的定义，有限自由模的坐标与对偶采用第一、三章的结果。除表示部分把底环取为域外，前面大部分构造只需交换底环。各类商构造的泛性质都要同时说明所除去的关系，以及从商出发的映射为何零化这些关系。

\ProseParagraphBreak
张量代数与对称代数的商构造可参阅 Etingof 等人的《Introduction to Representation Theory》\cite[Section 2.12, Definition 2.12.1, p. 35]{EtingofEtAl2011RepresentationTheory}；对偶与张量积表示的公式见\cite[Section 4.4, p. 65]{EtingofEtAl2011RepresentationTheory}。对称幂的商定义适用于一般交换环；用平均算子识别不变张量时，相关阶乘须可逆。

% !TeX root = ../../Book2.tex
% 纲要标题：### 8.1 多线性映射

\section{多线性映射}
\label{b2:sec:0801}

从双线性映射走向多线性映射，首先遇到的不是新运算，而是括号和标量的安排。一个三变量公式可以先合并前两个变量，也可以先合并后两个变量；只有证明两种构造通过指定同构相容，才可以省略括号。本节固定交换环 \(R\)，所有模均为幺性 \(R\) 模。这一交换性约定允许在不同因子间移动标量；第六章一般双模的作用侧别仍是它的基础。

% 纲要要点（供目录核对）：
%
% * 多重张量积
% * 自然同构与结合约定
% * 自由模张量积的基与有限秩情形的秩公式
%

\subsection{多重张量积}
\label{b2:subsec:0801:01}

\begin{definition}\label{b2:def:multilinear-map}
设 \(R\) 为交换环，\(M_1,\ldots,M_d,N\) 为 \(R\) 模，\(d\ge1\)。如果映射
\[
b:M_1\times\cdots\times M_d\longrightarrow N
\]
在固定其余变量后对每个变量都是 \(R\) 线性的，就称 \(b\) 为 \(R\) 上的\term[duoxianxing-yingshe]{多线性映射}[multilinear map]。具体说，对每个 \(1\le j\le d\)，任取 \(r,s\in R\)、\(x,y\in M_j\) 及其余各位置的 \(m_i\in M_i\)，都有
\[
\begin{aligned}
b(m_1,\ldots,rx+sy,\ldots,m_d)
&=r\cdot b(m_1,\ldots,x,\ldots,m_d)\\
&\quad+s\cdot b(m_1,\ldots,y,\ldots,m_d).
\end{aligned}
\]
\end{definition}

直积 \(M_1\times\cdots\times M_d\) 本身是逐坐标运算的 \(R\) 模。当 \(d=1\) 时，多线性就是线性；当 \(d\ge2\) 时，对每个变量分别线性与对这个直积模整体线性是不同的要求。多线性映射满足 \(b(rm_1,\ldots,rm_d)=r^d\cdot b(m_1,\ldots,m_d)\)，整体线性却要求右端只有一个标量 \(r\)。例如 \(R=\mathbb R\)，乘积 \((x,y,z)\mapsto xyz\) 对每个变量分别线性，但把三个变量同时乘二会使值乘八，而非乘二。

\ProseParagraphBreak
令 \(T_1=M_1\)，递归定义 \(T_d=T_{d-1}\otimes_RM_d\)。因 \(R\) 交换，第六章的剩余作用使每个 \(T_j\) 为 \(R\) 模，因而可以继续取张量积。把 \((m_1,\ldots,m_d)\) 映到
\[
\Bigl(\cdots\bigl((m_1\otimes m_2)\otimes m_3\bigr)\cdots\Bigr)\otimes m_d
\]
的映射记作 \(\tau\)。以后用 \(m_1\otimes\cdots\otimes m_d\) 表示这个元素，用 \(\bigotimes_{j=1}^d M_j\) 表示 \(T_d\)。左结合的递归构造先给出了一个具体模型；下面的多线性泛性质将刻画这个对象，使后续构造不再依赖选用哪一种括号形式。

\begin{theorem}[多重张量积的泛性质]\label{b2:thm:multiple-tensor-universal}
设 \(R\) 为交换环，\(d\ge1\)，\(M_1,\ldots,M_d\) 为 \(R\) 模。令 \(\tau:M_1\times\cdots\times M_d\to\bigotimes_{j=1}^dM_j\) 把 \((m_1,\ldots,m_d)\) 送到纯张量 \(m_1\otimes\cdots\otimes m_d\)，则 \(\tau\) 是多线性映射。对任意 \(R\) 模 \(N\) 与多线性映射 \(b:M_1\times\cdots\times M_d\rightarrow N\)，存在唯一 \(R\) 线性映射
\[
\widetilde b:\bigotimes_{j=1}^dM_j\longrightarrow N,
\qquad \widetilde b(m_1\otimes\cdots\otimes m_d)=b(m_1,\ldots,m_d).
\]
多重张量积由这些纯张量生成，并在保持指定多线性映射的唯一同构意义下唯一。
\end{theorem}
\begin{proof}
三变量时，给定三线性映射 \(b:M_1\times M_2\times M_3\rightarrow N\)，固定 \(z\in M_3\) 后，二重张量积的泛性质给出唯一线性映射
\[
b_z:M_1\otimes_RM_2\longrightarrow N,
\qquad b_z(x\otimes y)=b(x,y,z).
\]
对 \(r,s\in R\) 及 \(z,w\in M_3\)，映射 \(b_{rz+sw}\) 与 \(r b_z+s b_w\) 在每个 \(x\otimes y\) 上取值相同，故由纯张量生成性相等。因此 \((t,z)\mapsto b_z(t)\) 是 \((M_1\otimes_RM_2)\times M_3\) 上的双线性映射，再应用一次张量积的泛性质，就得到 \((M_1\otimes_RM_2)\otimes_RM_3\rightarrow N\)。这样先合并前两个变量，再合并最后一个变量，并没有改变原来的三线性取值。

对一般 \(d\) 按同一方法归纳。\(d=1\) 时是恒等映射，\(d=2\) 时是\cref{b2:thm:tensor-universal}及\cref{b2:prop:tensor-bimodule}给出的交换环线性版本。设结论对 \(d-1\) 成立，固定 \(m_d\) 后得到唯一线性映射 \(b_{m_d}:T_{d-1}\rightarrow N\)。这些映射随 \(m_d\) 线性变化：所需等式在前 \(d-1\) 重纯张量上由 \(b\) 的多线性成立，再由生成性推广到整个 \(T_{d-1}\)。故 \((t,m_d)\mapsto b_{m_d}(t)\) 为双线性映射。

应用\cref{b2:thm:tensor-universal}及\cref{b2:prop:tensor-bimodule}给出的交换环线性版本，得到所需 \(\widetilde b\)。任一 \(T_d\) 中的元素先写成有限和 \(\sum t_i\otimes m_i\)，再把各 \(t_i\) 展开成前 \(d-1\) 重纯张量，故 \(d\) 重纯张量生成全模。这证明映射唯一，也说明 \(\tau\) 的逐变量线性由每一步的双线性继承。

若另一个对象带有相同泛性质，把两者的指定多线性映射分别代入对方的泛性质，得到正反两个线性映射。两复合在每个指定纯张量上为恒等，故在全模为恒等；保持指定映射的同构也由生成性唯一确定。
\end{proof}

为了统一后面的次数与乘法，采用零重张量积为 \(R\) 的约定，并把没有输入变量的“多线性映射”约定为指定目标模中的一个元素 \(n\)。它对应唯一线性映射 \(R\rightarrow N\)、\(r\mapsto rn\)，其在 \(1\) 处的值为 \(n\)。次数零的这一约定将提供张量代数的单位元。

\subsection{自然同构与结合约定}
\label{b2:subsec:0801:02}

任意完整括号形式都可递归应用二重张量积构造。对因子个数归纳，把最外层写成 \(X\otimes_RY\)，其中 \(X,Y\) 是较少因子的括号积。固定后半组变量，由 \(X\) 的泛性质合并前半组；所得映射对后半组仍逐变量线性，再由 \(Y\) 的泛性质合并，得到 \(X\times Y\) 上的双线性映射。二重张量积的泛性质给出所需线性映射，唯一性则由有序纯张量生成全模得到。因此每种括号形式都满足\cref{b2:thm:multiple-tensor-universal}的同一个 \(d\) 变量泛性质，它们之间有唯一同构，将相同顺序的纯张量送成彼此。\cref{b2:thm:tensor-associativity-unit}的结合同构正是三因子时的这个映射。

\ProseParagraphBreak
起点与终点的括号形式给定后，任何逐步改变括号的路径都保持同一组有序纯张量。由泛性质的唯一性，这些复合映射必定相同；第六章的四因子计算就是这一事实的具体情形。因此下文可以省略括号：无括号的记号表示经这些典范同构识别的张量积，而非宣称不同递归模型字面相等。

\begin{proposition}\label{b2:prop:multiple-tensor-functoriality}
设 \(R\) 为交换环，\(d\ge1\)，\(M_j,N_j\) 为 \(R\) 模。给定 \(R\) 线性映射 \(f_j:M_j\rightarrow N_j\)（\(1\le j\le d\)），存在唯一 \(R\) 线性映射
\[
\bigotimes_{j=1}^d f_j:\bigotimes_{j=1}^d M_j\longrightarrow\bigotimes_{j=1}^d N_j,
\qquad \bigotimes_jm_j\longmapsto\bigotimes_j f_j(m_j).
\]
它保持恒等映射与逐因子复合，并与保持有序纯张量的结合同构相容。若 \(\sigma\in S_d\)，还存在置换因子的同构
\[
P_\sigma:M_1\otimes\cdots\otimes M_d
\longrightarrow M_{\sigma^{-1}(1)}\otimes\cdots\otimes M_{\sigma^{-1}(d)},
\quad \bigotimes_jm_j\longmapsto
m_{\sigma^{-1}(1)}\otimes\cdots\otimes m_{\sigma^{-1}(d)}.
\]
当所有因子相同时，有 \(P_\sigma P_\pi=P_{\sigma\pi}\)。
\end{proposition}
\begin{proof}
所给公式在原变量中多线性，所以由多重张量积的泛性质诱导线性映射。对逐因子复合，两个可能映射都把 \(\bigotimes_jm_j\) 送到 \(\bigotimes_j g_j\bigl(f_j(m_j)\bigr)\)，故相等；恒等映射与结合的相容性由相同的纯张量检验得到。

置换公式仍逐变量线性，因此也给出映射，反向用 \(\sigma^{-1}\) 即得逆。所有因子相同时，连续做 \(P_\pi\)、\(P_\sigma\)，第 \(j\) 个位置最终来自原来的 \(\pi^{-1}\sigma^{-1}(j)=(\sigma\pi)^{-1}(j)\) 位置，所以满足所述乘法公式。
\end{proof}

\symidx{\(P_\sigma\)：置换张量因子的算子}
逆置换使原来第 \(j\) 个位置的元素移动到第 \(\sigma(j)\) 个位置，从而给出左群作用。例如在 \(M^{\otimes3}\) 中，取 \(\sigma=(123)\)、\(\pi=(12)\)，并按先做 \(\pi\) 再做 \(\sigma\) 的次序复合，则
\[
\begin{aligned}
P_\sigma(a\otimes b\otimes c)&=c\otimes a\otimes b,\\
P_\sigma P_\pi(a\otimes b\otimes c)
&=P_\sigma(b\otimes a\otimes c)
=c\otimes b\otimes a.
\end{aligned}
\]
这正是 \(P_{\sigma\pi}=P_{(13)}\) 的取值。若把公式中的 \(\sigma^{-1}\) 换成 \(\sigma\)，复合反而会对应 \(\pi\sigma\)，次序被倒转。置换因子也没有交换代数中两个元素的含义：\(M\otimes N\cong N\otimes M\) 是指定同构，通常不能在一个尚未给乘法的模中把 \(m,n\) 相乘。

\subsection{有限自由情形的基与维数}
\label{b2:subsec:0801:03}

纯基张量的基定理不要求各因子的基有限；有限性只在把所得基的大小写成有限秩或维数公式时需要。

\begin{proposition}\label{b2:prop:multiple-tensor-basis}
设 \(R\) 为交换环，\(d\ge1\)，各 \(M_j\) 为自由 \(R\) 模，具有指定基 \((e_{j,i})_{i\in I_j}\)，则
\[
e_{1,i_1}\otimes\cdots\otimes e_{d,i_d},
\qquad (i_1,\ldots,i_d)\in I_1\times\cdots\times I_d,
\]
构成多重张量积的一组基。若 \(R\ne0\)，且各 \(I_j\) 有限，记 \(n_j=\lvert I_j\rvert\)，则其秩为 \(n_1\cdots n_d\)；在域上，这就是有限维空间的维数乘法公式。
\end{proposition}
\begin{proof}
逐因子展开 \(m_j=\sum_{i_j\in I_j}a_{j,i_j}e_{j,i_j}\)。每个展开只有有限多个非零系数，而因子的个数 \(d\) 有限，所以多线性给出有限和
\[
m_1\otimes\cdots\otimes m_d
=\sum_{i_1,\ldots,i_d}
a_{1,i_1}\cdots a_{d,i_d}\,
e_{1,i_1}\otimes\cdots\otimes e_{d,i_d}.
\]
纯张量生成全模，故所列元素生成。对每组下标 \((i_1,\ldots,i_d)\)，将各因子的第 \(i_j\) 个坐标相乘，得到多线性映射 \((m_1,\ldots,m_d)\mapsto a_{1,i_1}\cdots a_{d,i_d}\)。它诱导的线性泛函在对应的纯基张量上取一，在其他纯基张量上取零。对任意有限线性关系应用相应的坐标泛函，逐项得到所有系数为零，故线性无关。基有限时共有 \(n_1\cdots n_d\) 个纯基张量；秩的含义由\cref{b2:thm:free-rank-invariance}保证。
\end{proof}

例如 \(R^2\otimes_RR^3\otimes_RR^2\) 有十二个标准纯基张量。在 \(R=\mathbb Z\) 时，
\[
(e_1+2e_2)\otimes(f_1-f_3)\otimes e_2
=e_1\otimes f_1\otimes e_2-e_1\otimes f_3\otimes e_2
+2e_2\otimes f_1\otimes e_2-2e_2\otimes f_3\otimes e_2.
\]
这些系数是唯一的，但将一个张量写成几个任意纯张量之和并不唯一。\subsecref{b2:subsec:0601:03}中已经证明，\(k^2\otimes_k k^2\) 中的 \(e_1\otimes e_1+e_2\otimes e_2\) 不能写成一个纯张量。因此，基张量坐标和一般纯张量分解是两种不同的数据。

\ProseParagraphBreak
自由假设不能只替换成“有限生成”。例如 \(M=\mathbb Z/2\mathbb Z\)、\(N=\mathbb Z/3\mathbb Z\) 均可由一个元素生成，\cref{b2:prop:tensor-quotient}却给出 \(M\otimes_{\mathbb Z}N=0\)。没有基时，张量关系可能消去整个模，不能用生成元数相乘代替秩或维数。

% !TeX root = ../../Book2.tex
% 纲要标题：### 8.2 张量代数

\section{张量代数}
\label{b2:sec:0802}

固定一个模 \(M\) 后，可以把不同次数的张量同时放进一个对象：一次张量是 \(M\) 的元素，二次张量记录两个元素的有序相乘，三次张量记录三个元素的有序相乘。把这些次数作直和，再用连接张量串定义乘法，便得到张量代数。这里首先保留次序；交换性将作为下一节另外施加的关系。

% 纲要要点（供目录核对）：
%
% * 分次张量代数
% * 自由结合代数的泛性质
% * 商代数与按关系定义代数
% * 增广理想的一次商与底模恢复
%

\subsection{分次张量代数}
\label{b2:subsec:0802:01}

\begin{definition}\label{b2:def:associative-graded-algebra}
设 \(R\) 为交换环。给定含单位元的环 \(A\) 及保单位元环同态 \(R\rightarrow Z(A)\)，称 \(A\) 为 \(R\) 上的\term[jiehe-daishu]{结合代数}[associative algebra]。它因此是 \(R\) 模，乘法对两个变量分别 \(R\) 线性；结构同态不要求单射。两个 \(R\) 代数之间与各自结构同态相容的保单位元环同态，称为 \(R\) 代数同态。

若另有 \(R\) 模直和分解
\[
A=\bigoplus_{d\ge0}A_d,\qquad 1_A\in A_0,\qquad A_iA_j\subseteq A_{i+j},
\]
则称 \(A\) 为非负整数\term[fenci-daishu]{分次代数}[graded algebra]；\(A_d\) 中的元素称为次数 \(d\) 的\term[qici-yuansu]{齐次元素}[homogeneous element]。保持各次数分量的代数同态称为分次同态。
\end{definition}

结构同态给出标量的作用，并不表示 \(A\) 必须含有一个与 \(R\) 相同的子环。例如对 \(n>1\)，商映射 \(\mathbb Z\rightarrow\mathbb Z/n\mathbb Z\) 使后者成为 \(\mathbb Z\) 代数，而有限环 \(\mathbb Z/n\mathbb Z\) 不可能含有 \(\mathbb Z\) 子环。

\ProseParagraphBreak
直和要求一个元素只有有限多个非零齐次分量，因此不包含任意无穷和。分次也没有规定 \(ab=(-1)^{ij}ba\)；这种带符号的交换条件是另一个额外性质。普通交换的多项式代数以全次数分次，张量代数则一般不交换。

\begin{definition}\label{b2:def:tensor-algebra}
设 \(R\) 为交换环，\(M\) 为 \(R\) 模。令 \(M^{\otimes0}=R\)、\(M^{\otimes1}=M\)，并记
\[
T_R(M)=\bigoplus_{d\ge0}M^{\otimes d}.
\]
对正次数的纯张量规定
\[
(m_1\otimes\cdots\otimes m_i)(n_1\otimes\cdots\otimes n_j)
=m_1\otimes\cdots\otimes m_i\otimes n_1\otimes\cdots\otimes n_j,
\]
对 \(r\in R=M^{\otimes0}\) 与 \(x\in M^{\otimes d}\)（\(d\ge0\)），规定 \(rx=xr=r\cdot x\)，再按双线性扩张到整个直和。所得分次 \(R\) 代数称为 \(M\) 的\term[zhangliang-daishu]{张量代数}[tensor algebra]。
\end{definition}
\symidx{\(M^{\otimes d}\)：模的 \(d\) 次张量幂}
\symidx{\(T_R(M)\)：模的张量代数}

各个 \(M^{\otimes d}\) 原本只是不同次数的模；上述连接把两个张量串变成一个更长的串，才给这些次数之间定义了乘法。其良定义由张量积的泛性质保证。对每对 \(i,j>0\)，\cref{b2:thm:multiple-tensor-universal}与结合同构给出线性映射
\[
M^{\otimes i}\otimes_RM^{\otimes j}\longrightarrow M^{\otimes(i+j)}
\]
并给出上面的纯张量公式，所以它不依赖纯张量和的表达。任意两个元素仅有有限多个分量，将这些双线性乘法逐项相加仍为有限和。三个齐次元素相乘时，两种结合方式都连接同一个有序张量串；由纯张量生成性及分配律，结合律在所有元素上成立。\(1_R\) 是单位，标量与所有张量可交换，次数相加，因此定义确实给出分次 \(R\) 代数。

\ProseParagraphBreak
一次分量的包含记为 \(\iota:M\rightarrow T_R(M)\)，它单射，因为这是一个直和分量的包含。\(T_R(M)\) 由 \(\iota(M)\) 和 \(R\) 作为代数生成：每个纯张量都等于若干一次元素之积，任意元素又是这些纯张量与常数的有限线性组合。

\subsection{自由结合代数的泛性质}
\label{b2:subsec:0802:02}

\begin{theorem}[张量代数的泛性质]\label{b2:thm:tensor-algebra-universal}
设 \(R\) 为交换环，\(M\) 为 \(R\) 模，\(\iota:M\hookrightarrow T_R(M)\) 为张量代数的一次分量包含。对任意含单位元的结合 \(R\) 代数 \(A\) 及 \(R\) 线性映射 \(f:M\rightarrow A\)，存在唯一 \(R\) 代数同态
\[
\widehat f:T_R(M)\longrightarrow A,\qquad \widehat f\iota=f.
\]
其零次部分为 \(A\) 的结构同态，其正次部分满足
\[
\widehat f(m_1\otimes\cdots\otimes m_d)=f(m_1)\cdots f(m_d).
\]
\end{theorem}
\begin{proof}
由于 \(R\) 的像在 \(A\) 中心，右端对 \(m_1,\ldots,m_d\) 的每个变量都 \(R\) 线性。\cref{b2:thm:multiple-tensor-universal}给出唯一线性映射 \(f_d:M^{\otimes d}\rightarrow A\)。令 \(f_0\) 为结构同态，利用直和的泛性质把它们合成 \(\widehat f\)。在两个纯张量上，先连接再作用与分别作用后在 \(A\) 相乘相同；由分配律推广到有限和，故 \(\widehat f\) 保持乘法。它在零次部分保持单位元及标量，所以是 \(R\) 代数同态。

任何所求同态都必须在一次分量上等于 \(f\)，在纯张量上等于这些像的有序乘积，在 \(R\) 上等于结构同态。因此它在生成全代数的元素上已被确定，唯一性成立。
\end{proof}

该泛性质的目标不要求交换，也不要求 \(f\) 单射。张量代数的“自由”是说：\(M\) 中原有的线性关系仍然保留，但不再对它的元素另加乘法关系。因此任意到结合 \(R\) 代数的线性映射都能唯一延拓，而不需要额外检查这些像之间满足何种乘法关系。

\ProseParagraphBreak
将一次元素全部送到零，就得到\term[zengguang-tongtai]{增广同态}[augmentation] \(\varepsilon:T_R(M)\rightarrow R\)。它是取零次分量的投影：若 \(x_d\in M^{\otimes d}\)，则
\[
\varepsilon(r+x_1+\cdots+x_n)=r.
\]
其核为正次部分 \(I=T_R(M)_+=\bigoplus_{d\ge1}M^{\otimes d}\)，称为增广理想。\cref{b2:prop:tensor-augmentation-indecomposables}将说明，模掉其中至少含两个一次因子的乘积后，\(I/I^2\) 恰好恢复一次分量 \(M\)。

\ProseParagraphBreak
把 \(R\) 线性映射 \(f:M\rightarrow N\) 后接 \(N\) 到其张量代数的包含，得到分次同态 \(T_R(f):T_R(M)\rightarrow T_R(N)\)，它在 \(d\) 次分量上就是 \(f^{\otimes d}\)，在零次分量上为 \(R\) 的恒等映射。恒等及复合在一次生成元上核对，便知 \(T_R\) 是函子。

\ProseParagraphBreak
这个函子却不一定保持单射。第七章用来检验平坦性的双数环（\cref{b2:prop:dual-number-flat}）给出一个具体例子：设 \(k\) 为域，\(R=k[\delta]/(\delta^2)\)，取 \(M=(\delta)\) 及包含 \(i:M\hookrightarrow R\)。映射 \(R\rightarrow M\)、\(a\mapsto a\delta\) 的核为 \((\delta)\)，故 \(M\cong R/(\delta)\)。由\cref{b2:prop:tensor-quotient}，
\[
M\otimes_RM\cong R/(\delta),
\]
其中 \(\delta\otimes\delta\) 对应 \(\bar1\)，所以非零。它经 \(i\otimes i\) 映到 \(R\otimes_RR\) 后却变成
\[
\delta\otimes\delta=1\otimes\delta^2=0.
\]
因此 \(T_R(i)\) 已在二次分量中有非零核。泛性质保证线性映射可以延拓成代数同态，不能保证单射也延拓成单射。

\begin{proposition}\label{b2:prop:free-associative-algebra}
设 \(R\) 为交换环，\(M\) 为以 \((x_i)_{i\in I}\) 为基的自由 \(R\) 模，则 \(T_R(M)\) 以所有有限字
\[
1,\quad x_{i_1},\quad x_{i_1}x_{i_2},\quad\ldots
\]
为 \(R\) 模的一组基，乘法为连接字。给定任意 \(R\) 代数 \(A\) 中的一族元素 \((a_i)_{i\in I}\)，存在唯一 \(R\) 代数同态将每个 \(x_i\) 送到 \(a_i\)。因此称 \(T_R(M)\) 为生成元 \((x_i)_{i\in I}\) 上的\term[ziyou-jiehe-daishu]{自由结合代数}[free associative algebra]，记为 \(R\langle x_i\mid i\in I\rangle\)。
\end{proposition}
\begin{proof}
由适用于任意自由基的\cref{b2:prop:multiple-tensor-basis}，每个正次 \(M^{\otimes d}\) 以长度为 \(d\) 的字为基，零次部分 \(R\) 则以空字 \(1\) 为基。不同次数直和后，全部有限字构成基。任意指定 \(x_i\mapsto a_i\) 由自由模的泛性质唯一延拓为 \(M\rightarrow A\)，再由\cref{b2:thm:tensor-algebra-universal}唯一延拓为代数同态。
\end{proof}
\symidx{\(R\langle x_i\mid i\in I\rangle\)：非交换生成元上的自由结合代数}

当 \(R\ne0\) 且有两个不同基向量 \(x,y\) 时，\(xy\)、\(yx\) 是不同的基字，所以 \(xy-yx\ne0\)。只有一个生成元时，每个字由长度决定，\(R\langle x\rangle\cong R[x]\)。一般的 \(T_R(M)\) 自由地接受来自模 \(M\) 的线性映射；只有 \(M\) 是以 \((x_i)\) 为基的自由模时，才可以把它进一步看成符号集合 \(\{x_i\}\) 上不带线性关系的字代数。非自由 \(M\) 中已有的线性关系，在张量代数的一次分量中仍然成立。

\subsection{商代数与按关系定义代数}
\label{b2:subsec:0802:03}

设 \(A\) 是 \(R\) 代数，\(J\) 是双边理想。因为结构标量在中心，\(J\) 自动为 \(R\) 子模，商环 \(A/J\) 继承 \(R\) 代数结构。第三章的模展示只需把关系的线性组合置为零；代数展示还必须让关系在左右乘法下继续成立：若 \(r_\alpha=0\)，则对任意 \(a,b\in A\)，也必须有 \(ar_\alpha b=0\)。所以给定自由代数中的关系族 \((r_\alpha)\) 后，必须取其生成的双边理想 \((r_\alpha)\)，它由有限和 \(\sum_\nu a_\nu r_{\alpha_\nu}b_\nu\) 组成。商代数
\[
R\langle x_i\mid i\in I\rangle/(r_\alpha)
\]
就是按生成元与关系给出的\term[daishu-zhanshi]{代数展示}[presentation of an algebra]，也称代数呈示。\termidx[daishu-chengshi]{代数呈示}这里的关系是代数乘法关系，也可含常数项，不能只把它们看成一次模关系。

\begin{proposition}\label{b2:prop:algebra-relations-universal}
设 \(R\) 为交换环，\(I\) 为指标集，\((r_\alpha)\) 是自由结合 \(R\) 代数 \(R\langle x_i\mid i\in I\rangle\) 中的一族元素，并以 \((r_\alpha)\) 记其生成的双边理想。对任意 \(R\) 代数 \(B\)，从 \(R\langle x_i\mid i\in I\rangle/(r_\alpha)\) 到 \(B\) 的 \(R\) 代数同态，与满足全部关系 \(r_\alpha(b_i)=0\) 的元素族 \((b_i)_{i\in I}\subseteq B\) 一一对应。
\end{proposition}
\begin{proof}
记 \(J=(r_\alpha)\) 为陈述中的双边理想。由\cref{b2:prop:free-associative-algebra}，给定像族恰好确定自由代数上的同态 \(h\)。若全部 \(h(r_\alpha)=0\)，则
\[
h\left(\sum_\nu a_\nu r_{\alpha_\nu}b_\nu\right)
=\sum_\nu h(a_\nu)h(r_{\alpha_\nu})h(b_\nu)=0,
\]
故它将关系理想送到零。公式 \(\bar h(a+J)=h(a)\) 因而不依赖代表元，给出唯一商代数同态。反向把商同态与商映射复合，其像族必满足全部关系。两个过程互逆。
\end{proof}

因此，从展示代数定义同态时，给定生成元的像就确定了唯一性，检查这些像满足全部关系则保证存在性。若还希望商代数保留分次，就必须保证关系理想包含其每个元素的全部齐次分量；否则一条混合不同次数的关系可能把原本应当分开的分量联系起来。由齐次关系生成理想可以保证这一点。

\begin{proposition}\label{b2:prop:homogeneous-ideal-quotient}
设 \(R\) 为交换环，\(A=\bigoplus_{d\ge0}A_d\) 为非负整数分次 \(R\) 代数。由 \(A\) 中齐次元素族生成的双边理想 \(J\) 满足
\[
J=\bigoplus_{d\ge0}(J\cap A_d).
\]
因此商代数有分次 \(A/J\cong\bigoplus_{d\ge0}A_d/(J\cap A_d)\)，商映射保持次数。
\end{proposition}
\begin{proof}
设生成关系 \(r_\alpha\) 为齐次。\(J\) 中元素是 \(a r_\alpha b\) 的有限和；把 \(a,b\) 各自分解成有限个齐次分量，每个乘积都为某个次数的齐次元素，且仍在 \(J\)。所以一个 \(J\) 元素的每个齐次分量都属于 \(J\)，得到所述直和等式。

把各 \(A_d\) 的商映到 \(A/J\) 后，每个商元素都有有限齐次表达。若这样的和为零，则其原像之和在 \(J\)，刚证的分量性质使每个原像分量都在 \(J\cap A_d\)，故各商分量为零。乘法仍按次数相加，由此得到所述分次。
\end{proof}

这样的理想称为\term[qici-lixiang]{齐次理想}[homogeneous ideal]。例如关系 \(xy-yx\) 齐次，次数为二。关系 \(yx-xy-1\) 却混合次数二与零，商后不能继续把 \(x,y\) 都放在次数一，除非单位元也变成了零。为排除这种零环坍缩，下面只需构造一个满足关系且单位元非零的代数。在任意域 \(k\) 上，令 \(X:k[t]\rightarrow k[t]\) 为乘 \(t\)，\(D\) 为形式求导，则对多项式 \(f\)，
\[
(DX-XD)(f)=(tf)'-tf'=f.
\]
所以由\cref{b2:prop:algebra-relations-universal}，赋值 \(x\mapsto X\)、\(y\mapsto D\) 给出 \(k\langle x,y\rangle/(yx-xy-1)\) 到 \(\operatorname{End}_k(k[t])\) 的代数同态，将单位元映成非零恒等算子。于是商中的单位元确实非零；若这个商继承了原次数分次，等式 \(yx-xy=1\) 将使一个非零元素同时属于次数二和零，与分次直和矛盾。微分算子在这里的作用就是证明商非零，从而确认不同次数相混合的障碍实际存在。

\ProseParagraphBreak
在张量代数中，增广理想记录所有正次部分，而它的平方只保留至少含两个一次因子的部分。因此取一次商可以在 \(R\) 模同构的意义下，从带有增广同态的代数恢复原来的模。

\begin{proposition}\label{b2:prop:tensor-augmentation-indecomposables}
设 \(R\) 为交换环，\(M\) 为 \(R\) 模，\(T_R(M)=\bigoplus_{d\ge0}M^{\otimes d}\) 为张量代数，\(\iota:M\hookrightarrow T_R(M)\) 为一次分量包含。由零映射 \(M\rightarrow R\) 诱导的增广同态 \(\varepsilon:T_R(M)\rightarrow R\) 是零次分量的投影。其核 \(I\) 满足
\[
I=\bigoplus_{d\ge1}M^{\otimes d},
\qquad I^2=\bigoplus_{d\ge2}M^{\otimes d},
\]
且映射
\[
M\longrightarrow I/I^2,
\qquad m\longmapsto\iota(m)+I^2
\]
是 \(R\) 模同构，且对 \(M\) 的 \(R\) 线性映射自然。
\end{proposition}
\begin{proof}
取零次分量的投影 \(p:T_R(M)\rightarrow R\) 保持单位元及标量。两个元素相乘时，由非负次数相加，乘积的零次分量恰为两者零次分量之积，所以 \(p\) 是 \(R\) 代数同态。它在一次分量上为零，故由\cref{b2:thm:tensor-algebra-universal}的唯一性等于 \(\varepsilon\)，并给出所述核 \(I\)。

任意两个 \(I\) 元素相乘，所有分量的次数至少为二，所以 \(I^2\subseteq\bigoplus_{d\ge2}M^{\otimes d}\)。反过来，对 \(d\ge2\)，每个纯张量都能写成
\[
m_1\otimes\cdots\otimes m_d
=\iota(m_1)(m_2\otimes\cdots\otimes m_d),
\]
右端两个因子均属于 \(I\)。由纯张量生成性，整个 \(M^{\otimes d}\) 都在 \(I^2\) 中，故得反向包含。于是 \(I=\iota(M)\oplus I^2\) 作为 \(R\) 模成立，取商即得所述同构。这个证明不要求 \(M\) 自由。

给定 \(R\) 线性映射 \(f:M\rightarrow N\)，记 \(I'\) 为 \(T_R(N)\) 的增广理想，\(\iota_N\) 为其一次分量包含。分次同态 \(T_R(f)\) 将 \(I\) 映入 \(I'\)，并将 \(I^2\) 映入 \((I')^2\)；它在一次商上把 \(\iota(m)+I^2\) 送到 \(\iota_N(f(m))+(I')^2\)。这与先作用 \(f\) 再取上述同构相同，故同构自然。
\end{proof}

% !TeX root = ../../Book2.tex
% 纲要标题：### 8.3 对称代数

\section{对称代数}
\label{b2:sec:0803}

张量代数保留因子的排列顺序。若把 \(uv\) 与 \(vu\) 识别，就得到只记录交换乘积的代数。这个商构造在任意交换底环上都成立；把它认成张量幂中的“不变部分”，却通常需要能除以阶乘。本节把这两个对象分别定义，再说明何时存在指定的同构。

% 纲要要点（供目录核对）：
%
% * 对称代数的泛性质
% * 有限自由模上的多项式解释
% * 对称幂的余不变量解释、对称化及特征 p 的取幂
% * 三次对称化的轨道计算及特征二、三的比较
% * 张量代数与对称代数的扩张标量相容性
%

\subsection{对称代数的泛性质}
\label{b2:subsec:0803:01}

\begin{definition}\label{b2:def:symmetric-algebra}
设 \(R\) 为交换含幺环，\(M\) 为幺性 \(R\) 模。令 \(J\) 为 \(T_R(M)\) 中由全部
\[
u\otimes v-v\otimes u,\qquad u,v\in M,
\]
生成的双边理想。商代数
\[
\operatorname{Sym}_R(M)\coloneqq T_R(M)/J
\]
称为 \(M\) 的\term[duichen-daishu]{对称代数}[symmetric algebra]。由齐次二次关系及\cref{b2:prop:homogeneous-ideal-quotient}，它有自然分次；对每个 \(d\ge0\)，其次数 \(d\) 的分量记为 \(\operatorname{Sym}_R^d(M)\)，称为 \(M\) 的 \(d\) 次\term[duichen-mi]{对称幂}[symmetric power]。张量 \(m_1\otimes\cdots\otimes m_d\) 的像记为 \(m_1\cdots m_d\)。
\end{definition}
\symidx{\(\operatorname{Sym}_R(M)\)：模的对称代数}
\symidx{\(\operatorname{Sym}_R^d(M)\)：模的 \(d\) 次对称幂}

关系在次数二，故次数零与一次部分仍分别为 \(R,M\)。全部一次生成元在商中彼此交换；任意两个字可逐个交换其中的生成元，故整个商代数交换。任何从 \(T_R(M)\) 到交换 \(R\) 代数的同态都将这些交换关系送到零，因而唯一经过这个商。所以 \(\operatorname{Sym}_R(M)\) 是强制一次生成元彼此交换得到的普遍交换商。一般模中的线性关系仍予保留，并没有因交换乘法而消失。

\begin{theorem}[对称代数的泛性质]\label{b2:thm:symmetric-algebra-universal}
设 \(R\) 为交换含幺环，\(M\) 为幺性 \(R\) 模。对任意交换含幺 \(R\) 代数 \(A\) 及 \(R\) 线性映射 \(f:M\rightarrow A\)，存在唯一保幺 \(R\) 代数同态
\[
\widehat f:\operatorname{Sym}_R(M)\longrightarrow A
\]
在一次分量上等于 \(f\)。它将 \(m_1\cdots m_d\) 送到 \(f(m_1)\cdots f(m_d)\)。
\end{theorem}
\begin{proof}
\cref{b2:thm:tensor-algebra-universal}先把 \(f\) 延拓为 \(T_R(M)\rightarrow A\)。每个关系的像为 \(f(u)f(v)-f(v)f(u)=0\)，因为 \(A\) 交换。因此延拓同态将 \(J\) 送到零，诱导唯一商同态。反过来，商中的一次分量和结构标量生成全代数，所需同态在它们上的值已指定，故唯一。
\end{proof}

对线性映射 \(f:M\rightarrow N\)，上述定理给出分次同态 \(\operatorname{Sym}_R(f)\)，其在 \(d\) 次部分的公式为
\[
\operatorname{Sym}_R^d(f)(m_1\cdots m_d)=f(m_1)\cdots f(m_d).
\]
这些映射保持恒等与复合，因为在一次生成元上保持。因此对称代数与每次对称幂都是函子。目标交换的条件不能删去：若 \(M=R^2\)，把两个基向量任意映到非交换代数中的两个元素，只有它们交换时才可经 \(\operatorname{Sym}_R(M)\) 延拓。

\ProseParagraphBreak
在次数 \(d\) 中，乘积 \(m_1\cdots m_d\) 不再记录各因子的排列。因此要把张量与它的置换像识别，也就是商去它们的差；这不同于从张量幂中选出已经被每个置换固定的元素。

\begin{proposition}\label{b2:prop:symmetric-power-universal}
设 \(R\) 为交换含幺环，\(M,N\) 为幺性 \(R\) 模，\(d\ge1\)。令 \(S_d\) 通过置换因子作用于 \(M^{\otimes d}\)，\(q_d:M^{\otimes d}\rightarrow\operatorname{Sym}_R^d(M)\) 为取商映射，则 \(q_d\) 的核由所有 \(t-P_\sigma t\)（\(t\in M^{\otimes d}\)、\(\sigma\in S_d\)）生成。因此 \(q_d\) 给出自然同构
\[
\operatorname{Sym}_R^d(M)\cong
M^{\otimes d}/\langle t-P_\sigma t\mid t\in M^{\otimes d},\ \sigma\in S_d\rangle_R.
\]
从 \(\operatorname{Sym}_R^d(M)\) 到 \(N\) 的 \(R\) 线性映射，与满足
\[
b(m_{\sigma(1)},\ldots,m_{\sigma(d)})=b(m_1,\ldots,m_d)
\quad(\sigma\in S_d,\ m_i\in M)
\]
的多线性映射 \(b:M^d\rightarrow N\) 一一对应。这样的 \(b\) 称为\term[duichen-duoxianxing-yingshe]{对称多线性映射}[symmetric multilinear map]。
\end{proposition}
\begin{proof}
设 \(K_d\) 为所述差生成的子模。先看相邻两个位置的交换：在纯张量上，\(t-P_{(i\ i+1)}t\) 是一个二次交换关系左右各乘其余纯张量的结果，因此在 \(J\) 中。任意置换均可通过相邻交换完成。若这些交换依次给出纯张量 \(t=t_0,t_1,\ldots,t_r=P_\sigma t\)，则每个 \(t_{j-1}-t_j\) 都属于 \(J\)，而 \(t-P_\sigma t=\sum_{j=1}^r(t_{j-1}-t_j)\)，故它也属于 \(J\)。这先对纯张量成立，再由线性对所有 \(t\) 成立，故 \(K_d\subseteq J\cap M^{\otimes d}\)。

反向，将 \(J\) 的元素写成 \(a(u\otimes v-v\otimes u)b\) 的有限和，再取次数 \(d\) 的分量，并把 \(a,b\) 展开成纯张量的有限和。所得每项都是某个纯张量与相邻交换后的张量之差，所以属于 \(K_d\)。因此核恰为 \(K_d\)。对任意 \(R\) 线性映射 \(f:M\to M'\)，有 \(f^{\otimes d}P_\sigma=P_\sigma f^{\otimes d}\)，故该商同构与 \(\operatorname{Sym}_R^d(f)\) 相容，确为自然同构。

多线性映射先由\cref{b2:thm:multiple-tensor-universal}唯一合成 \(M^{\otimes d}\rightarrow N\)。变量置换不改变值，恰好要求它在上述核上为零，因而唯一经 \(q_d\) 分解。反向复合显然逐变量线性且对称。
\end{proof}

\subsection{有限自由模的多项式解释}
\label{b2:subsec:0803:02}

\begin{theorem}[多项式解释]\label{b2:thm:symmetric-polynomial-algebra}
设 \(R\) 为交换含幺环，\(M\) 为以 \(e_1,\ldots,e_r\) 为基的幺性有限自由 \(R\) 模，则存在唯一分次 \(R\) 代数同构
\[
\operatorname{Sym}_R(M)\longrightarrow R[X_1,\ldots,X_r],
\qquad e_i\longmapsto X_i.
\]
\end{theorem}
\begin{proof}
将基向量送到不定元给出线性映射 \(M\rightarrow R[X_1,\ldots,X_r]\)。目标交换，故\cref{b2:thm:symmetric-algebra-universal}给出正向同态。反向把多项式 \(\sum_\alpha a_\alpha X^\alpha\) 送到 \(\sum_\alpha a_\alpha e_1^{\alpha_1}\cdots e_r^{\alpha_r}\)。和为有限和，商中的 \(e_i\) 彼此交换，因此多项式的加法与乘法公式保证它是代数同态。两复合分别在各 \(e_i\) 与 \(X_i\) 上为恒等，并在结构标量上为恒等；生成性使两复合在全代数为恒等。两映射保持生成元的次数，故保持分次。
\end{proof}

\begin{proposition}\label{b2:prop:symmetric-power-basis}
设 \(R\) 为交换含幺环，\(M\) 为以 \(e_1,\ldots,e_r\) 为基的幺性有限自由 \(R\) 模，\(d\ge0\)，则次数 \(d\) 的单项式
\[
e_1^{a_1}\cdots e_r^{a_r},\qquad a_i\ge0,\quad a_1+\cdots+a_r=d,
\]
构成 \(\operatorname{Sym}_R^d(M)\) 的一组基。若 \(R\ne0\)、\(r\ge1\)，则其秩为 \(\binom{r+d-1}{d}\)。当 \(r=0\) 时，零次部分为 \(R\)，所有正次部分为零。
\end{proposition}
\begin{proof}
多项式作为有限形式和，其不同单项式的系数唯一；上一同构把全次数 \(d\) 的单项式送到题述元素，所以它们构成基。计算数目时，把 \(d\) 个相同的记号排成一行，并插入 \(r-1\) 道隔板。第 \(i\) 段记号的个数就是 \(a_i\)，允许空段。这与从 \(r+d-1\) 个位置选出 \(d\) 个记号位置一一对应，故数目为所述二项式系数。零模情形直接由定义得到。
\end{proof}

例如 \(R^2\) 的三次对称幂以 \(e_1^3,e_1^2e_2,e_1e_2^2,e_2^3\) 为基，秩为四；相应三次张量幂则有八个有序纯基张量。对称商把 \(e_1\otimes e_1\otimes e_2\) 的三个不同排列都映到 \(e_1^2e_2\)；只有把这三个张量先相加，其像才是 \(3e_1^2e_2\)。单个张量的像不需要乘上排列数。

\ProseParagraphBreak
对有限维 \(k\) 向量空间 \(V\)，\(\operatorname{Sym}_k(V)\) 的一次生成元是向量，而 \(V\) 上的线性坐标函数属于对偶空间 \(V^*\)。因此，由线性坐标函数生成的形式多项式代数应写为 \(\operatorname{Sym}_k(V^*)\)：线性泛函 \(\lambda\) 在 \(v\) 处取值为 \(\lambda(v)\)，乘积 \(\lambda_1\cdots\lambda_d\) 则取值为 \(\prod_j\lambda_j(v)\)。选定 \(V\) 的基后，其对偶基对应熟悉的坐标不定元。两空间虽同维，\(V\) 与 \(V^*\) 却没有由向量空间本身指定的识别，不能把向量生成元直接当作坐标函数。

\ProseParagraphBreak
形式多项式也不等同于它定义的函数。例如在 \(\mathbb F_p\) 上，\(X^p-X\) 是非零多项式，却在每个域元素处取零。对称代数采用形式多项式的乘法和系数，不因所有取值为零就把它的元素置零；\secref{b2:sec:0804}将说明这一区别如何影响不变环的定义。

\subsection{对称张量、对称商与对称化}
\label{b2:subsec:0803:03}

先考虑一个交换两个基向量的作用：在自由模 \(Ru\oplus Rv\) 中，不变子模为 \(R(u+v)\)，余不变量则为 \((Ru\oplus Rv)/R(u-v)\)。前者要求两个坐标本来相同，后者让 \(u,v\) 的类相同。下面把这两种构造用于张量因子的置换。

\ProseParagraphBreak
设 \(R\) 为交换含幺环，\(M\) 为幺性 \(R\) 模，\(d\ge1\)。令 \(S_d\) 按\cref{b2:prop:multiple-tensor-functoriality}的公式作用在 \(M^{\otimes d}\) 上。这一作用的\term[bubianliang]{不变量}[invariants]是子模
\[
(M^{\otimes d})^{S_d}
=\{\,t\in M^{\otimes d}\mid P_\sigma t=t\;\text{对所有}\;\sigma\in S_d\,\},
\]
其中的元素称为\term[duichen-zhangliang]{对称张量}[symmetric tensor]。同一作用的\term[yububianliang]{余不变量}[coinvariants]是商模
\[
(M^{\otimes d})_{S_d}
\coloneqq M^{\otimes d}/\langle t-P_\sigma t\mid t\in M^{\otimes d},\ \sigma\in S_d\rangle_R.
\]
由\cref{b2:prop:symmetric-power-universal}，商映射给出 \(\operatorname{Sym}_R^d(M)\cong(M^{\otimes d})_{S_d}\)。因此包含与取商组成
\[
(M^{\otimes d})^{S_d}\hookrightarrow M^{\otimes d}
\xrightarrow{q_d}\operatorname{Sym}_R^d(M)
\cong(M^{\otimes d})_{S_d},
\]
其中 \(q_d\) 满射，总有指定映射
\[
q_d\big|_{(M^{\otimes d})^{S_d}}:
(M^{\otimes d})^{S_d}\longrightarrow\operatorname{Sym}_R^d(M).
\]
线性映射的张量幂保持不变子模，并与 \(q_d\) 相容，所以这个限制映射是自然的。问题在于它何时成为同构，而不是两个模是否偶然具有相同的秩或维数。
\symidx{\((M^{\otimes d})^{S_d},(M^{\otimes d})_{S_d}\)：对称张量的不变量与余不变量}

\ProseParagraphBreak
在次数二，若 \(2\) 可逆，令 \(P\) 为交换两个因子的算子，则 \((t+Pt)/2\) 被 \(P\) 固定；若 \(t\) 本来对称，这个平均仍为 \(t\)。一般次数下，对全部 \(d!\) 个置换同样取平均，就能得到一个固定不变张量的投影。

\begin{proposition}\label{b2:prop:symmetrization-characteristic}
设 \(R\) 为交换含幺环，\(M\) 为幺性 \(R\) 模，\(d\ge1\)。若 \(d!\) 在 \(R\) 中可逆，则\term[duichen-hua]{对称化}[symmetrization]算子
\[
E_d=\frac1{d!}\sum_{\sigma\in S_d}P_\sigma
\]
是投影到 \((M^{\otimes d})^{S_d}\) 的线性映射，并诱导自然同构
\[
\overline E_d:\operatorname{Sym}_R^d(M)\xrightarrow{\ \cong\ }(M^{\otimes d})^{S_d}.
\]
其逆是商映射 \(q_d\) 在不变子模上的限制。
\end{proposition}
\begin{proof}
左乘一个固定置换只重新排列求和项，因此 \(P_\pi E_d=E_d\)，像都不变；若 \(t\) 已不变，则 \(E_dt=t\)。所以像正是不变子模，且 \(E_d^2=E_d\)。

同样，\(E_dP_\pi=E_d\)，故 \(E_d\) 零化全部 \(t-P_\pi t\)。由\cref{b2:prop:symmetric-power-universal}，它唯一经 \(q_d\) 分解为 \(\overline E_d\)。在对称商中 \(q_dP_\sigma=q_d\)，因而 \(q_dE_d=q_d\)，说明 \(q_d\overline E_d\) 为恒等；在不变子模上 \(E_d\) 为恒等，又说明 \(\overline E_dq_d\) 为恒等。线性映射的张量幂与每个 \(P_\sigma\) 交换，所以也与平均算子相容；所得同构因此自然。
\end{proof}

阶乘可逆保证上述同构，但它不可逆时须另行检查，不可据此一概断言同构失败。特征零也不保证阶乘可逆：取 \(R=\mathbb Z\)、\(M=\mathbb Ze\oplus\mathbb Zf\)。交换两个因子的不变张量恰为
\[
a(e\otimes e)+b(e\otimes f+f\otimes e)+c(f\otimes f),
\qquad a,b,c\in\mathbb Z,
\]
因为在纯基张量的展开中，\(e\otimes f\) 与 \(f\otimes e\) 的系数必须相同。它在 \(\operatorname{Sym}_{\mathbb Z}^2(M)\) 中的像为
\[
ae^2+2bef+cf^2.
\]
\(e^2,ef,f^2\) 是这个对称幂的一组基，所以 \(ef\) 不在像中，否则需要整数 \(b\) 满足 \(2b=1\)。这里商映射在不变子模上的限制不是满射，障碍是 \(2\) 在 \(\mathbb Z\) 中不可逆。

\ProseParagraphBreak
正特征下还可能出现单射性失败。取 \(k\) 特征为素数 \(p\)，\(V=ke\oplus kf\)，在 \(V^{\otimes p}\) 中令
\[
s=\sum_{j=1}^p e\otimes\cdots\otimes
\underset{\text{第 }j\text{ 个位置}}{f}
\otimes\cdots\otimes e.
\]
各项是互不相同的纯基张量，所以 \(s\ne0\)；置换只重排这些项，所以 \(s\) 不变。但
\[
q_p(s)=p\,e^{p-1}f=0.
\]
因此商映射在不变张量上不是单射。最小例子是特征二下的 \(e\otimes f+f\otimes e\)：它非零且对称，却在二次对称幂中为零。反过来，\(ef\) 仍是对称幂的一个非零基元素。

\begin{example}\label{b2:exm:cubic-symmetrization-orbits}
设 \(k\) 是域，\(V=ke\oplus kf\)。在 \(V^{\otimes3}\) 的八个纯基张量上，\(S_3\) 的置换作用有四个轨道，按 \(f\) 出现零、一、二、三次分类。相应轨道和为
\[
\begin{aligned}
s_0&=e\otimes e\otimes e,\\
s_1&=e\otimes e\otimes f+e\otimes f\otimes e+f\otimes e\otimes e,\\
s_2&=e\otimes f\otimes f+f\otimes e\otimes f+f\otimes f\otimes e,\\
s_3&=f\otimes f\otimes f.
\end{aligned}
\]
置换作用只重排纯基张量，因而一个张量不变，当且仅当同一轨道上的系数相同。四个轨道的支集互不相交，所以 \(s_0,s_1,s_2,s_3\) 线性无关，并生成全部不变张量。它们是一组基；\(\operatorname{Sym}_k^3(V)\) 则以 \(e^3,e^2f,ef^2,f^3\) 为基。限制商映射把 \(s_0,s_1,s_2,s_3\) 分别送到 \(e^3,3e^2f,3ef^2,f^3\)，因此相对于这两组基的矩阵为
\[
\operatorname{diag}(1,3,3,1).
\]
若 \(k=\mathbb Q\)，它是同构，且 \(e^2f\) 的逆像为 \(s_1/3\)。这也等于六个置换的平均：\(e\otimes e\otimes f\) 的每个不同排列出现两次，所以平均为 \(2s_1/6=s_1/3\)。

若 \(\operatorname{char}k=2\)，上述矩阵为单位矩阵，限制映射仍是同构，尽管 \(3!=0\) 不可逆。若 \(\operatorname{char}k=3\)，它的核为 \(ks_1\oplus ks_2\)，像为 \(ke^3\oplus kf^3\)，秩为二；两端都是四维，也不保证指定映射可逆。因此阶乘可逆是统一保证同构的充分条件，并不是一个特定模上限制映射可逆的必要条件。
\end{example}

删去平均公式中的分母，就得到始终有定义的求和算子 \(N_d=\sum_\sigma P_\sigma\)。但它可能把本来的不变张量也送到零。例如在特征 \(p\) 的域 \(k\) 上，一维空间 \(ke\) 的次数 \(p\) 张量幂中，全部置换作用都是恒等，故 \(N_p=p!\operatorname{id}=0\)，而张量幂、不变子空间和对称幂都非零。因此只删去分母不能补救平均公式；这里限制商映射反而是同构，失败的是这个求和算子。

\ProseParagraphBreak
正特征不仅影响平均过程，还使对称代数内部的取幂运算出现新的加法性。设 \(p\) 为素数，\(R\) 为特征 \(p\) 的交换含幺环，\(M\) 为幺性 \(R\) 模。对称代数交换，故
\[
(v+w)^p=v^p+w^p,\qquad (av)^p=a^pv^p.
\]
第一式由二项式系数 \(\binom pi\) 对 \(0<i<p\) 都被 \(p\) 整除得到。因此 \(\Phi_p:M\rightarrow\operatorname{Sym}_R^p(M)\)、\(v\mapsto v^p\) 是可加映射，并对 Frobenius 环同态 \(F_R:R\rightarrow R\)、\(a\mapsto a^p\) 半线性：\(\Phi_p(av)=F_R(a)\Phi_p(v)\)。这里半线性指映射可加，而标量 \(a\) 须先经过 \(F_R\) 才作用于像；若 \(F_R\) 为恒等，就得到通常的 \(R\) 线性。
\symidx{\(F_R\)：特征 \(p\) 环上的 Frobenius 同态}

\ProseParagraphBreak
若进一步 \(R=k\) 为域且 \(M\ne0\)，对非零 \(v\in M\)，先在 \(kv\) 上定义 \(av\mapsto a\)，再由\cref{b2:lem:functional-extension}延拓为 \(\ell:M\to k\)，便有 \(\ell(v)=1\)。对称多线性映射 \((m_1,\ldots,m_p)\mapsto\prod_i\ell(m_i)\) 由\cref{b2:prop:symmetric-power-universal}诱导线性映射 \(\operatorname{Sym}_k^p(M)\to k\)，将 \(v^p\) 送到 \(1\)，故 \(v^p\ne0\)。若 \(\Phi_p\) 是 \(k\) 线性的，便有 \((a^p-a)v^p=0\) 对每个 \(a\in k\) 成立；域上的非零向量不能被非零标量零化，故 \(a^p=a\)。因此在这个域且模非零的范围内，\(\Phi_p\) 线性当且仅当 \(F_k\) 为恒等。一般环上不能由 \(\Phi_p\) 线性就这样消去 \(v^p\)。因子对称、取对称商和取 \(p\) 次幂是相互联系但不同的运算。

\ProseParagraphBreak
第六章把扩张标量写成 \(S\otimes_R-\)。现在的代数构造只用有限张量、直和与生成关系，改变底环后也能与扩张标量相容；不必为模选择一组基。

\begin{proposition}\label{b2:prop:tensor-symmetric-base-change}
设 \(R,S\) 是交换含幺环，\(\varphi:R\to S\) 是保幺环同态，\(M\) 是幺性 \(R\) 模，记 \(P=S\otimes_RM\)。存在自然的分次 \(S\) 代数同构
\[
\begin{aligned}
\alpha_T&:S\otimes_RT_R(M)\xrightarrow{\ \cong\ }T_S(P),\\
\alpha_{\mathrm{Sym}}&:S\otimes_R\operatorname{Sym}_R(M)
\xrightarrow{\ \cong\ }\operatorname{Sym}_S(P).
\end{aligned}
\]
左侧的乘法均由 \((s\otimes a)(t\otimes b)=st\otimes ab\) 给出，次数 \(d\) 的分量分别为 \(S\otimes_RM^{\otimes_Rd}\) 与 \(S\otimes_R\operatorname{Sym}_R^d(M)\)。零次的同构均为 \(S\otimes_RR\to S\)、\(s\otimes r\mapsto s\varphi(r)\)。对 \(d\ge1\)，其公式为
\[
\begin{aligned}
\alpha_T\bigl(s\otimes(m_1\otimes_R\cdots\otimes_Rm_d)\bigr)
&=(s\otimes_Rm_1)\otimes_S(1\otimes_Rm_2)
\otimes_S\cdots\otimes_S(1\otimes_Rm_d),\\
\alpha_{\mathrm{Sym}}\bigl(s\otimes(m_1\cdots m_d)\bigr)
&=(s\otimes_Rm_1)(1\otimes_Rm_2)\cdots(1\otimes_Rm_d).
\end{aligned}
\]
当 \(d=1\) 时，右侧只保留首因子。不要求 \(S\) 作为 \(R\) 模平坦。
\end{proposition}
\begin{proof}
先给左侧代数结构。对任意含幺 \(R\) 代数 \(A\)，因 \(R\) 的结构像在 \(A\) 中心，所述乘法保持两处 \(R\) 平衡关系。例如
\[
(s\varphi(r)\otimes a)(t\otimes b)
=st\varphi(r)\otimes ab
=st\otimes (ra)b
=(s\otimes ra)(t\otimes b),
\]
另一处由 \(arb=abr\) 同样核对。因此乘法良定义；结合律与单位元 \(1\otimes1\) 由 \(S,A\) 的对应等式继承，\(s\mapsto s\otimes1\) 的像在中心，故得到 \(S\) 代数。若 \(A\) 分次，\cref{b2:prop:tensor-direct-sums}给出 \(S\otimes_RA\cong\bigoplus_d(S\otimes_RA_d)\)，且上述乘法使次数相加。

令 \(\iota_R:M\to T_R(M)\)、\(\iota_S:P\to T_S(P)\) 为一次分量包含。把 \(T_S(P)\) 经 \(\varphi\) 看成 \(R\) 代数，映射 \(m\mapsto\iota_S(1\otimes m)\) 为 \(R\) 线性：平衡关系 \(1\otimes rm=\varphi(r)\otimes m\) 保证标量相容。\cref{b2:thm:tensor-algebra-universal}给出 \(R\) 代数同态 \(h:T_R(M)\to T_S(P)\)。公式
\[
\alpha_T(s\otimes a)=s\cdot h(a)
\]
对 \(R\) 平衡，因为 \(h(ra)=\varphi(r)h(a)\)，所以诱导 \(S\) 线性映射。\(S\) 标量在目标中心，故
\[
\alpha_T\bigl((s\otimes a)(t\otimes b)\bigr)
=st\cdot h(ab)
=(s\cdot h(a))(t\cdot h(b)),
\]
并且它保单位元。\(h\) 在一次生成元上的值给出题述纯张量公式；将标量 \(s\) 移入第一个 \(P\) 因子即得到所写形式。

反向，\(S\) 线性映射
\[
j:P\to S\otimes_RT_R(M),\qquad s\otimes m\longmapsto s\otimes\iota_R(m)
\]
由 \(\operatorname{id}_S\otimes\iota_R\) 给出，故良定义。再次用张量代数的泛性质，得到 \(S\) 代数同态 \(\beta_T:T_S(P)\to S\otimes_RT_R(M)\)。它在纯张量上的公式为
\[
\beta_T\bigl((s_1\otimes m_1)\otimes_S\cdots\otimes_S(s_d\otimes m_d)\bigr)
=(s_1\cdots s_d)\otimes(m_1\otimes_R\cdots\otimes_Rm_d),
\]
在零次部分则为 \(s\mapsto s\otimes1\)。两个复合都固定 \(S\) 标量；\(\alpha_T\beta_T\) 在每个一次生成元 \(s\otimes m\in P\) 上为恒等，\(\beta_T\alpha_T\) 在 \(s\otimes\iota_R(m)\) 上为恒等。这些元素连同 \(S\) 生成相应代数，故两复合在全代数上为恒等。两向映射把一次生成元送到一次分量，把标量送到零次分量，所以保持分次。

对称情形中，把 \(m\) 送到 \(\operatorname{Sym}_S(P)\) 的一次生成元 \(1\otimes m\)，由\cref{b2:thm:symmetric-algebra-universal}延拓为 \(R\) 代数同态 \(h_{\mathrm{Sym}}:\operatorname{Sym}_R(M)\to\operatorname{Sym}_S(P)\)。再定义 \(\alpha_{\mathrm{Sym}}(s\otimes a)=s\cdot h_{\mathrm{Sym}}(a)\)；与张量情形相同的平衡等式和乘法等式说明它是 \(S\) 代数同态，且给出题述公式。

反向，\(S\otimes_R\operatorname{Sym}_R(M)\) 是交换 \(S\) 代数。将 \(s\otimes m\in P\) 送到目标代数一次分量 \(S\otimes_R\operatorname{Sym}_R^1(M)\) 中的对应元素，得到 \(S\) 线性映射；再用对称代数的泛性质，得到 \(\beta_{\mathrm{Sym}}:\operatorname{Sym}_S(P)\to S\otimes_R\operatorname{Sym}_R(M)\)。具体地，
\[
\beta_{\mathrm{Sym}}\bigl((s_1\otimes m_1)\cdots(s_d\otimes m_d)\bigr)
=(s_1\cdots s_d)\otimes(m_1\cdots m_d).
\]
两个复合在 \(S\) 标量和一次生成元上都为恒等，故在全代数上为恒等；生成元的次数也得到保留，因而同构分次。这里所用的是两个交换代数的泛性质，并未要求张量积保持单射。

对任意 \(R\) 线性映射 \(f:M\to M'\)，先扩张标量再作用于一次生成元，与先作用于 \(m\) 再按上述公式映射，都得到 \(s\otimes f(m)\)；在 \(S\) 标量上也相同。生成性因此给出
\[
\begin{aligned}
\alpha_{T,M'}(\operatorname{id}_S\otimes T_R(f))
&=T_S(\operatorname{id}_S\otimes f)\alpha_{T,M},\\
\alpha_{\mathrm{Sym},M'}(\operatorname{id}_S\otimes\operatorname{Sym}_R(f))
&=\operatorname{Sym}_S(\operatorname{id}_S\otimes f)\alpha_{\mathrm{Sym},M}.
\end{aligned}
\]
这证明两个同构对 \(M\) 自然。若 \(S'\) 也是交换含幺 \(R\) 代数，并给定保幺 \(R\) 代数同态 \(\psi:S\to S'\)，两侧由 \(s\mapsto\psi(s)\)、\(s\otimes m\mapsto\psi(s)\otimes m\) 诱导的映射，也与这些同构相容：上述纯张量公式两种复合均只将系数 \(s\) 换成 \(\psi(s)\)，生成性使相容性在全代数上成立。
\end{proof}

% !TeX root = ../../Book2.tex
% 纲要标题：### 8.4 张量运算与表示
% 本轮补充的纲要细目：
% * 由对偶空间的对称代数定义形式多项式上的群作用，写清逆向作用、不变环及有限域上与点集函数的区别
% * 以反射及单项式权计算不变环，不将逐次不动空间的计算自动视为全部生成关系已知

\section{张量运算与表示}
\label{b2:sec:0804}

若群同时作用于几个向量空间，多线性构造也随之得到群作用。这个过程使一个表示产生新的表示，并把同态问题转为不变张量问题。本节以任意域 \(k\) 和任意群 \(G\) 开始；凡需有限维或特征零处都另行说明。因子置换由 \(S_d\) 作用，表示的不变量由 \(G\) 作用；这两个群作用应始终区分。

% 纲要要点（供目录核对）：
%
% * 对偶、张量积与对称幂表示
% * 不变量和对称张量的区别
% * 第五卷 张量表示与 Schur–Weyl 理论
%
% ---
%

\subsection{对偶、张量积与对称幂表示}
\label{b2:subsec:0804:01}

\begin{definition}\label{b2:def:linear-representation}
设 \(k\) 为域，\(G\) 为群，\(V\) 为 \(k\) 向量空间。如果给定群同态
\[
\rho_V:G\longrightarrow\operatorname{GL}_k(V),
\]
就称 \(\rho_V\) 为 \(G\) 在 \(V\) 上的\term[xianxing-biaoshi]{线性表示}[linear representation]。
常把 \(\rho_V(g)v\) 写成 \(gv\)。在本节给定的表示中，显式作用记号 \(g\cdot v\) 也表示同一个 \(\rho_V(g)v\)；对另一个表示 \(W\)，则按它自己的 \(\rho_W\) 解释。若 \(V,W\) 均为 \(G\) 的 \(k\) 线性表示，\(k\) 线性映射 \(f:V\rightarrow W\) 满足 \(f(gv)=g\cdot f(v)\) 对每个 \(g\in G\)、\(v\in V\) 成立，则称 \(f\) 为表示同态或等变线性映射；全部这类映射组成空间 \(\operatorname{Hom}_G(V,W)\)。
\end{definition}
\symidx{\(\operatorname{Hom}_G(V,W)\)：表示之间的等变线性映射空间}
\symidx{\(\rho_V(g)\)：群元素在表示空间上的线性作用}

群作用公理在这里具体为 \(1v=v\)、\((gh)v=g(hv)\)，并要求每个 \(g\) 的作用可逆且线性。定义没有要求表示忠实，也没有要求群有限。例如无穷循环群由一个生成元 \(g\) 生成，在 \(\mathbb Q^2\) 上可以规定 \(g\) 的作用矩阵为 \(\operatorname{diag}(2,1/2)\)，再由整数次幂定义全部群元素的作用。

\ProseParagraphBreak
向量与线性泛函同时变换时，自然的求值配对应保持不变，也就是要求 \((g\lambda)(gv)=\lambda(v)\)，其中 \(k\) 上采用平凡作用。把 \(gv\) 记为 \(w\)，这个要求便强迫 \((g\lambda)(w)=\lambda(g^{-1}w)\)。因此对偶作用中的逆元来自求值配对的相容性：向量的变换与泛函的变换必须互相抵消。

\begin{proposition}\label{b2:prop:tensor-dual-symmetric-representations}
设 \(k\) 为域，\(G\) 为群，\(V,W\) 为 \(G\) 的 \(k\) 线性表示，\(d\ge0\)。以下公式分别定义 \(V^*\)、\(V\otimes_kW\) 及 \(\operatorname{Sym}_k^d(V)\) 上的表示：
\[
\begin{aligned}
(g\lambda)(v)&=\lambda(g^{-1}v),\\
g(v\otimes w)&=gv\otimes gw,\\
g(v_1\cdots v_d)&=(gv_1)\cdots(gv_d).
\end{aligned}
\]
它们分别称为\term[duiou-biaoshi]{对偶表示}[dual representation]、\term[zhangliangji-biaoshi]{张量积表示}[tensor product representation]和\term[duichenmi-biaoshi]{对称幂表示}[symmetric power representation]。
\end{proposition}
\begin{proof}
对偶公式对 \(v\) 线性，也对 \(\lambda\) 线性。对 \(g,h\in G\)，
\[
\bigl(g(h\lambda)\bigr)(v)
=(h\lambda)(g^{-1}v)
=\lambda(h^{-1}g^{-1}v)
=\bigl((gh)\lambda\bigr)(v).
\]
单位元作用为恒等，\(g^{-1}\) 的作用为逆，故为表示。使用逆元正是为了抵消对偶映射反转复合次序的效应。

张量公式由 \(\rho_V(g)\otimes\rho_W(g)\) 给出，良定义来自\cref{b2:prop:tensor-maps}；该命题的复合公式保证群乘法，逆由 \(g^{-1}\) 给出。对称幂的映射是 \(\operatorname{Sym}_k^d\bigl(\rho_V(g)\bigr)\)，对称幂函子保持恒等和复合，因此同样为表示。也可直接观察：张量幂作用将每个交换关系送到另一个交换关系，所以它通过对称商。
\end{proof}

对偶构造把特征值取逆，张量构造则把各因子的特征值相乘。具体地，若 \(g\) 在有限维 \(V\) 的基 \(e_1,\ldots,e_r\) 上对角作用，\(ge_i=a_i e_i\)，则它在对偶基上的特征值为 \(a_i^{-1}\)，在张量基 \(e_{i_1}\otimes\cdots\otimes e_{i_d}\) 上的特征值为 \(a_{i_1}\cdots a_{i_d}\)，在对称基 \(e_1^{b_1}\cdots e_r^{b_r}\) 上的特征值为 \(a_1^{b_1}\cdots a_r^{b_r}\)。这里各 \(a_i\ne0\)，因为群元素作用可逆。若一个张量被 \(g\) 固定，其基展开中只能出现上述特征值等于一的项。

\ProseParagraphBreak
\cref{b2:prop:finite-free-hom-tensor}已在有限自由模上建立求值同构；有限维向量空间当然满足该条件。给 \(V,W\) 加上群作用后，还要使这个同构与作用相容。对线性映射 \(A:V\to W\)，变换输入时先由 \(g^{-1}\) 拉回，施加 \(A\) 后再由 \(g\) 推进输出，这给出下面的 Hom 作用。

\begin{proposition}\label{b2:prop:hom-as-tensor-representation}
设 \(k\) 为域，\(V\) 为有限维 \(k\) 向量空间，\(W\) 为任意 \(k\) 向量空间，则
\[
\Theta:V^*\otimes_kW\longrightarrow\operatorname{Hom}_k(V,W),
\qquad
\lambda\otimes w\longmapsto\bigl(v\mapsto\lambda(v)w\bigr)
\]
是自然线性同构。若群 \(G\) 在 \(V,W\) 上均有 \(k\) 线性表示，则它是 \(G\) 表示同构，其中右端的作用为
\[
(g\cdot A)(v)=\rho_W(g)\bigl(A(\rho_V(g)^{-1}v)\bigr).
\]
\end{proposition}
\begin{proof}
取\cref{b2:prop:finite-free-hom-tensor}中的 \(R=k\)、\(P=V\)、\(M=W\)，即得所写自然线性同构。该命题还给出坐标逆式
\[
\Theta^{-1}(A)=\sum_i e_i^*\otimes A(e_i),
\]
其中 \(e_i\) 是 \(V\) 的任意基，\(e_i^*\) 是其对偶基；该表达式所确定的张量不依赖所选基。

群作用的相容性直接在纯张量上核对：
\[
\Theta(g\lambda\otimes gw)(v)
=\lambda(\rho_V(g)^{-1}v)\,\rho_W(g)w
=\rho_W(g)\bigl(\Theta(\lambda\otimes w)(\rho_V(g)^{-1}v)\bigr).
\]
这就是右端所给的作用。它满足群乘法公理，因为对 \(A:V\to W\) 先作用 \(h\) 再作用 \(g\)，所得为 \(\rho_W(gh)A\rho_V(gh)^{-1}\)。
\end{proof}

当 \(W=V\) 时，恒等算子对应的张量 \(\sum_i e_i^*\otimes e_i\) 不依赖基，并被所有 \(g\) 固定，因为 \(\rho_V(g)\operatorname{id}_V\rho_V(g)^{-1}=\operatorname{id}_V\)。向量与对偶向量的变换互相抵消，并不要求每个基向量被固定。一般算子的矩阵也可按列作这种转换：若 \(V\) 有基 \(e_1,e_2\)，且 \(A\) 的矩阵为 \(\begin{pmatrix}a&b\\c&d\end{pmatrix}\)，则
\[
\Theta^{-1}(A)
=e_1^*\otimes(ae_1+ce_2)+e_2^*\otimes(be_1+de_2).
\]
恒等矩阵便给出 \(e_1^*\otimes e_1+e_2^*\otimes e_2\)。

\subsection{不变量与对称张量}
\label{b2:subsec:0804:02}

\begin{definition}\label{b2:def:invariant-vector}
设 \(k\) 为域，\(G\) 为群，\(V\) 为 \(G\) 的 \(k\) 线性表示。若 \(v\in V\) 满足 \(gv=v\) 对每个 \(g\in G\) 成立，就称 \(v\) 为 \(G\) 的\term[bubian-xiangliang]{不变向量}[invariant vector]。全部不变向量组成的子空间记为
\[
V^G=\{\,v\in V\mid gv=v\;\text{对每个}\;g\in G\,\}.
\]
若子空间 \(U\subseteq V\) 满足 \(gU\subseteq U\) 对每个 \(g\in G\) 成立，就称 \(U\) 为 \(G\) 不变子空间，也称稳定子空间；这只要求作用保留整个子空间，不能因此断言其中每个向量都被逐点固定。
\end{definition}
\symidx{\(V^G\)：群作用下逐点不变的向量空间}

不变条件由一族线性方程给出，故 \(V^G\) 确为子空间。对任意两个 \(G\) 表示 \(V,W\)，上一命题所写的公式都在 \(\operatorname{Hom}_k(V,W)\) 上定义作用，不要求 \(V\) 有限维。由此有
\[
\operatorname{Hom}_k(V,W)^G=\operatorname{Hom}_G(V,W).
\]
确切地说，\(A:V\to W\) 在上述作用下不变，当且仅当对每个 \(g\in G\) 都有
\[
\rho_W(g)A\rho_V(g)^{-1}=A,
\qquad\text{即}\qquad
\rho_W(g)A=A\rho_V(g).
\]
后一个等式正是等变性。因此当 \(V\) 有限维时，\cref{b2:prop:hom-as-tensor-representation}给出
\[
(V^*\otimes_kW)^G\cong\operatorname{Hom}_G(V,W).
\]
寻找等变映射，可以转为求张量表示中的固定向量。

\ProseParagraphBreak
对称张量采用的是 \(S_d\) 对因子位置的作用，它只重新排列向量；\(G\) 不变张量采用的则是 \(G\) 的对角作用，它保持位置而同时改变各位置上的向量。即使两种作用发生在同一个空间，也不是同一个条件。取 \(G=\langle g\rangle\) 为无穷循环群，\(V=\mathbb Qe\oplus\mathbb Qf\)，规定 \(ge=2e\)、\(gf=\tfrac12f\)。在 \(V^{\otimes2}\) 中，\(e\otimes e\) 对交换两因子保持不变，却被 \(g\) 乘四，因而不是 \(G\) 不变量；\(e\otimes f\) 被 \(G\) 固定，却被交换因子送到不同的基张量 \(f\otimes e\)，因而不是对称张量。

\ProseParagraphBreak
本例中
\[
(V^{\otimes2})^G=\operatorname{span}_{\mathbb Q}\{e\otimes f,f\otimes e\},
\]
因为四个纯基张量的特征值依次为 \(4,1,1,1/4\)。其中的对称部分是一维空间
\(\mathbb Q(e\otimes f+f\otimes e)\)。二次对称幂的 \(G\) 不变部分则为 \(\mathbb Q ef\)，两者由平均同构对应；该对应确实使用了二的可逆性。

\ProseParagraphBreak
固定向量也可以在对称代数的各个齐次分量中寻找。\secref{b2:sec:0803}已说明，\(V^*\) 的元素是 \(V\) 上的线性坐标函数，而 \(\operatorname{Sym}_k(V^*)\) 用交换乘法组合这些函数。把各次不变量连同乘法一起考察，便得到多项式不变环。

\begin{definition}\label{b2:def:polynomial-invariant-ring}
设 \(k\) 为域，\(V\) 为有限维 \(k\) 向量空间。称对称代数
\[
k[V]\coloneqq\operatorname{Sym}_k(V^*)
\]
为 \(V\) 上的多项式代数。若群 \(G\) 线性作用于 \(V\)，就在生成元 \(\lambda\in V^*\) 上规定 \(g\lambda=\lambda\circ g^{-1}\)，并按代数同态延拓，得到 \(G\) 在 \(k[V]\) 上的作用。被所有群元素固定的多项式称为\term[bubian-duoxiangshi]{不变多项式}[invariant polynomial]；它们组成的子代数
\[
k[V]^G=\{\,f\in k[V]\mid gf=f\;\text{对每个}\;g\in G\,\}
\]
称为这一作用的\term[bubian-huan]{不变环}[invariant ring]。
\end{definition}
\symidx{\(k[V]^G\)：表示空间上的多项式不变量环}
\symidx{\(k[V]=\operatorname{Sym}_k(V^*)\)：向量空间上的多项式代数}

\cref{b2:prop:tensor-dual-symmetric-representations}保证这个延拓满足群乘法，并保持每个齐次分量。在一组基及其对偶坐标 \(x_1,\ldots,x_n\) 下，\cref{b2:thm:symmetric-polynomial-algebra}把 \(k[V]\) 识别为 \(k[x_1,\ldots,x_n]\)。把多项式在 \(v\in V\) 处求值，便有
\[
(gf)(v)=f(g^{-1}v).
\]
这是函数沿 \(g^{-1}:V\to V\) 的拉回作用。使用逆元使拉回后的复合符合群乘法的顺序；在形式多项式代数中，它就是由线性坐标函数的对偶作用诱导的代数自同构。因而不变多项式的和与积仍不变，单位元也不变，定义中的集合确为子代数。

\ProseParagraphBreak
逐次求不动空间，只得到不变环的各个向量空间分量；要说明整个环如何生成，还须研究这些分量之间的乘法。例如，即使已知二次不变量恰为 \(kq\)，其中 \(q\ne0\)，也不能据此断言 \(k[V]^G=k[q]\)：还须知道其他次数是否出现不能由 \(q\) 的幂表示的不变量。找到多个生成元后，还要确定它们之间的关系。第五卷第22章将把这些问题与 Hilbert 级数、生成关系及商联系起来。\footnote{希尔伯特（David Hilbert，1862-1943），德国数学家，证明多项式环的基定理，并推进不变量论与几何公理化。}

\ProseParagraphBreak
在无限域上，不同多项式给出不同的点集函数。一元非零多项式的根数不超过次数；多元情形将多项式按最后一个变量展开，固定其余变量后应用一元结论，再对系数多项式归纳，即知在 \(k^n\) 上处处为零的多项式只能为零。在有限域上则不然。例如 \(k=\mathbb F_q\)、\(V=k\) 时，\(x^q-x\) 的最高次系数为一，因而是非零多项式；但 \(k^\times\) 是阶为 \(q-1\) 的群，Lagrange 定理给出每个非零 \(a\in k\) 都满足 \(a^{q-1}=1\)，加上 \(a=0\) 便得所有 \(a\) 都满足 \(a^q-a=0\)。所以这个非零多项式定义零函数。这里的 \(k[V]\) 始终指形式多项式代数；不变性要求代数中的等式，而不能一概只检查有限个点上的取值。

\newpage
\begin{example}\label{b2:exm:reflection-polynomial-invariants}
设 \(\operatorname{char}k\ne2\)，群 \(C_2=\{1,s\}\) 在 \(V=k^2\) 上由 \(s(a,b)=(-a,b)\) 作用。则
\[
k[V]^{C_2}=k[x^2,y],\qquad
\bigl(\operatorname{Sym}^2_k(V^*)\bigr)^{C_2}
=kx^2\oplus ky^2.
\]
\end{example}
\begin{proof}
由于 \(s^{-1}=s\)，对偶坐标满足 \(sx=-x,sy=y\)。将 \(f=\sum_{i,j}c_{ij}x^iy^j\) 代入 \(sf=f\)，逐个比较单项式系数，得到奇数 \(i\) 对应的 \(2c_{ij}=0\)，故这些系数全零。剩下的单项式都是 \((x^2)^r y^j\)，所以每个不变多项式都是 \(x^2,y\) 的多项式；反过来，\(x^2,y\) 都不变，它们的每个多项式也都不变。这给出了整个不变环，因为所分类的单项式覆盖任意次数。次数二时只剩 \(x^2,y^2\)。这个系数比较适用于有限域，无需用点集函数的相等代替多项式相等。
\end{proof}

\cref{b2:tab:08-group-actions}汇总本节出现的作用；同一张量幂上的两行分别来自 \(G\) 和 \(S_d\)。
\begin{table}[htbp]
\centering
\caption{表示及其张量构造上的群作用}\label{b2:tab:08-group-actions}
\begin{tabularx}{\linewidth}{@{}l l X@{}}
\toprule
空间 & 作用群 & 作用公式 \\
\midrule
\(V\) & \(G\) & \(v\mapsto \rho_V(g)v\) \\
\(V^*\) & \(G\) & \(\lambda\mapsto\lambda\circ\rho_V(g)^{-1}\) \\
\(V^{\otimes d}\) & \(G\) & \(\rho_V(g)^{\otimes d}(\bigotimes_i v_i)=\bigotimes_i\rho_V(g)v_i\) \\
\(V^{\otimes d}\) & \(S_d\) & \(P_\sigma(\bigotimes_i v_i)=\bigotimes_i v_{\sigma^{-1}(i)}\) \\
\(k[V]\) & \(G\) & \(g\lambda=\lambda\circ\rho_V(g)^{-1}\) 对 \(\lambda\in V^*\)，再按代数同态延拓 \\
\bottomrule
\end{tabularx}
\end{table}
\cref{b2:tab:08-group-actions}中的 \(G\) 在张量幂上同时变换各个因子，\(S_d\) 则只调换位置。下一命题先对整个 \(\operatorname{GL}_k(V)\) 证明两种作用相交换，因此也适用于通过 \(\rho_V\) 作用的任意群 \(G\)。

\subsection{张量表示与 Schur–Weyl 理论}
\label{b2:subsec:0804:03}

\begin{proposition}\label{b2:prop:commuting-tensor-actions}
设 \(k\) 为域，\(V\) 为 \(k\) 向量空间，\(d\ge1\)。在 \(V^{\otimes d}\) 上，\(\operatorname{GL}_k(V)\) 的对角作用 \(g\mapsto g^{\otimes d}\) 与 \(S_d\) 的位置置换作用 \(\sigma\mapsto P_\sigma\) 相交换：
\[
g^{\otimes d}P_\sigma=P_\sigma g^{\otimes d}.
\]
因此对称张量子空间和定义对称幂的关系子空间都在一般线性群作用下稳定。
\end{proposition}
\begin{proof}
对 \(v_1\otimes\cdots\otimes v_d\)，两个复合都得到
\[
gv_{\sigma^{-1}(1)}\otimes\cdots\otimes gv_{\sigma^{-1}(d)}.
\]
纯张量生成全空间，所以两算子相等。若 \(t\) 对所有置换不变，则 \(P_\sigma g^{\otimes d}t=g^{\otimes d}t\)；若一个关系为 \(t-P_\sigma t\)，其像为 \(g^{\otimes d}t-P_\sigma g^{\otimes d}t\)，仍是关系。于是两个稳定性结论分别成立。
\end{proof}

这已经说明两种对称性可以同时研究，但“作用交换”并不等于“每一种作用已经给出了另一种作用的全部对称算子”。为写清后一问题，设 \(k\) 为域，\(E\) 为 \(k\) 向量空间。若 \(\mathcal A\) 为 \(\operatorname{End}_k(E)\) 的子代数，记
\[
\operatorname{End}_{\mathcal A}(E)
=\bigl\{\,u\in\operatorname{End}_k(E)\mid ua=au\;\text{对每个}\;a\in\mathcal A\,\bigr\},
\]
称为该作用的\term[zhongxinhua-zi]{中心化子}[centralizer]。它是子代数，因为恒等、和、积都继续与 \(\mathcal A\) 交换。

\ProseParagraphBreak
也可以从一个算子族 \(\mathcal S\) 出发，令 \(\mathcal A\) 为它生成的含单位子代数。算子 \(u\) 若与每个 \(s\in\mathcal S\) 交换，就与这些算子的任意有限乘积及线性组合交换，因而与整个 \(\mathcal A\) 交换；反向由 \(\mathcal S\subseteq\mathcal A\) 立即成立。因此求中心化子时只须检验生成算子。对于张量幂上的两种作用，所要问的是：位置置换生成的代数是否已经包含所有与一般线性群作用交换的算子，反方向又是否成立？

\begin{claim}[Schur–Weyl 对偶]\label{b2:claim:schur-weyl-preview}
设 \(V\) 为有限维复向量空间，\(d\ge1\)。在 \(E=V^{\otimes d}\) 上，令 \(\mathcal A\) 为位置置换算子 \(P_\sigma\) 的复线性生成空间，\(\mathcal B\) 为算子 \(g^{\otimes d}\)、\(g\in\operatorname{GL}(V)\) 的复线性生成空间。此时
\[
\mathcal A=\operatorname{End}_{\mathcal B}(E),\qquad
\mathcal B=\operatorname{End}_{\mathcal A}(E).
\]
\end{claim}

两种线性生成空间确为含单位子代数，因为各自生成算子的乘积仍属于同类算子。\cref{b2:prop:commuting-tensor-actions}及算子线性组合的分配律只给出
\[
\mathcal A\subseteq\operatorname{End}_{\mathcal B}(E),\qquad
\mathcal B\subseteq\operatorname{End}_{\mathcal A}(E).
\]
两个反向包含才构成\term[Schur-Weyl-duiou]{Schur–Weyl 对偶}[Schur–Weyl duality]的实质：它们断言中心化子中没有其他算子。若 \(V=0\)，两等式的两侧均为零；非零情形的证明见第五卷第7.3节，也可参阅 Etingof 等人的《Introduction to Representation Theory》\cite[Theorem 5.18.4, p. 121; Proposition 5.19.1, p. 122]{EtingofEtAl2011RepresentationTheory}。

\ProseParagraphBreak
交换性本身不能推出这些等号。例如在 \(V=k^2\)、\(E=V^{\otimes2}\) 上，把一般线性群的作用换成平凡群的作用，其生成代数便只有 \(\mathcal B=k\operatorname{id}_E\)。位置置换代数为 \(\mathcal A=k\operatorname{id}_E+kP\)，其中 \(P\) 交换两个因子。两种作用当然交换，但 \(\operatorname{End}_{\mathcal B}(E)=\operatorname{End}_k(E)\) 的维数为十六，而 \(\mathcal A\) 的维数至多为二，所以第一个包含严格。

\ProseParagraphBreak
对域 \(k\) 上的向量空间 \(V\)，若 \(2\) 在 \(k\) 中可逆，令 \(P\) 为 \(V^{\otimes2}\) 上交换两个因子的算子，则两个算子
\[
E_+=\frac{\operatorname{id}+P}{2},\qquad
E_-=\frac{\operatorname{id}-P}{2}
\]
满足 \(E_++E_-=\operatorname{id}\)、\(E_\pm^2=E_\pm\)、\(E_+E_-=E_-E_+=0\)，因为 \(P^2=\operatorname{id}\)。这里 \(E_+\) 正是\cref{b2:prop:symmetrization-characteristic}在 \(S_2\) 上的平均算子，\(E_-\) 则是它的互补投影。因此
\[
V^{\otimes2}=\operatorname{im}E_+\oplus\operatorname{im}E_-,
\]
两项分别是 \(P\) 特征值为一的对称张量子空间和特征值为负一的反对称张量子空间；由\cref{b2:prop:commuting-tensor-actions}，它们在一般线性群作用下都稳定。下一章将通过外代数解释第二部分，并专门处理特征二时这两种条件不再分离的情形。


\begin{chaptersummary}
多重张量积由多线性映射的泛性质确定。结合和位置置换都由纯张量上的明确公式给出，生成性保证各种相容等式。张量代数把所有次数作直和并连接张量串，因而适合向任意结合代数指定线性生成元的像；只有底模具有一组自由基时，才得到以所有有限字为基的自由结合代数。

\ProseParagraphBreak
对称代数由交换关系取商，有限自由模的基使它成为多项式代数。增广理想模去其平方恢复一次生成模；这两种代数的构造也与扩张标量相容，无需平坦假设。对称代数的次数分量是对称商，张量幂中的对称张量则是不变子模；阶乘可逆时，平均算子给出两者的自然同构。即使在特征零的整数环上，不可逆的阶乘也会妨碍这个同构；正特征的轨道和还可能在商中消失，取幂也只满足相应的标量幂次规律。

\ProseParagraphBreak
群表示与张量运算相容，对偶作用中的逆元使求值配对保持不变，并保证复合次序正确。有限维条件把 \(V^*\otimes W\) 识别为 \(\operatorname{Hom}(V,W)\)，其不变量正是等变映射。多项式代数 \(k[V]=\operatorname{Sym}_k(V^*)\) 使用同一种对偶作用；其各次不动空间连同乘法组成不变环，有限域上仍须保留形式多项式与点集函数的区别。位置置换与一般线性群的作用交换；Schur–Weyl 对偶在有限维复空间上进一步给出双中心化子等式，其证明归于第五卷第7.3节。下一章改用交替关系，研究外幂与行列式。
\end{chaptersummary}

\BookExercises

\begin{exercise}\label{b2:ex:08-trilinear-coordinates}
设 \(M_1=R^2,M_2=R^3,M_3=R^2\)，\(N\) 为 \(R\) 模。给定十二个元素 \(a_{ijk}\in N\)。证明存在唯一三线性映射，使标准基向量组 \((e_i,f_j,e_k)\) 的像为 \(a_{ijk}\)，写出任意三元组的像。若 \(R=k\) 为域且 \(\dim_kN=s<\infty\)，求全部这些三线性映射组成的空间的维数。
\end{exercise}
\begin{solution}
多线性强迫
\[
b\left(\sum_i x_i e_i,\sum_j y_jf_j,\sum_k z_ke_k\right)
=\sum_{i,j,k}x_iy_jz_k a_{ijk}.
\]
各坐标唯一，因此公式良定义；对每一个变量分别核对加法与标量乘法，得到三线性，所以存在且唯一。由\cref{b2:prop:multiple-tensor-basis}，定义域张量积的基有十二个元素，其线性映射到 \(N\) 可以在这十二个元素上任意取值，故同构于 \(N^{12}\)，维数为 \(12s\)。
\end{solution}

\begin{exercise}\label{b2:ex:08-permutation-convention}
设 \(V\) 有三个不同基向量 \(e_1,e_2,e_3\)，取 \(\sigma=(12)\)、\(\pi=(23)\)。分别计算 \(P_\sigma P_\pi(e_1\otimes e_2\otimes e_3)\) 与 \(P_\pi P_\sigma(e_1\otimes e_2\otimes e_3)\)。若将定义改成按 \(\sigma(1),\ldots,\sigma(d)\) 排列，复合次序将如何变化？
\end{exercise}
\begin{solution}
先 \(P_\pi\) 后 \(P_\sigma\) 得到 \(e_3\otimes e_1\otimes e_2\)；先 \(P_\sigma\) 后 \(P_\pi\) 得到 \(e_2\otimes e_3\otimes e_1\)。两者不同，因为它们是不同的纯基张量。若改用 \(Q_\sigma(\bigotimes_jv_j)=\bigotimes_jv_{\sigma(j)}\)，则 \(Q_\sigma=P_{\sigma^{-1}}\)，故
\[
Q_\sigma Q_\pi=P_{\sigma^{-1}\pi^{-1}}
=P_{(\pi\sigma)^{-1}}=Q_{\pi\sigma}.
\]
于是得到反序的复合法则。置换作用的公式必须与采用的左、右作用约定一起说明。
\end{solution}

\begin{exercise}\label{b2:ex:08-multilinear-dual}
设 \(M_1,\ldots,M_d\) 为有限自由 \(R\) 模。证明
\[
M_1^*\otimes_R\cdots\otimes_RM_d^*
\longrightarrow\operatorname{Hom}_R(M_1\otimes_R\cdots\otimes_RM_d,R)
\]
由 \((\lambda_1\otimes\cdots\otimes\lambda_d)(m_1\otimes\cdots\otimes m_d)=\prod_j\lambda_j(m_j)\) 定义，是自然同构，其中 \(M_j^*=\operatorname{Hom}_R(M_j,R)\)。
\end{exercise}
\begin{solution}
先固定各 \(\lambda_j\)，右端对 \(m_j\) 多线性，故定义张量积上的泛函；所得泛函又对各 \(\lambda_j\) 多线性，所以诱导所写线性映射。对每个 \(M_j\) 选有限基及坐标对偶基。由\cref{b2:prop:multiple-tensor-basis}，纯基张量构成定义域基，而纯对偶基张量的像恰是目标中对应纯基张量的坐标泛函，故映射为同构。给定各 \(f_j:N_j\rightarrow M_j\)，两侧的预复合在纯张量上都取值为 \(\prod_j\lambda_j\bigl(f_j(n_j)\bigr)\)，证明自然性。
\end{solution}

\begin{exercise}\label{b2:ex:08-three-tensor-not-pure}
在 \((k^2)^{\otimes3}\) 中证明 \(t=e_1\otimes e_1\otimes e_1+e_2\otimes e_2\otimes e_2\) 不是纯张量。由此说明任意张量虽能写成纯张量的有限和，也不能总压缩成一项。
\end{exercise}
\begin{solution}
若 \(t=a\otimes b\otimes c\)，写 \(a=a_1e_1+a_2e_2\)，类似写 \(b,c\)。纯基张量的坐标唯一，给出 \(a_1b_1c_1=1\)、\(a_2b_2c_2=1\)，故六个系数都非零。但 \(e_1\otimes e_1\otimes e_2\) 的系数应为零，即 \(a_1b_1c_2=0\)，与域中无零因子矛盾。因此 \(t\) 不是纯张量，而题中已经给出两项纯张量之和。
\end{solution}

\begin{exercise}\label{b2:ex:08-tensor-torsion-module}
对整数 \(n>1\)，证明 \(T_{\mathbb Z}(\mathbb Z/n\mathbb Z)\cong\mathbb Z[X]/(nX)\) 作为分次代数成立，生成元 \(\bar1\) 对应 \(X\) 的像。指出每个次数的分量，以及它为何不是一个自由 \(\mathbb Z\) 模。
\end{exercise}
\begin{solution}
由\cref{b2:prop:tensor-quotient}，每个正次张量幂均为 \(\mathbb Z/n\mathbb Z\)，由 \(\bar1^{\otimes d}\) 生成；零次为 \(\mathbb Z\)。乘法把这两个纯幂连接为更高次纯幂。另一方面，理想 \((nX)\) 恰由常数项为零、所有正次系数均被 \(n\) 整除的多项式组成。因此商多项式的零次系数在 \(\mathbb Z\)，各正次系数在 \(\mathbb Z/n\mathbb Z\)，乘法正好相同。逐次数对应即给出分次同构。一次分量中的非零元素 \(\bar1\) 被 \(n\) 零化，故整个底模有挠；非零自由 \(\mathbb Z\) 模无挠，所以不是自由模。
\end{solution}

\begin{exercise}\label{b2:ex:08-augmentation}
设 \(R\) 为交换含幺环，\(M\) 为幺性 \(R\) 模，\(I\) 为 \(T_R(M)\) 的增广理想。对每个 \(r\ge1\)，求 \(I^r\) 的各次分量，证明 \(I^r/I^{r+1}\cong M^{\otimes r}\)，并证明这些同构使由乘法诱导的映射
\[
(I^r/I^{r+1})\otimes_R(I^s/I^{s+1})\longrightarrow I^{r+s}/I^{r+s+1}
\]
对应于连接张量串的同构，其中 \(r,s\ge1\)。再取 \(M=kx\)、\(R=k\) 为域，证明由 \(x\mapsto x+x^2\) 定义的增广代数自同态在 \(I/I^2\) 上诱导恒等，却不是代数自同构。
\end{exercise}
\begin{solution}
每个 \(I\) 因子只有正次分量，\(r\) 个这样的因子之积没有低于 \(r\) 次的分量，故 \(I^r\subseteq\bigoplus_{d\ge r}M^{\otimes d}\)。反过来，当 \(d\ge r\) 时，每个 \(d\) 重纯张量可以分成 \(r\) 个非空张量串的乘积，例如把前 \(r-1\) 个一次因子分别取出，其余合为最后一串。这些因子均在 \(I\) 中，而纯张量生成整个次数分量，故
\[
I^r=\bigoplus_{d\ge r}M^{\otimes d}.
\]
取次数 \(r\) 的分量，便给出 \(I^r/I^{r+1}\cong M^{\otimes r}\)。如果改变代表元，乘积之差落在 \(I^{r+s+1}\)，所以题述乘法在商上良定义。它在纯张量上连接两个串，由\cref{b2:thm:tensor-associativity-unit}得到所述同构。整个论证只用了纯张量的生成性，不要求 \(M\) 自由。

当 \(M=kx\) 时，\(T_k(M)=k[x]\)。赋值 \(x\mapsto x+x^2\) 给出保增广的自同态，且 \(x+x^2\equiv x\pmod{I^2}\)，所以一次商上的映射为恒等。但若 \(x\) 在此自同态的像中，就有多项式 \(f\) 满足 \(f(x+x^2)=x\)。\(f\) 不可能是常数；若其次数为 \(d\ge1\)，左侧的次数为 \(2d\)，不可能等于一。因此自同态不满射。一次商记录线性部分，并不足以判定整个代数映射是否可逆。
\end{solution}

\begin{exercise}\label{b2:ex:08-tensor-algebra-injection-failure}
令 \(R=k[t]\)，\(M=(t)\subset R\)，其中 \(k\) 为域。写出包含 \(i:M\hookrightarrow R\) 诱导的张量代数同态的生成元公式，证明它单射但不满射，并求其各次分量上的余核。再把系数环扩张到 \(k(t)\)，证明所得代数同态为同构。
\end{exercise}
\begin{solution}
\(M\) 是以 \(t\) 为基的秩一自由 \(R\) 模，而目标 \(R\) 以 \(1\) 为基。把它们在各自张量代数中的一次像分别记为 \(X,Y\)，由\cref{b2:prop:free-associative-algebra}，所求同态写成
\[
R[X]\longrightarrow R[Y],\qquad X\longmapsto tY.
\]
一次生成元 \(X\) 代表模中的基元素 \(t\)，不能把它与零次系数 \(t\) 混同。次数 \(d\) 的映射是 \(RX^d\to RY^d\)、\(aX^d\mapsto at^dY^d\)。\(R\) 为整环，故每个分量上的映射均单射，整个直和上的映射也单射。次数零的余核为零；\(d\ge1\) 时余核为 \(R/(t^d)\)。像是 \(R[tY]\)，而 \(Y\) 不在像中，因为它的一次系数 \(1\) 不属于 \(tR\)，所以映射不满射。

由\cref{b2:prop:tensor-symmetric-base-change}，扩张标量后得到 \(k(t)[X]\to k(t)[Y]\)、\(X\mapsto tY\)。此时 \(t\) 可逆，反向赋值 \(Y\mapsto t^{-1}X\) 给出逆同态。\secref{b2:sec:0802}的双数环反例说明张量代数不普遍保持单射；本题则说明必须检查具体底环、底模及各次张量映射，不能把该反例误读成每个非平凡包含都失去单射性。
\end{solution}

\begin{exercise}\label{b2:ex:08-inhomogeneous-relation}
设 \(k\) 为域，\(A=k[X]/(X^2-X)\)。证明 \(A\) 中 \(1,x\) 为一组基，其中 \(x\) 是 \(X\) 的像；证明这个商不能继承把 \(x\) 放在次数一的多项式分次。
\end{exercise}
\begin{solution}
由首一多项式的带余除法，每个商元素唯一表示为 \(a+bx\)，所以 \(1,x\) 是基，尤其 \(x\ne0\)。若继承原分次，则 \(x\) 在次数一、\(x^2\) 在次数二。关系 \(x^2=x\ne0\) 使同一个非零元素同时落入两个不同次数分量，与直和分解矛盾。商仍然是代数，失去的只是所指定次数的分次。
\end{solution}

\begin{exercise}\label{b2:ex:08-square-zero-algebra}
证明 \(k\langle x,y\rangle/(x^2,xy,yx,y^2)\) 有基 \(1,x,y\)。除了用关系约化说明生成性，还须证明线性无关性。
\end{exercise}
\begin{solution}
任意长度至少二的字都含一个所列二字关系，故在商中为零，得到生成性。要检测这三个元素是否独立，可以构造一个仍保留单位元及两个独立一次元素、而把所有二次乘积都置零的代数。在向量空间 \(B=k\oplus k^2\) 上定义
\[
(a,u)(b,v)=(ab,av+bu).
\]
单位为 \((1,0)\)。两种三项乘积的第二坐标都等于 \(ab w+ac v+bc u\)，第一坐标都为 \(abc\)，故结合；它是 \(k\) 代数。令 \(x\mapsto(0,e_1)\)、\(y\mapsto(0,e_2)\)，所有二字关系的像为零，故\cref{b2:prop:algebra-relations-universal}给出商到 \(B\) 的同态。\(1,x,y\) 的像是 \(B\) 的一组基，所以它们在原商中线性无关。
\end{solution}

\begin{exercise}\label{b2:ex:08-symmetric-quotient}
设 \(N\subseteq M\) 为子模。令 \((N)\) 表示 \(\operatorname{Sym}_R(M)\) 中由 \(N\) 的一次像生成的理想。证明
\[
\operatorname{Sym}_R(M/N)\cong\operatorname{Sym}_R(M)/(N)
\]
为自然分次代数同构。
\end{exercise}
\begin{solution}
商映射 \(M\rightarrow M/N\) 诱导对称代数同态，并将 \(N\) 送到零，故得到右端商代数到左端的同态。反过来，映射 \(M\rightarrow\operatorname{Sym}_R(M)/(N)\) 在 \(N\) 上为零，故公式 \(m+N\mapsto\bar m\) 良定义且线性。由\cref{b2:thm:symmetric-algebra-universal}，它延拓为左端到右端的代数同态。两复合在一次生成元和结构标量上均为恒等，所以互逆。构造保持一次分量和次数，也与保持子模的线性映射相容，故为自然分次同构。
\end{solution}

\begin{exercise}\label{b2:ex:08-symmetric-direct-sum}
对交换 \(R\) 代数 \(A,B\)，先核对 \(A\otimes_RB\) 由 \((a\otimes b)(a'\otimes b')=aa'\otimes bb'\) 成为交换 \(R\) 代数。再证明
\[
\operatorname{Sym}_R(M\oplus N)\cong
\operatorname{Sym}_R(M)\otimes_R\operatorname{Sym}_R(N),
\]
使 \((m,0)\) 对应 \(m\otimes1\)、\((0,n)\) 对应 \(1\otimes n\)。
\end{exercise}
\begin{solution}
四变量公式 \(aa'\otimes bb'\) 对每个变量线性，且把两对因子中的 \(R\) 平衡关系分别送成相等元素，因此诱导张量积上的双线性乘法。在纯张量上核对结合、交换及单位 \(1\otimes1\)，再由生成性推广，便得交换代数。

题述一次映射由 \(M\oplus N\) 的直和性质定义，继而由对称代数泛性质延拓为正向同态。两次包含 \(M,N\rightarrow M\oplus N\) 分别诱导同态 \(\alpha,\beta\) 到左端；公式 \(a\otimes b\mapsto\alpha(a)\beta(b)\) 平衡，且目标交换，故诱导反向代数同态。两复合在 \(M,N\) 的一次像上为恒等；右端由 \(a\otimes1\)、\(1\otimes b\) 生成，而这些又由相应一次像及标量生成，因此两个复合在全代数为恒等。
\end{solution}

\begin{exercise}\label{b2:ex:08-symmetric-dimension}
若 \(\dim_kV=3\)，求 \(V^{\otimes4}\)、\(\operatorname{Sym}_k^4(V)\) 和商映射 \(q_4\) 的核的维数。说明这些数是否依赖域的特征，并给出核中的一个非零元素。
\end{exercise}
\begin{solution}
张量幂维数为 \(3^4=81\)，对称幂维数为 \(\binom64=15\)。\(q_4\) 满射，故核维数为 \(81-15=66\)。前两项的基分别由有序下标与全次数四的单项式给出，都不涉及除法，因此任何特征下相同。取两个不同基向量 \(e_1,e_2\)，则
\[
e_1\otimes e_2\otimes e_1\otimes e_1
-e_2\otimes e_1\otimes e_1\otimes e_1
\]
是两个不同纯基张量之差，非零且在对称商中为零，特征二时也仍非零。
\end{solution}

\begin{exercise}\label{b2:ex:08-symmetrization-orbits}
设 \(k\) 为域，\(V=ke\oplus kf\)。求 \((V^{\otimes4})^{S_4}\) 的一组基，并写出商映射在这组基上相对于 \(e^4,e^3f,e^2f^2,ef^3,f^4\) 的矩阵。分别计算 \(k=\mathbb Q\)、特征二与特征三时的核和余核，并在 \(\mathbb Q\) 上写出 \(e^2f^2\) 对应的平均张量。若把底环改为 \(\mathbb Z\)，求同一限制商映射的核和余核。
\end{exercise}
\begin{solution}
十六个纯基张量按 \(f\) 出现的次数 \(j=0,1,2,3,4\) 分成五个轨道。记 \(u_j\) 为含 \(j\) 个 \(f\) 的全部不同纯基张量之和。不变性要求同一轨道的系数相同，各轨道的支集互不相交，故 \(u_0,\ldots,u_4\) 是不变部分的一组基。第 \(j\) 个轨道有 \(\binom4j\) 项，每项的商像均为 \(e^{4-j}f^j\)，所以矩阵为
\[
\operatorname{diag}(1,4,6,4,1).
\]
在 \(\mathbb Q\) 上这五个数可逆，核与余核均为零。\(e^2f^2\) 对应 \(u_2/6\)：在二十四项平均中，每个不同排列重复 \(2!2!=4\) 次，所以其系数为 \(4/24=1/6\)。

特征二时矩阵为 \(\operatorname{diag}(1,0,0,0,1)\)，核以 \(u_1,u_2,u_3\) 为基，余核以 \(e^3f,e^2f^2,ef^3\) 的类为基，均为三维。特征三时矩阵为 \(\operatorname{diag}(1,1,0,1,1)\)，核为 \(ku_2\)，余核由 \(e^2f^2\) 的类生成，均为一维。这与\cref{b2:exm:cubic-symmetrization-orbits}中三次、特征二时的同构不同：不能只由阶乘不可逆判定某个具体次数的映射。

在整数环上，同一轨道系数的论证仍成立，矩阵的五个对角元都非零，故核为零；逐坐标取商得到余核
\[
\mathbb Z/4\mathbb Z\oplus\mathbb Z/6\mathbb Z\oplus\mathbb Z/4\mathbb Z.
\]
\end{solution}

\begin{exercise}\label{b2:ex:08-invariants-over-integers}
令 \(M=\mathbb Z^r\)、\(r\ge2\)，标准基为 \(e_1,\ldots,e_r\)。求 \((M^{\otimes2})^{S_2}\) 的一组基，并写出限制商映射到 \(\operatorname{Sym}_{\mathbb Z}^2(M)\) 的矩阵。求其核和余核，再求把此矩阵模二后所得线性映射的核与余核的维数。
\end{exercise}
\begin{solution}
交换因子固定每个 \(e_i\otimes e_i\)，并交换 \(e_i\otimes e_j\) 与 \(e_j\otimes e_i\)。不变性要求每对交换项的系数相同，所以不变量以
\[
e_i\otimes e_i\quad(1\le i\le r),\qquad
e_i\otimes e_j+e_j\otimes e_i\quad(1\le i<j\le r)
\]
为基。目标依次以 \(e_i^2\) 和 \(e_ie_j\) 为基时，矩阵为
\[
\operatorname{diag}(\underbrace{1,\ldots,1}_{r\text{ 项}},
\underbrace{2,\ldots,2}_{\binom r2\text{ 项}}).
\]
所有对角元非零，限制映射单射，余核为 \((\mathbb Z/2\mathbb Z)^{\binom r2}\)，由不同 \(e_ie_j\) 的类给出独立的二阶元素。模二后，对角项 \(1\) 保留、\(2\) 消失，故核与余核的维数均为 \(\binom r2\)。章首的秩二计算是这一计算的最小情形；增加基向量后，每一对不同下标都贡献一个不能在整数环上平均的混合项。
\end{solution}

\begin{exercise}\label{b2:ex:08-frobenius-semilinear}
在 \(k=\mathbb F_2[a]/(a^2+a+1)\) 上，令 \(V=ke\)，并定义 \(\Phi_2:V\rightarrow\operatorname{Sym}_k^2(V)\)、\(v\mapsto v^2\)。求 Frobenius 同态 \(F_k:c\mapsto c^2\) 在 \(a\) 上的值，证明 \(\Phi_2\) 可加且对 \(F_k\) 半线性，但不是 \(k\) 线性。若将底域换成 \(\mathbb F_2\)，同一公式有何变化？
\end{exercise}
\begin{solution}
多项式 \(a^2+a+1\) 在 \(\mathbb F_2\) 的两个元素 \(0\) 与 \(1\) 处取值均非零，所以商是四元域，且 \(F_k(a)=a^2=a+1\ne a\)。对称代数交换，特征二下交叉项消失，故 \(\Phi_2(u+v)=\Phi_2(u)+\Phi_2(v)\)；对任意 \(c\in k\) 又有 \(\Phi_2(cv)=c^2\Phi_2(v)=F_k(c)\Phi_2(v)\)。但
\[
(ae)^2=a^2e^2\ne ae^2,
\]
其中 \(e^2\ne0\) 由对称幂的单项式基保证，所以不是 \(k\) 线性。在 \(\mathbb F_2\) 上，每个标量满足 \(c^2=c\)，公式同时保持加法与标量乘法，因而成为线性映射。
\end{solution}

\begin{exercise}\label{b2:ex:08-polynomial-versus-function}
在 \(\mathbb F_3[X,Y]\) 中，证明齐次多项式 \(X^3Y-XY^3\) 非零，但在 \(\mathbb F_3^2\) 的每一点都取零。说明这对把 \(\operatorname{Sym}_{\mathbb F_3}^4\bigl((\mathbb F_3^2)^*\bigr)\) 当作函数空间有什么限制。
\end{exercise}
\begin{solution}
\(X^3Y\) 与 \(XY^3\) 是不同的次数四单项式，系数分别为一和负一，故形式多项式非零。对 \(\mathbb F_3\) 中的每个元素都有 \(a^3=a\)，所以在 \((a,b)\) 处取值为 \(ab-ab=0\)。选取对偶基后，对称幂与这些形式齐次多项式同构，所给非零元素却成为零函数。因此取值映射不单射，不能用函数相等代替对称代数中的元素相等。
\end{solution}

\begin{exercise}\label{b2:ex:08-traces}
设 \(A:V\rightarrow V\)、\(B:W\rightarrow W\) 为有限维空间的线性算子。证明
\(\operatorname{tr}(A\otimes B)=\operatorname{tr}(A)\operatorname{tr}(B)\)。
若 \(2\ne0\) 于 \(k\)，进一步证明
\[
\operatorname{tr}\bigl(\operatorname{Sym}^2(A)\bigr)
=\frac{(\operatorname{tr}A)^2+\operatorname{tr}(A^2)}2.
\]
\end{exercise}
\begin{solution}
在张量基 \(e_i\otimes f_j\) 上，\(A\otimes B\) 的对应对角元为 \(a_{ii}b_{jj}\)，求和得第一式。令 \(P\) 为交换两因子的算子。\(A\otimes A\) 与 \(P\) 交换，因而保持 \(P\) 的两个特征子空间。由\cref{b2:prop:symmetrization-characteristic}，投影 \(E_+=(1+P)/2\) 的像自然识别为对称商。在对称子空间上，\(E_+(A\otimes A)\) 等于 \(\operatorname{Sym}^2(A)\)；在反对称子空间上它为零。两部分的直和就是整个张量空间，所以整个空间上的迹等于对称部分上的迹，即
\[
\operatorname{tr}\bigl(\operatorname{Sym}^2(A)\bigr)
=\operatorname{tr}\bigl(E_+(A\otimes A)\bigr).
\]
\(P(A\otimes A)\) 在 \(e_i\otimes e_j\) 处的对角元为 \(a_{ji}a_{ij}\)，其总和为 \(\operatorname{tr}(A^2)\)。结合第一式并取平均，得到结论，无需假定 \(A\) 可对角化。
\end{solution}

\begin{exercise}\label{b2:ex:08-symmetric-unipotent-matrix}
设 \(A(e)=e\)、\(A(f)=e+f\)。写出 \(\operatorname{Sym}^2(A)\) 在 \(e^2,ef,f^2\) 上的矩阵，求特征二时与特征不为二时的固定子空间维数。
\end{exercise}
\begin{solution}
三列分别由
\[
A(e)^2=e^2,\quad A(e)A(f)=e^2+ef,\quad
A(f)^2=e^2+2ef+f^2
\]
给出，所以矩阵为 \(\begin{pmatrix}1&1&1\\0&1&2\\0&0&1\end{pmatrix}\)。固定向量的坐标 \((a,b,c)\) 满足 \(b+c=0\)、\(2c=0\)。特征不为二时 \(b=c=0\)，固定空间为 \(ke^2\)，维数一；特征二时只剩 \(b=c\)，固定空间由 \(e^2,ef+f^2\) 生成，维数二。对称幂的维数均为三，改变的是诱导算子的固定空间。
\end{solution}

\begin{exercise}\label{b2:ex:08-hom-invariant-jordan}
在 \(V=\mathbb Qe_1\oplus\mathbb Qe_2\) 上，让无穷循环群的生成元作用为
\(\begin{pmatrix}1&1\\0&1\end{pmatrix}\)。
求 \(\operatorname{End}_G(V)\)，并通过 \(\Theta\) 写出 \((V^*\otimes V)^G\) 的一组基。
\end{exercise}
\begin{solution}
设 \(C=\begin{pmatrix}a&b\\c&d\end{pmatrix}\)。与生成元矩阵相交换等价于 \(c=0,d=a\)，所以等变自同态恰为 \(aI+bE_{12}\)，构成二维空间。与生成元相交换也与其正负整数幂相交换，因此已核对全部群元素。由\cref{b2:prop:hom-as-tensor-representation}，\(\Theta\) 将不变张量对应到等变自同态，故所求不变张量的一组基是
\[
e_1^*\otimes e_1+e_2^*\otimes e_2,\qquad e_2^*\otimes e_1.
\]
它们分别对应 \(I\) 和 \(E_{12}\)。
\end{solution}

\begin{exercise}\label{b2:ex:08-commuting-not-centralizers}
设 \(V=\mathbb Q^2\)，\(E=V^{\otimes2}\)，\(P\) 为交换两因子的算子，\(\mathcal A=\operatorname{span}_{\mathbb Q}\{\operatorname{id}_E,P\}\)。利用对称与反对称子空间分解，求 \(\operatorname{End}_{\mathcal A}(E)\) 及其中心化子，证明后者恰为 \(\mathcal A\)。若令平凡群作用在 \(E\) 上，其生成代数为 \(\mathcal B=\mathbb Q\operatorname{id}_E\)，指出 \(\mathcal B\subseteq\operatorname{End}_{\mathcal A}(E)\) 是否为等号。
\end{exercise}
\begin{solution}
记 \(E_+=\ker(P-\operatorname{id})\)、\(E_-=\ker(P+\operatorname{id})\)。若 \(e_1,e_2\) 是 \(V\) 的标准基，则 \(E_+\) 有基
\[
e_1\otimes e_1,\quad e_1\otimes e_2+e_2\otimes e_1,\quad e_2\otimes e_2,
\]
而 \(E_-\) 以 \(e_1\otimes e_2-e_2\otimes e_1\) 为基。与 \(P\) 交换的算子恰为分别保持这两个特征子空间的算子，所以
\[
\mathcal C\coloneqq\operatorname{End}_{\mathcal A}(E)
\cong M_3(\mathbb Q)\oplus\mathbb Q,
\]
其维数为十。两块上的投影均在 \(\mathcal C\) 中；与 \(\mathcal C\) 交换的算子因而也须分块。其三维块须与全部三阶矩阵交换：与对角矩阵单位交换迫使该块为对角矩阵，与非对角矩阵单位交换又迫使三个对角元相等。故
\[
\operatorname{End}_{\mathcal C}(E)
=\{\,a\operatorname{id}_{E_+}\oplus b\operatorname{id}_{E_-}\mid a,b\in\mathbb Q\,\}
=\mathcal A,
\]
最后一个等号由 \(P\) 在两块上分别取 \(1,-1\) 得出。另一方面，\(\mathcal B\) 只有一维，严格小于十维的 \(\mathcal C\)。双中心化子的计算确实恢复了 \(\mathcal A\)，但与它交换的一个任意群作用，并不必生成整个 \(\mathcal C\)。
\end{solution}

\begin{exercise}\label{b2:ex:08-polynomial-weight-invariants}
令无穷循环群 \(G=\langle g\rangle\) 在 \(V=\mathbb Q^2\) 上由 \(g(a,b)=(2a,b/2)\) 作用。求对偶坐标 \(x,y\) 上的作用，证明 \(\mathbb Q[V]^G=\mathbb Q[xy]\)，并求每个次数 \(d\) 的齐次不变空间维数。
\end{exercise}
\begin{solution}
对偶作用使用 \(g^{-1}\)，故 \(gx=x/2,gy=2y\)。单项式 \(x^iy^j\) 的特征值是 \(2^{j-i}\)，在有理数域上恰当 \(i=j\) 时为一。逐项比较系数可知不变多项式正是有限和 \(\sum_i c_i(xy)^i\)。所以次数 \(d\) 为偶数时，不变空间由 \((xy)^{d/2}\) 生成，维数为一；次数为奇数时空间为零。这同时说明生成元 \(xy\) 在原分次下的次数为二。
\end{solution}

\begin{exercise}\label{b2:ex:08-algebra-base-change}
设 \(k\) 为域，\(R=k[\varepsilon]/(\varepsilon^2)\)、\(M=(\varepsilon)\subset R\)、\(S=R/(\varepsilon)=k\)。求 \(T_R(M)\) 与 \(\operatorname{Sym}_R(M)\) 的生成元关系表达，并分别计算它们扩张标量到 \(S\) 后的代数。证明 \(S\otimes_RM\) 非零，再计算包含 \(i:M\hookrightarrow R\) 在扩张标量后诱导的张量代数同态
\[
S\otimes_RT_R(M)\longrightarrow S\otimes_RT_R(R)
\]
的核及像。
\end{exercise}
\begin{solution}
模 \(M\) 由 \(\varepsilon\) 生成。写 \(a=a_0+a_1\varepsilon\in R\)，便有 \(a\varepsilon=a_0\varepsilon\)，故 \(\operatorname{Ann}_R(\varepsilon)=(\varepsilon)\)。因此它的全部线性关系由 \(\varepsilon\cdot\varepsilon=0\) 生成，\(M\cong R/(\varepsilon)\)。把这一模生成元记为代数中的 \(X\)，则
\[
T_R(M)\cong R[X]/(\varepsilon X).
\]
可用泛性质核对：到任意结合 \(R\) 代数 \(A\) 的线性赋值由一个满足 \(\varepsilon a=0\) 的元素 \(a\) 决定；一个生成元的自由结合代数就是 \(R[X]\)，该关系恰好允许赋值下降。这个代数已经交换，故 \(\operatorname{Sym}_R(M)\) 也为同一个商。

由\cref{b2:prop:tensor-quotient}，\(S\otimes_RM\cong M/\varepsilon M=M\)，作为 \(k\) 空间是一维，生成元 \(1\otimes\varepsilon\) 非零。\cref{b2:prop:tensor-symmetric-base-change}于是给出
\[
S\otimes_RT_R(M)\cong k[X],\qquad
S\otimes_R\operatorname{Sym}_R(M)\cong k[X].
\]
原来对一次生成元的 \(\varepsilon\) 关系，在新系数域中已自动成立；零次的 \(\varepsilon\) 也被商去。

另一方面，\(S\otimes_RT_R(R)\cong T_k(k)=k[Y]\)。扩张标量后的包含把 \(1\otimes\varepsilon\) 送到 \(S\otimes_RR\cong S\) 中的零，故诱导代数同态为 \(k[X]\to k[Y]\)、\(X\mapsto0\)。它把多项式送到常数项，核为 \((X)\)，像为常数子代数 \(k\)。扩张标量相容的同构在这里仍成立；失去单射性的是由 \(i\) 诱导的映射。
\end{solution}

% !TeX root = ../../Book2.tex
% 纲要标题：## 第9章　外代数与行列式
% 纲要位置：计划纲要.md，第 1066 行。

\chapter{外代数与行列式}
\label{b2:ch:09}
\BookChapterNumber{b2:ch:09}

\begin{chapterabstract}
外幂把交替多线性映射化为线性映射，外代数则把不同次数的交替运算统一起来。本章先区分交替性与反对称性，在任意交换环上构造外幂并证明有限自由模的外幂基定理；随后从最高外幂定义行列式，推导乘法性与 Cauchy–Binet 公式。伴随矩阵、Cayley–Hamilton 定理、四维空间中的 Plücker 关系以及 Koszul 微分，分别展示外代数在线性算子、子空间坐标和正合性问题中的作用。
\end{chapterabstract}

\begin{learninggoals}
\begin{itemize}
\item 分清交替性与反对称性，说明特征二或目标模中的二挠为何影响它们的等价性。
\item 用泛性质构造外幂和楔积，证明外幂基、分次交换律及直和分解，不借助除以阶乘的公式。
\item 从行列式线推导行列式乘法性与子式公式，并在交换环上证明伴随矩阵恒等式和 Cayley–Hamilton 定理。
\item 在四维向量空间中用 Plücker 坐标判定二向量是否可分解，写出收缩算子和 Koszul 微分，区分复合为零与正合。
\end{itemize}
\end{learninggoals}

在二维自由模的基 \(e_1,e_2\) 下，令 \(v=ae_1+ce_2\)、\(w=be_1+de_2\)。我们希望构造一个双线性运算，记作 \(v\wedge w\)，要求两个相同向量的积为零。这样的交替运算在交换输入时变号，按双线性展开后只剩
\[
v\wedge w=(ad-bc)e_1\wedge e_2.
\]
系数 \(ad-bc\) 正是二阶行列式。此时楔积还只是待构造运算的记号；这个计算说明交替双线性规则会产生面积式的系数。下面用张量积的商实现这种运算，再证明最高外幂上的线性作用恰由这个标量描述。

\ProseParagraphBreak

要把这种计算推广到任意交换环，不能预先把“交换时变号”当作“同一向量的楔积为零”的同义说法：在特征二它们并不等价。从交替关系构造外幂后，基和符号规则都须由这些关系推出。这样得到的运算还可用于多行多列的子式，以及线性算子所满足的多项式关系。

\ProseParagraphBreak

本章需要第六章的张量积泛性质与第八章的多重张量、张量代数。主要构造允许任意交换环；只在把楔积与线性无关、子空间相联系时改用域。行列式的坐标公式由外幂的交替求和得到。

\ProseParagraphBreak
松村英之的《Commutative Ring Theory》\cite[Appendix C, (1)--(6), pp. 283--286]{Matsumura1989CommutativeRingTheory}讨论一般交换环上的外幂、直和与收缩，适合与本章的泛性质证明对照。Roman 的《Advanced Linear Algebra》\cite[Chapter 14, The Determinant, pp. 403--405]{Roman2008AdvancedLinearAlgebra}从多线性形式解释行列式；交替与反对称的定义在不同特征下须加以区别。

% !TeX root = ../../Book2.tex
% 纲要标题：### 9.1 交替多线性映射
% 纲要位置：计划纲要.md，第 1068 行。

\section{交替多线性映射}
\label{b2:sec:0901}

面积在交换两条边时改变符号，在两条边重合时变为零。这两种性质在熟悉的实数情形中相互联系，但在一般系数环上并不等价。外幂所表达的是“有两个输入相同时为零”的交替性。本章始终设 \(R\) 为含单位元的交换环，所有张量积和多线性映射均在 \(R\) 上；需要域或额外特征条件时另行说明。

% 纲要要点（供目录核对）：
%
% * 交替性与反对称性的区别；双线性矩阵的对角线条件
% * 特征二与目标模中的二挠现象
% * 外幂的泛性质；先用张量积表达多线性，再按重复输入取商
% * 外幂不保持一般单射的理想反例，完整证明放在正文
%

\subsection{交替性与反对称性}
\label{b2:subsec:0901:01}

\begin{definition}\label{b2:def:alternating-multilinear}
设 \(R\) 为交换环，\(M,L\) 为 \(R\) 模，\(p\ge1\)。对逐变量 \(R\) 线性的映射 \(b:M^p\to L\)，若任意两个输入相同便有
\[
b(m_1,\ldots,m_p)=0,
\]
就称 \(b\) 为\term[jiaoti-duoxianxing-yingshe]{交替多线性映射}[alternating multilinear map]。若对任意两个位置，交换输入后 \(b\) 的值都变为原值的负元，就称 \(b\) 为\term[fanduichen-duoxianxing-yingshe]{反对称多线性映射}[skew-symmetric multilinear map]。\(p=1\) 时，两种条件均为空条件，此时这两类映射就是线性映射。
\end{definition}

\begin{proposition}\label{b2:prop:alternating-implies-skew}
设 \(R\) 为交换环，\(M,L\) 为 \(R\) 模，\(p\ge1\)。从 \(M^p\) 到 \(L\) 的交替多线性映射必反对称。反过来，若目标模 \(L\) 中 \(2\ell=0\) 蕴含 \(\ell=0\)，则从 \(M^p\) 到 \(L\) 的反对称多线性映射也交替。特别地，当 \(2\) 在 \(R\) 中可逆时，两种条件等价。
\end{proposition}
\begin{proof}
\(p=1\) 时两个条件均为空条件。以下设 \(p\ge2\)。
固定除第 \(i,j\) 个位置以外的所有输入，把这两个位置都取为 \(x+y\)。交替性使函数值为零；按这两个变量展开，两个含 \((x,x)\)、\((y,y)\) 的项分别为零，故剩下的两项满足
\[
b(\ldots,x,\ldots,y,\ldots)
+b(\ldots,y,\ldots,x,\ldots)=0.
\]
这就是反对称性。

若反对称映射在两个位置取相同的 \(x\)，交换它们不改变输入，却使值变号，所以这个值 \(\ell\) 满足 \(2\ell=0\)。所给假设使 \(\ell=0\)，即交替。当 \(2\) 可逆时，乘以 \(2^{-1}\) 就能从 \(2\ell=0\) 推出 \(\ell=0\)。
\end{proof}

将置换写成对换的复合，交替性还给出
\begin{equation}\label{b2:eq:alternating-permutation}
b(m_{\sigma(1)},\ldots,m_{\sigma(p)})
=\operatorname{sgn}(\sigma)b(m_1,\ldots,m_p)
\qquad(\sigma\in S_p).
\end{equation}
所用的符号同态及其在换位上取值为 \(-1\) 的性质，已由第一卷定理3.2.3证明。符号 \(\operatorname{sgn}(\sigma)\) 只由置换决定，因此表达不依赖所选的对换分解；公式在 \(p!\) 不可逆的环上也成立。

\ProseParagraphBreak
例如，任意交换环上的
\[
b\bigl((x_1,x_2),(y_1,y_2)\bigr)=x_1y_2-x_2y_1
\]
是交替双线性映射。其交替性应直接由 \(x_1x_2-x_2x_1=0\) 核对；不能先用“交换变号”，再在可能存在二挠的环上消去一个 \(2\)。

\subsection{特征二与二挠现象}
\label{b2:subsec:0901:02}

在域 \(k=\mathbb F_2\) 上，双线性映射 \(b:k\times k\to k\)、\(b(x,y)=xy\) 满足
\[
b(y,x)=b(x,y)=-b(x,y),
\]
所以反对称，却有 \(b(1,1)=1\)，因而不交替。在特征二中，“交换变号”变成“交换不变”，完全没有强迫对角线取值为零。

\ProseParagraphBreak
这一区别不只由环的特征决定。取 \(R=\mathbb Z\)、\(M=\mathbb Z\)、\(L=\mathbb Z/2\mathbb Z\)，公式 \(b(a,c)=\overline{ac}\) 同样反对称而不交替。环本身特征零，目标模却有二挠。因此，使用\cref{b2:prop:alternating-implies-skew}从反对称性推出交替性时，要检查的是实际取值所在的目标模能否消去 \(2\)，而非仅仅检查底环的特征。

\ProseParagraphBreak
对于交换环 \(R\) 上基为 \(e_1,\ldots,e_n\) 的有限自由模 \(M\)，一个双线性形式 \(b:M\times M\rightarrow R\) 由矩阵 \(A=(a_{ij})\)、\(a_{ij}=b(e_i,e_j)\) 确定。反对称性等价于 \(A^{\mathsf t}=-A\)，而交替性还要求对角线为零，即
\[
a_{ii}=0,\qquad a_{ij}=-a_{ji}.
\]
必要性由交替性及\cref{b2:prop:alternating-implies-skew}得到。反过来，对 \(x=\sum_i x_ie_i\)，展开
\[
b(x,x)=\sum_i a_{ii}x_i^2+
\sum_{i<j}(a_{ij}+a_{ji})x_ix_j=0.
\]
所以条件充分。特征二中，这种矩阵是对称且对角线为零的矩阵；一般对称矩阵并不足够。后续讨论交替双线性型时，仍以对角线取值为零作为交替性的定义。

\subsection{外幂的泛性质}
\label{b2:subsec:0901:03}

前面的区别决定了外幂应当施加哪些关系。如果在二重张量积中只把 \(x\otimes y+y\otimes x\) 置为零，特征二时这只强制交换不变，并不强制 \(x\otimes x=0\)。例如 \(M=\mathbb F_2 e\) 时，所有这样的反对称关系本来就是零，非零的 \(e\otimes e\) 会被保留下来。因此，要表达交替性，必须直接把重复输入所对应的张量置为零。

\begin{definition}\label{b2:def:exterior-power}
设 \(R\) 为交换环，\(M\) 为 \(R\) 模。对 \(p\ge2\)，在 \(p\) 重张量积 \(M^{\otimes p}\) 中，令 \(J_p\) 为所有具有两个相同因子的纯张量生成的子模。将商模
\[
\bigwedge^p M\coloneqq M^{\otimes p}/J_p,
\]
称为 \(M\) 的 \(p\) 次\term[waimi]{外幂}[exterior power]，并把 \(m_1\otimes\cdots\otimes m_p\) 的像记为 \(m_1\wedge\cdots\wedge m_p\)。约定 \(\bigwedge^0M=R\)、\(\bigwedge^1M=M\)，空楔积为 \(1\in R\)。
\end{definition}
\symidx{\(\bigwedge^p M\)：模的第 \(p\) 次外幂}
\symidx{\(m_1\wedge\cdots\wedge m_p\)：外幂中的纯楔积}

多重张量积已经把 \(p\) 个输入的多线性关系转化为模中的线性关系；再模掉 \(J_p\)，就是在这个模中加入“重复输入为零”的关系。这与第三章按关系取模的商、以及第八章通过张量商施加代数关系的构造相同。所得映射 \((m_1,\ldots,m_p)\mapsto m_1\wedge\cdots\wedge m_p\) 多线性且交替，所有这些纯楔积生成外幂。

\begin{proposition}[外幂的泛性质]\label{b2:prop:exterior-alternating-universal}
设 \(R\) 为交换环，\(M,L\) 为 \(R\) 模，\(p\ge1\)。对任意交替 \(p\) 线性映射 \(b:M^p\to L\)，存在唯一 \(R\) 线性映射
\[
\widetilde b:\bigwedge^pM\longrightarrow L,
\qquad \widetilde b(m_1\wedge\cdots\wedge m_p)=b(m_1,\ldots,m_p).
\]
反过来，每个这样的线性映射都给出交替 \(p\) 线性映射。具有这一性质的对象连同其指定映射，在唯一的相容同构意义下确定。
\end{proposition}
\begin{proof}
\(p=1\) 时，\(\bigwedge^1M=M\)，结论就是线性映射本身。以下设 \(p\ge2\)。由\cref{b2:thm:multiple-tensor-universal}，多线性映射 \(b\) 唯一延伸为 \(M^{\otimes p}\to L\)。交替性恰好说明 \(J_p\) 的每个生成元都被送到零，所以它唯一地通过商模分解。反过来，商映射本身交替，与线性映射复合仍然交替。唯一性来自纯楔积的生成性。

若另一个对象也具有同样的性质，两次应用泛性质得到双向线性映射；它们的复合在每个指定纯楔积上为恒等，故处处为恒等，证明唯一同构。\(p=0\) 时把零变量的值看成指定的 \(\ell\in L\)，它唯一对应 \(R\to L\)、\(r\mapsto r\ell\)。
\end{proof}

因此，要定义一个从 \(\bigwedge^pM\) 到 \(L\) 的线性映射，可以先给出 \(M^p\) 上的公式，检查它对各变量线性且在重复输入时为零，再由泛性质取得唯一的线性映射。后面读取外幂坐标、定义楔积及收缩算子时，使用的都是这一方法。

\ProseParagraphBreak
设 \(f:M\to N\) 线性。交替公式
\[
(m_1,\ldots,m_p)\longmapsto
f(m_1)\wedge\cdots\wedge f(m_p)
\]
诱导 \(\bigwedge^p f:\bigwedge^pM\to\bigwedge^pN\)。在纯楔积上计算得到
\begin{equation}\label{b2:eq:exterior-functoriality}
\bigwedge^p(gf)=\left(\bigwedge^p g\right)\left(\bigwedge^p f\right),
\qquad \bigwedge^p\operatorname{id}_M=\operatorname{id}_{\bigwedge^pM}.
\end{equation}
于是每个外幂都是一个函子；\(p=0\) 时规定所有 \(\bigwedge^0 f\) 都为 \(\operatorname{id}_R\)。与张量积相似，纯楔积的表达也不唯一，外幂的一般元素是有限和，未必能写成一个纯楔积。

\ProseParagraphBreak
函子化保证每个线性映射都诱导外幂上的映射，却不保证它保留单射。由张量积取商得到外幂，并没有消除第七章讨论过的这种障碍。

\begin{proposition}\label{b2:prop:exterior-injection-failure}
设 \(k\) 为域，\(R=k[x,y]\)，\(I=(x,y)\) 为 \(R\) 的理想，并把 \(I\) 视为 \(R\) 模。包含映射 \(i:I\hookrightarrow R\) 是单射，但诱导映射
\[
\bigwedge^2i:\bigwedge^2_R I\longrightarrow\bigwedge^2_R R
\]
不是单射：\(x\wedge y\ne0\) 于 \(\bigwedge^2_R I\)，而其像为零。
\end{proposition}
\begin{proof}
为了检测 \(x\wedge y\)，只需构造一个在 \((x,y)\) 处不为零的交替双线性映射。把 \(k\) 视为经 \(R\rightarrow R/I\cong k\)、\(r\mapsto r(0,0)\) 取得标量作用的 \(R\) 模。任意 \(f,g\in I\) 唯一地写成
\[
f=ax+by+f_{\ge2},\qquad g=cx+dy+g_{\ge2},
\qquad a,b,c,d\in k,
\]
其中余下的项全次数至少为二。定义 \(\beta(f,g)=ad-bc\)。读取一次项系数是可加的，所以 \(\beta\) 对两个变量均可加。对任意 \(r\in R\)，\(rf\) 的一次项恰为 \(r(0,0)(ax+by)\)，因此
\[
\beta(rf,g)=r(0,0)\beta(f,g),
\qquad \beta(f,rg)=r(0,0)\beta(f,g).
\]
这就是以 \(k\cong R/I\) 为目标的 \(R\) 双线性。又 \(\beta(f,f)=ab-ba=0\)，所以 \(\beta\) 交替。由\cref{b2:prop:exterior-alternating-universal}，它诱导线性映射 \(\widetilde\beta:\bigwedge^2_RI\rightarrow k\)，满足 \(\widetilde\beta(x\wedge y)=\beta(x,y)=1\)。于是 \(x\wedge y\ne0\)。

另一方面，\(R\) 作为自身的模由 \(1\) 生成；任意纯楔积 \(r\wedge s=rs(1\wedge1)=0\)。纯楔积生成二次外幂，所以 \(\bigwedge^2_R R=0\)。故 \(\bigwedge^2i\) 将上述非零元素送到零，不是单射。
\end{proof}

% !TeX root = ../../Book2.tex
% 纲要标题：### 9.2 外代数
% 纲要位置：计划纲要.md，第 1074 行。

\section{外代数}
\label{b2:sec:0902}

第八章的对称代数把生成元交换前后的差规定为零；外代数则把每个一次元素的平方规定为零，由此得到反对称关系。两者都是张量代数的商，却分别对应对称与交替的多线性运算。即使在特征二中，外代数的平方为零关系也不会消失。

\ProseParagraphBreak
不同次数的外幂不是彼此孤立的模。把输入首尾连接，可以在它们之间定义楔积，组成外代数。有基时，其计算归结为把指标排列成递增次序；没有基时，构造仍由交替映射的泛性质保证。本节把这两种描述联系起来，并核对它们在任意特征下都成立。

% 纲要要点（供目录核对）：
%
% * 基、维数与楔积；用交替坐标函数检测外幂基的系数
% * 分次交换律与奇数次平方为零；代数泛性质要求所有一次元素平方为零
% * 直和的外代数；分次张量积的符号与按次数分解
% * 域上纯楔积非零与线性无关的判据；一般环上非零与成基的区别
%

\subsection{基、维数与楔积}
\label{b2:subsec:0902:01}

设 \(\xi\in\bigwedge^pM\)、\(\eta\in\bigwedge^qM\)。希望用下式定义乘法：
\[
(m_1\wedge\cdots\wedge m_p)\wedge(n_1\wedge\cdots\wedge n_q)
=m_1\wedge\cdots\wedge m_p\wedge n_1\wedge\cdots\wedge n_q.
\]
纯楔积并非自由符号，因而需要先证明这个公式与它们之间的关系相容。

\begin{proposition}\label{b2:prop:wedge-product}
设 \(R\) 为交换环，\(M\) 为 \(R\) 模，\(p,q\ge0\)。规定两个纯楔积的积为
\[
(m_1\wedge\cdots\wedge m_p)\wedge(n_1\wedge\cdots\wedge n_q)
=m_1\wedge\cdots\wedge m_p\wedge n_1\wedge\cdots\wedge n_q.
\]
这一公式唯一延伸为 \(R\) 双线性映射
\[
\bigwedge^pM\times\bigwedge^qM\longrightarrow\bigwedge^{p+q}M,
\qquad(\xi,\eta)\longmapsto\xi\wedge\eta.
\]
这些映射满足结合律，\(R=\bigwedge^0M\) 中的 \(1\) 为单位元。因此
\[
\bigwedge M\coloneqq\bigoplus_{p\ge0}\bigwedge^pM
\]
成为分次 \(R\) 代数，称为 \(M\) 的\term[waidaishu]{外代数}[exterior algebra]；其乘法称为\term[xieji]{楔积}[wedge product]。
\end{proposition}
\begin{proof}
在 \(p=2,q=1\) 时，固定 \(n\in M\)，公式 \((m_1,m_2)\mapsto m_1\wedge m_2\wedge n\) 是交替双线性的。由\cref{b2:prop:exterior-alternating-universal}，它唯一地给出线性映射 \(L_n:\bigwedge^2M\rightarrow\bigwedge^3M\)，满足
\[
L_n(m_1\wedge m_2)=m_1\wedge m_2\wedge n.
\]
对 \(r,s\in R\) 及 \(n,n'\in M\)，\(L_{rn+sn'}=r L_n+s L_{n'}\) 在纯楔积上由多线性成立，故在整个 \(\bigwedge^2M\) 上成立。因此 \((\xi,n)\mapsto L_n(\xi)\) 双线性，且不依赖 \(\xi\) 的纯楔积表达。这就是把前两个变量先合成外幂，再对第三个变量使用一次外幂 \(\bigwedge^1M=M\) 的泛性质。

一般地，先设 \(p,q\ge1\)。固定 \(n_1,\ldots,n_q\)，右端关于 \(m_1,\ldots,m_p\) 交替且多线性，所以对 \(\xi\) 唯一线性延伸。所得公式关于各 \(n_j\) 仍多线性且交替：在纯 \(\xi\) 上已有这些等式，再由线性延伸到所有 \(\xi\)。第二次应用泛性质，即得关于 \(\eta\) 的唯一线性延伸。

若任一次数为零，则定义 \(r\wedge\xi=\xi\wedge r=r\cdot\xi\)，其中 \(r\in R=\bigwedge^0M\)。结合律在三个正次纯楔积上只是同一串因子的首尾连接，三线性与生成性使其在任意三个元素上成立；含零次因子时则由标量作用的结合律成立。\(1\) 因而为单位。直和中的元素只有有限多个非零次数分量，乘法扩张后的求和仍有限，故给出整个直和上的代数结构。
\end{proof}
\symidx{\(\bigwedge M\)：模 \(M\) 的外代数}

\begin{theorem}[有限自由模的外幂基]\label{b2:thm:exterior-power-basis}
设 \(R\) 为交换环，\(M\) 为有基 \(e_1,\ldots,e_n\) 的有限自由 \(R\) 模。对 \(0\le p\le n\)，全部递增指标楔积
\[
e_I\coloneqq e_{i_1}\wedge\cdots\wedge e_{i_p},
\qquad I=\{i_1<\cdots<i_p\}\subseteq\{1,\ldots,n\},
\]
组成 \(\bigwedge^pM\) 的基；\(p=0\) 时只有 \(e_{\varnothing}=1\)。对 \(p>n\)，有 \(\bigwedge^pM=0\)。特别地，在域上
\(\dim_k\bigwedge^p k^n=\binom np\)。
\end{theorem}
\begin{proof}
将各因子按 \(e_i\) 展开，再用多线性，任意纯楔积成为基向量楔积的有限线性组合。含重复指标的项为零，其余项由\eqref{b2:eq:alternating-permutation}化为带符号的递增指标项，故所列元素生成。当 \(p>n\) 时所有项都重复，得到零模。

为证明线性无关，希望对每个 \(e_I\) 构造一个只读取其系数的线性泛函，即在 \(e_I\) 上取一，在其他递增指标楔积上取零。先固定 \(1\le p\le n\) 及一个递增指标组 \(I\)。若 \(v_j=\sum_{a=1}^n x_{aj}e_a\)，每次从第 \(j\) 个输入选一个 \(I\) 中的基向量，按所选指标的排列符号求和，得到
\[
\Delta_I(v_1,\ldots,v_p)
=\sum_{\sigma\in S_p}\operatorname{sgn}(\sigma)
\prod_{j=1}^p x_{i_{\sigma(j)},j}.
\]
每个变量只出现一次，所以它多线性。若第 \(r,s\) 个输入相等，则把置换 \(\sigma\) 与 \(\sigma\circ(r\ s)\) 配对：二者不同，对应乘积相同而符号相反。每对的和为零，所以 \(\Delta_I\) 交替。这一成对消去在特征二中仍然有效。

泛性质给出线性映射 \(\widetilde\Delta_I:\bigwedge^pM\to R\)。对递增指标 \(J\)，当 \(I=J\) 时，求和中只有恒等置换对应的项为一，其余项均为零；当 \(I\ne J\) 时，每项都含有一个为零的坐标。因此 \(\widetilde\Delta_I(e_J)=1\) 当 \(I=J\)，否则为零。于是若 \(\sum_J a_Je_J=0\)，应用 \(\widetilde\Delta_I\) 便得 \(a_I=0\)。这证明线性无关，且没有除以任何整数。\(p=0\) 的结论由 \(\bigwedge^0M=R\) 直接得到。
\end{proof}

例如，在 \(R^3\) 中，二次外幂有基 \(e_1\wedge e_2,e_1\wedge e_3,e_2\wedge e_3\)。直接展开给出
\[
(e_1+2e_2)\wedge(e_2+e_3)
=e_1\wedge e_2+e_1\wedge e_3+2e_2\wedge e_3.
\]
系数 \(2\) 是否为零由所处的环决定，基的构造本身没有因此改变。

\subsection{分次交换律}
\label{b2:subsec:0902:02}

\begin{proposition}[分次交换律]\label{b2:prop:exterior-graded-commutativity}
设 \(R\) 为交换环，\(M\) 为 \(R\) 模，\(p,q\ge0\)。若 \(\xi\in\bigwedge^pM\)、\(\eta\in\bigwedge^qM\)，则
\begin{equation}\label{b2:eq:graded-wedge-sign}
\xi\wedge\eta=(-1)^{pq}\eta\wedge\xi.
\end{equation}
此外，每个奇数次齐次元素 \(\xi\) 都满足 \(\xi\wedge\xi=0\)，无须假设 \(2\) 可逆。
\end{proposition}
\begin{proof}
先取两个纯楔积。把前面的 \(p\) 个因子依次移过后面的 \(q\) 个因子，共交换 \(pq\) 次；每次由\cref{b2:prop:alternating-implies-skew}产生一个负号。这证明纯楔积情形，双线性再给出一般情形。

若 \(p\) 为奇数，分次交换律只直接给出 \(2\xi^2=0\)，在可能有二挠的环上不能据此推出平方为零。需要用外代数中重复因子为零的关系。写 \(\xi=\sum_{i=1}^t a_iw_i\)，其中 \(w_i\) 为 \(p\) 次纯楔积。每个 \(w_i\wedge w_i\) 都含重复因子，所以为零；不同指标的两项合为
\[
a_ia_j(w_i\wedge w_j+w_j\wedge w_i)=0
\]
，因为 \((-1)^{p^2}=-1\)。因此整个平方为零。
\end{proof}

若 \(2=0\) 于 \(R\)，分次交换律中的符号全部成为 \(1\)，外代数作为普通代数也是交换的，但仍保留每个一次元素的平方为零。例如
\[
\bigwedge(\mathbb F_2^2)
\cong\mathbb F_2[x,y]/(x^2,y^2),
\qquad e_1\longleftrightarrow x,\quad e_2\longleftrightarrow y.
\]
双方均由 \(1,x,y,xy\) 或相应楔积组成基，乘法关系一致，因而公式确为同构。它不同于无平方关系的对称代数 \(\mathbb F_2[x,y]\)。

外代数把每个一次元素的平方规定为零，而不只是把某一组基向量的平方规定为零。在线性映射的像中，\(u(x+y)^2=0\) 还要求两个交叉乘积之和为零。因此，即使 \(M\) 自由，仅检查基向量像的平方为零也不足以定义外代数上的同态。

\begin{proposition}[外代数的代数泛性质]\label{b2:prop:exterior-algebra-universal}
设 \(R\) 为交换环，\(M\) 为 \(R\) 模，\(A\) 为含单位元的结合 \(R\) 代数。若 \(R\) 线性映射 \(u:M\to A\) 满足 \(u(m)^2=0\) 对所有 \(m\in M\) 成立，则存在唯一 \(R\) 代数同态 \(\bigwedge M\to A\) 延伸 \(u\)。等价地，外代数是张量代数的商：
\[
\bigwedge M\cong T_R(M)/\langle m\otimes m\mid m\in M\rangle,
\]
分母表示这些二次元素生成的双边理想。
\end{proposition}
\begin{proof}
由假设，\(u(x)^2=u(y)^2=u(x+y)^2=0\)。由于 \(u\) 线性，展开第三个等式得
\[
0=(u(x)+u(y))^2
=u(x)^2+u(x)u(y)+u(y)u(x)+u(y)^2,
\]
因此 \(u(x)u(y)=-u(y)u(x)\)。乘积映射
\((m_1,\ldots,m_p)\mapsto u(m_1)\cdots u(m_p)\) 多线性；若两个输入相同，把其间的因子逐个交换，便出现一个 \(u(m)^2=0\)，所以它交替。逐次应用\cref{b2:prop:exterior-alternating-universal}得到各次外幂上的线性映射。首尾连接的公式保证它们合起来保持乘法和单位，故为代数同态。\(M\) 生成整个外代数，所以延伸唯一。

再由\cref{b2:thm:tensor-algebra-universal}，线性映射 \(M\to\bigwedge M\) 诱导满代数同态 \(T_R(M)\to\bigwedge M\)，并杀掉全部 \(m\otimes m\)，故经过上述商代数。反过来，商代数中的 \(M\) 的像都平方为零，由刚证明的泛性质得到反向代数同态。两复合在 \(M\) 与单位元上为恒等，因生成性而处处为恒等。
\end{proof}

\subsection{直和的外代数}
\label{b2:subsec:0902:03}

\begin{proposition}[直和的外幂]\label{b2:prop:exterior-direct-sum}
设 \(R\) 为交换环，\(M,N\) 为 \(R\) 模，\(p\ge0\)，则存在自然同构
\[
\bigoplus_{a+b=p}\left(\bigwedge^aM\otimes_R\bigwedge^bN\right)
\longrightarrow\bigwedge^p(M\oplus N),
\]
其在纯楔积上的公式为
\[
(m_1\wedge\cdots\wedge m_a)\otimes(n_1\wedge\cdots\wedge n_b)
\longmapsto
(m_1,0)\wedge\cdots\wedge(m_a,0)\wedge(0,n_1)\wedge\cdots\wedge(0,n_b).
\]
\end{proposition}
\begin{proof}
把 \(E=\bigwedge M\otimes_R\bigwedge N\) 按两侧次数之和分次。乘积原来的因子次序是 \(\xi,\eta,\xi',\eta'\)；要把 \(M\) 侧与 \(N\) 侧分别合在一起，就要让 \(N\) 侧次数为 \(b\) 的 \(\eta\) 越过 \(M\) 侧次数为 \(c\) 的 \(\xi'\)。在纯楔积上共发生 \(bc\) 次一次因子交换，所以应有符号 \((-1)^{bc}\)。据此采用带符号的分次张量积乘法，规定
\[
(\xi\otimes\eta)(\xi'\otimes\eta')
=(-1)^{bc}(\xi\wedge\xi')\otimes(\eta\wedge\eta'),
\qquad \deg\eta=b,\quad\deg\xi'=c.
\]
固定次数后此式分别线性且对两次张量关系平衡，故良定义。三因子结合律中，两种括号的符号指数分别为
\(bc+(b+d)e\) 与 \(de+b(c+e)\)，其中 \(d=\deg\eta'\)、\(e=\deg\xi''\)，二者相同；各外代数的结合律保证剩余因子也相同。单位是 \(1\otimes1\)。

映射 \(M\oplus N\to E\)、\((m,n)\mapsto m\otimes1+1\otimes n\) 的每个像都平方为零：两个同侧平方为零，两项交叉乘积分别为 \(m\otimes n\) 与 \(-m\otimes n\)。因此\cref{b2:prop:exterior-algebra-universal}给出分次代数同态 \(\alpha:\bigwedge(M\oplus N)\to E\)。

反方向，把 \(\xi\otimes\eta\) 送到两侧包含所诱导的楔积，得到线性映射 \(\beta:E\to\bigwedge(M\oplus N)\)。将中间的 \(\eta\) 与 \(\xi'\) 交换时，由\eqref{b2:eq:graded-wedge-sign}恰好得到 \((-1)^{bc}\)，所以 \(\beta\) 也是代数同态。两复合在所有 \(m\in M,n\in N\) 及单位元上都为恒等；这些元素生成双方代数，故 \(\alpha,\beta\) 互逆。取全次数 \(p\) 的分量，就得到命题中的直和同构。对 \(M,N\) 的线性映射，公式逐因子作用，故同构自然。
\end{proof}

一个 \(p\) 重楔积的输入来自 \(M\oplus N\)。将每个输入写成 \((m,n)=(m,0)+(0,n)\) 后按多线性展开，每一项恰有 \(a\) 个因子来自 \(M\)、\(b\) 个来自 \(N\)，其中 \(a+b=p\)。把 \(M\) 因子移到前面、\(N\) 因子移到后面，就落在上述直和的 \((a,b)\) 分量。因此直和公式是按两侧各取了多少个因子分类。特别地，
\[
\bigwedge^2(M\oplus N)
\cong\bigwedge^2M\oplus(M\otimes_RN)\oplus\bigwedge^2N.
\]
例如 \((m,0)\wedge(0,n)\) 属于中间项，而反向次序对应 \(-m\otimes n\)。这项负号已经由外代数的乘法决定，不是对直和分解额外指定的约定。若 \(M,N\) 分别具有 \(m,n\) 个基向量的有限自由模，比较基的数目便得到
\(\binom{m+n}{p}=\sum_{a+b=p}\binom ma\binom nb\)，其代数来源正是把一个递增指标组按两组基拆开。

\ProseParagraphBreak
在域上，纯楔积的非零性还能检测其因子是否独立。这里不需要先假定整个向量空间有限维：有限个因子只涉及一个有限维生成空间，再用直和公式比较它与母空间中的外幂。

\begin{proposition}\label{b2:prop:wedge-linear-independence}
设 \(k\) 为任意域，\(V\) 为 \(k\) 向量空间，\(p\ge1\)，\(v_1,\ldots,v_p\in V\)。则
\[
v_1\wedge\cdots\wedge v_p\ne0
\quad\Longleftrightarrow\quad
v_1,\ldots,v_p\text{ 线性无关}.
\]
\end{proposition}
\begin{proof}
若这些向量线性相关，取关系 \(\sum_{i=1}^p c_iv_i=0\) 中的非零系数 \(c_j\)，由域中可除性把 \(v_j\) 写成其余向量的线性组合。代入楔积并按多线性展开，每一项都含重复因子，因而为零。\(p=1\) 时这就是说 \(v_1=0\)，结论同样成立。

若这些向量线性无关，令 \(W=\operatorname{span}_k\{v_1,\ldots,v_p\}\)。它们构成有限维空间 \(W\) 的一组基，由\cref{b2:thm:exterior-power-basis}，\(v_1\wedge\cdots\wedge v_p\) 在 \(\bigwedge^pW\) 中是非零的基向量。由\cref{b2:thm:vector-basis-extension}将它们扩充为 \(V\) 的基，令 \(U\) 为其余基向量的线性生成空间，则 \(V=W\oplus U\)，即使 \(V\) 无限维也成立。\cref{b2:prop:exterior-direct-sum}给出
\[
\bigwedge^pV\cong
\bigoplus_{a+b=p}\bigl(\bigwedge^aW\otimes_k\bigwedge^bU\bigr).
\]
由该同构的纯楔积公式，\(W\hookrightarrow V\) 诱导的 \(\bigwedge^pW\rightarrow\bigwedge^pV\)，恰把 \(\bigwedge^pW\cong\bigwedge^pW\otimes_k k\) 送入 \((a,b)=(p,0)\) 的直和分量。因此它是单射，上述纯楔积在 \(\bigwedge^pV\) 中仍非零。
\end{proof}

在域上，独立向量还可以扩充成基；一般环上，非零楔积并不保证其因子能够扩充成母模的一组基。例如在 \(\mathbb Z^2\) 中，
\[
(2e_1)\wedge e_2=2(e_1\wedge e_2)\ne0,
\]
因为 \(e_1\wedge e_2\) 是自由 \(\mathbb Z\) 模 \(\bigwedge^2\mathbb Z^2\) 的基向量。但 \(2e_1,e_2\) 只能生成第一坐标为偶数的子模，不能生成整个 \(\mathbb Z^2\)。它们虽独立，楔积中的系数 \(2\) 却不是单位；“系数非零”与“系数可逆”的差别将在行列式的可逆性判据中再次出现。

% !TeX root = ../../Book2.tex
% 纲要标题：### 9.3 行列式的内在定义
% 纲要位置：计划纲要.md，第 1080 行。

\section{行列式的内在定义}
\label{b2:sec:0903}

一个 \(n\) 维空间中的 \(n\) 重楔积只剩下一条线，所以线性变换在这条线上只能表现为乘以一个标量。这个标量就是行列式。以最高外幂为起点，乘法性来自函子复合，子式公式来自较低次外幂的坐标。本节同时允许任意交换环上的有限自由模，不以域上的除法或特征值作为前提。

% 纲要要点（供目录核对）：
%
% * 最高外幂与行列式线
% * 行列式的乘法性
% * Cauchy–Binet 公式的外代数解释
%

\subsection{最高外幂与行列式线}
\label{b2:subsec:0903:01}

\begin{definition}\label{b2:def:determinant-line}
设 \(R\) 为交换环，\(M\) 为有 \(n\) 个基向量的有限自由 \(R\) 模。它的最高外幂
\[
\det M\coloneqq\bigwedge^nM
\]
是秩一自由 \(R\) 模，称为 \(M\) 的\term[hanglieshi-xian]{行列式线}[determinant line]。\(M=0\) 且取空基时，约定 \(\det M=R\)。
\end{definition}
\symidx{\(\det M=\bigwedge^nM\)：有限自由模的行列式线}

由\cref{b2:thm:exterior-power-basis}，任一有序基 \(e_1,\ldots,e_n\) 给出行列式线的基向量 \(\omega=e_1\wedge\cdots\wedge e_n\)。对于线性自同态 \(f:M\to M\)，存在唯一 \(a\in R\) 使
\begin{equation}\label{b2:eq:intrinsic-determinant}
\left(\bigwedge^n f\right)(\omega)=a\omega.
\end{equation}
这个 \(a\) 不依赖 \(\omega\) 的选择。事实上，同一自由模的另一个单元素基必为 \(\omega'=u\omega\)，其中 \(u\) 为单位：若 \(\omega=v\omega'\)，则由基坐标唯一性得 \(uv=1\)。现在
\(\bigwedge^n f(\omega')=u a\omega=a\omega'\)，使用的正是 \(R\) 的交换性。

\ProseParagraphBreak
据此定义 \(f\) 的\term[hanglieshi]{行列式}[determinant]为 \(\det f=a\)。对 \(n\times n\) 矩阵 \(A=(a_{ij})\)，将其看成 \(R^n\) 上的线性映射并记其行列式为 \(\det A\)。把列向量 \(f(e_j)=\sum_i a_{ij}e_i\) 代入\eqref{b2:eq:intrinsic-determinant}，多线性展开后只有没有重复指标的项保留，故
\begin{equation}\label{b2:eq:determinant-leibniz}
\det A=\sum_{\sigma\in S_n}\operatorname{sgn}(\sigma)
\prod_{j=1}^n a_{\sigma(j),j}.
\end{equation}
这与熟悉的 Leibniz 公式相同；本节的定义没有改变原来的行列式，而是解释了它为何具有坐标无关的含义。空矩阵的行列式为 \(1\)，与空楔积的约定一致。
\symidx{\(\det f,\det A\)：线性自同态及矩阵的行列式}

\ProseParagraphBreak
公式立即表明行列式对各列分别线性，若两列相同则为零，交换两列变号。对转置矩阵，用置换 \(\sigma^{-1}\) 重排求和，并利用标量乘法交换，得到
\begin{equation}\label{b2:eq:determinant-transpose}
\det(A^{\mathsf t})=\det A.
\end{equation}
因此相同性质也适用于行。上三角矩阵的行列式等于对角元之积：非恒等置换必有某个 \(\sigma(j)>j\)，相应乘积含下三角的零元，只有恒等置换可能贡献。

\subsection{行列式的乘法性}
\label{b2:subsec:0903:02}

\begin{theorem}[行列式的乘法性]\label{b2:thm:determinant-multiplicativity}
设 \(R\) 为交换环，\(M\) 为有限自由 \(R\) 模，\(f,g:M\to M\) 为 \(R\) 线性自同态，则
\[
\det(gf)=\det(g)\det(f),\qquad
\det(\operatorname{id})=1.
\]
特别地，\(n\times n\) 矩阵满足 \(\det(BA)=\det B\det A\)。若 \(P\) 可逆，则
\[
\det(P^{-1}AP)=\det A.
\]
\end{theorem}
\begin{proof}
对行列式线的基 \(\omega\)，由\eqref{b2:eq:exterior-functoriality}，
\[
\begin{aligned}
\left(\bigwedge^n(gf)\right)(\omega)
&=\left(\bigwedge^n g\right)\left(\bigwedge^n f\right)(\omega)\\
&=\det(f)\det(g)\omega.
\end{aligned}
\]
基系数唯一，得到乘法公式；恒等映射保持 \(\omega\)，得到第二式。若 \(P\) 可逆，乘法性给出 \(\det(P^{-1})\det P=1\)，所以相似矩阵公式由三次乘法直接推出。
\end{proof}

在域上，\(\det f\ne0\) 等价于 \(f\) 可逆：若像中的 \(n\) 个向量线性相关，其楔积为零；若线性无关，它们构成一组基，从而 \(f\) 可逆。这里“相关则楔积为零”由把一个向量写成其他向量的线性组合得到；反方向把无关组扩充为基，再用\cref{b2:thm:exterior-power-basis}可知其楔积非零。

\ProseParagraphBreak
在一般交换环上，非零行列式不保证可逆。例如 \(\mathbb Z\) 上乘二的行列式为 \(2\ne0\)，但乘二不满射。正确条件是行列式为单位；必要性已经由乘法性得到，充分性将在\cref{b2:prop:adjugate-identity}中用伴随矩阵证明。单射、满射与可逆也不宜再用域上的秩直觉一概替代。

\subsection{Cauchy–Binet 公式的外代数解释}
\label{b2:subsec:0903:03}

设 \(A:R^n\to R^m\) 的矩阵仍记为 \(A\)。对递增指标集 \(I\subseteq\{1,\ldots,m\}\)、\(J\subseteq\{1,\ldots,n\}\)，且 \(|I|=|J|=p\)，记 \(A_{I,J}\) 为对应的 \(p\times p\) 子矩阵；它的行列式称为一个 \(p\) 阶\term[zishi]{子式}[minor]。基向量楔积的展开与\eqref{b2:eq:determinant-leibniz}给出
\begin{equation}\label{b2:eq:exterior-minor-action}
\left(\bigwedge^p A\right)(e_J)
=\sum_{|I|=p}\det(A_{I,J})e_I.
\end{equation}
这里目标与来源分别使用各自的标准基；每个系数都是选定 \(p\) 列后，对相应 \(p\) 行坐标所作的交替求和。因此，较低次外幂把所有同阶子式同时组织成一个线性映射。

\begin{theorem}[Cauchy–Binet]\label{b2:thm:cauchy-binet}
设 \(R\) 为交换环，\(m,n,r\ge0\)，\(A\in M_{m\times n}(R)\)、\(B\in M_{n\times r}(R)\)，\(p\) 为满足 \(0\le p\le\min(m,r)\) 的整数。对大小均为 \(p\) 的递增指标集 \(I\subseteq\{1,\ldots,m\}\)、\(J\subseteq\{1,\ldots,r\}\)，有
\begin{equation}\label{b2:eq:cauchy-binet}
\det\bigl((AB)_{I,J}\bigr)
=\sum_{\substack{K\subseteq\{1,\ldots,n\}\\|K|=p}}
\det(A_{I,K})\det(B_{K,J}).
\end{equation}
若 \(p>n\)，右端为空和，等于零。特别地，\(A\) 为 \(p\times n\)、\(B\) 为 \(n\times p\) 矩阵时，此式给出 \(\det(AB)\) 的全部最大子式展开。
\end{theorem}
\begin{proof}
先对 \(\bigwedge^pB\) 使用\eqref{b2:eq:exterior-minor-action}，再对每个结果使用 \(\bigwedge^pA\)，得到
\[
\left(\bigwedge^pA\right)\left(\bigwedge^pB\right)(e_J)
=\sum_I\left(\sum_K\det(A_{I,K})\det(B_{K,J})\right)e_I.
\]
另一方面，函子性说明左边等于 \(\bigwedge^p(AB)(e_J)\)，其 \(e_I\) 系数是 \(\det\bigl((AB)_{I,J}\bigr)\)。外幂基的坐标唯一性给出公式。若 \(p>n\)，中间模 \(\bigwedge^p R^n\) 为零，复合映射为零，故全部这些子式都为零。
\end{proof}

这一公式称为\term[Cauchy-Binet-gongshi]{Cauchy–Binet 公式}[Cauchy--Binet formula]。例如取
\[
A=\begin{pmatrix}1&0&1\\0&1&1\end{pmatrix},\qquad
B=\begin{pmatrix}1&2\\0&1\\1&0\end{pmatrix}.
\]
直接相乘得到 \(AB=\begin{pmatrix}2&2\\1&1\end{pmatrix}\)，行列式为零。按列指标 \(12,13,23\)，\(A\) 的三个二阶子式为 \(1,1,-1\)，\(B\) 对应行指标的子式为 \(1,-2,-1\)，故 Cauchy–Binet 求和为 \(1-2+1=0\)。这个算例同时核对了指标次序和符号。

\ProseParagraphBreak
当 \(M,N\) 分别具有 \(m,n\) 个基向量时，\cref{b2:prop:exterior-direct-sum}在次数 \(m+n\) 只剩一项，给出
\[
\det(M\oplus N)\cong\det M\otimes_R\det N.
\]
在这个同构下，\(f\oplus g\) 的最高外幂是两条行列式线上的作用之张量积，所以块对角矩阵的行列式为两个块的行列式之积。这也是直和外代数公式在最高次数的直接应用。

% !TeX root = ../../Book2.tex
% 纲要标题：### 9.4 外代数的应用
% 纲要位置：计划纲要.md，第 1086 行。

\section{外代数的应用}
\label{b2:sec:0904}

外代数既解释行列式，也提供记录子空间与构造复形的方法。本节先从子式推出伴随矩阵公式，并在任意交换环上证明 Cayley–Hamilton 定理；然后在域上用一个四维例子说明 Plücker 关系；最后给出 Koszul 微分的构造。三个应用分别用到最高外幂、二次外幂和整个分次外代数。

% 纲要要点（供目录核对）：
%
% * 余子式、伴随矩阵与 Cayley–Hamilton；系数递推和相邻项抵消
% * Plücker 坐标的初步例子；非零纯楔积恢复子空间、四维二次关系与特征二
% * Koszul 复形的外代数背景；收缩由分次乘法规则产生，反交换关系完整证明
% * 三变量自然关系与关系之间的相容性；复合为零不等于正合，自由消解的含义
% * 含单位序列的 Koszul 复形正合，构造及证明放在正文；正则序列理论归第三卷
%
% ---
%

\subsection{余子式、伴随矩阵与 Cayley–Hamilton 定理}
\label{b2:subsec:0904:01}

设 \(R\) 为交换环，\(A=(a_{ij})\in M_n(R)\)、\(n\ge1\)。删去第 \(i\) 行、第 \(j\) 列所得子矩阵记为 \(A^{\widehat i,\widehat j}\)。数
\[
C_{ij}=(-1)^{i+j}\det A^{\widehat i,\widehat j}
\]
称为对应的\term[daishu-yuzishi]{代数余子式}[cofactor]。当 \(n=1\) 时，删去后是空矩阵，行列式取 \(1\)。定义
\[
\operatorname{adj}(A)=(C_{ji})_{i,j}
\]
为 \(A\) 的伴随矩阵（adjugate matrix）\termidx[ban-sui-juzhen]{伴随矩阵}。这是代数余子式矩阵的转置，与第四章由首一多项式构造的友矩阵（companion matrix）是不同概念。
\symidx{\(\operatorname{adj}(A)\)：矩阵 \(A\) 的伴随矩阵}

\begin{proposition}[余子式展开与伴随矩阵]\label{b2:prop:adjugate-identity}
设 \(R\) 为交换环，\(A=(a_{ij})\in M_n(R)\)、\(n\ge1\)。令 \(C_{ij}=(-1)^{i+j}\det A^{\widehat i,\widehat j}\) 为代数余子式，\(\operatorname{adj}(A)=(C_{ji})_{i,j}\) 为伴随矩阵，其中 \(n=1\) 时空矩阵的行列式取 \(1\)。固定任意行 \(i\) 或列 \(j\)，有
\[
\det A=\sum_{k=1}^n a_{ik}C_{ik}
=\sum_{k=1}^n a_{kj}C_{kj}.
\]
进而
\begin{equation}\label{b2:eq:adjugate-identity}
A\operatorname{adj}(A)=\operatorname{adj}(A)A=(\det A)I_n.
\end{equation}
所以 \(A\) 可逆当且仅当 \(\det A\) 是 \(R\) 的单位，此时
\(A^{-1}=(\det A)^{-1}\operatorname{adj}(A)\)。
\end{proposition}
\begin{proof}
在\eqref{b2:eq:determinant-leibniz}中，按某一固定行所使用的列位置分组。含 \(a_{ik}\) 的项，其余因子的排列恰好遍历删去该行、该列后的置换。把第 \(i\) 行和第 \(k\) 列分别移到首位，共产生 \(i-1+k-1\) 个换位，故这一组的符号因子为 \((-1)^{i+k}\)。组内求和于是等于 \(a_{ik}C_{ik}\)，得到行展开。对转置矩阵使用行展开，再用\eqref{b2:eq:determinant-transpose}，得到列展开。

乘积 \(A\operatorname{adj}(A)\) 的 \((i,j)\) 项为 \(\sum_k a_{ik}C_{jk}\)。当 \(i=j\) 时是 \(\det A\)。当 \(i\ne j\) 时，我们仍要把这个和解释成一次行列式展开：把 \(A\) 的第 \(j\) 行替换为第 \(i\) 行，沿第 \(j\) 行展开时，代数余子式仍为 \(C_{jk}\)，行中的系数则变为 \(a_{ik}\)。所得行列式正是上述和，但矩阵有两行相同，故为零。这证明第一个乘积公式。第二个乘积公式通过替换第 \(i\) 列为第 \(j\) 列并作列展开得到。

若行列式为单位，式\eqref{b2:eq:adjugate-identity}给出双侧逆。若 \(A\) 可逆，则\cref{b2:thm:determinant-multiplicativity}给出 \(\det A\det A^{-1}=1\)，所以行列式为单位。
\end{proof}

下面通过比较多项式的系数重新证明 Cayley–Hamilton 定理\termidx[Cayley-Hamilton-dingli]{Cayley--Hamilton 定理}。\cref{b2:prop:operator-characteristic-minimal}利用域上的多项式环模分类，证明特征多项式是最小多项式的倍数，因而零化算子；这里直接使用行列式与伴随矩阵恒等式，把系数推广到任意交换环。

\begin{theorem}[Cayley–Hamilton]\label{b2:thm:exterior-cayley-hamilton}
若 \(R\) 为交换环、\(A\in M_n(R)\)、\(n\ge1\)，令
\[
\chi_A(t)=\det(tI_n-A)=t^n+c_{n-1}t^{n-1}+\cdots+c_0.
\]
则 \(\chi_A(A)=A^n+c_{n-1}A^{n-1}+\cdots+c_0I_n=0\)。
\end{theorem}
\begin{proof}
在交换环 \(R[t]\) 上应用\cref{b2:prop:adjugate-identity}，得到
\[
(tI_n-A)\operatorname{adj}(tI_n-A)=\chi_A(t)I_n.
\]
对标量多项式 \(\chi_A(t)\)，\(\chi_A(A)\) 当然有意义。但若要在左端按普通多项式代入的规则令 \(t=A\)，还须先核对矩阵系数与 \(A\) 交换：在矩阵系数多项式中，\(t\) 与所有系数交换，换成 \(A\) 后却不能未经证明地保留这一性质。因此这里保留每个矩阵因子的次序，改由系数关系消去伴随矩阵中的未知系数。

伴随矩阵的每个条目是 \(n-1\) 阶子式，所以次数至多 \(n-1\)。将它按 \(t\) 展开为
\[
\operatorname{adj}(tI_n-A)=B_0+B_1t+\cdots+B_{n-1}t^{n-1},
\qquad B_j\in M_n(R).
\]
比较 \(t\) 的各次系数，依次得到
\[
-AB_0=c_0I_n,\qquad
B_{j-1}-AB_j=c_jI_n\ (1\le j<n),\qquad B_{n-1}=I_n.
\]
这些关系把 \(c_jI_n\) 写成相邻两个 \(B_j\) 的组合，不必逐个求出 \(B_j\)。将中间第 \(j\) 个等式在左侧乘以 \(A^j\)，最后一式在左侧乘以 \(A^n\)，再连同首式求和，相邻的 \(B_j\) 项便有相同的左侧幂次而符号相反。右端为 \(\chi_A(A)\)，左端为
\[
-AB_0+\sum_{j=1}^{n-1}(A^jB_{j-1}-A^{j+1}B_j)
+A^nB_{n-1}=0.
\]
相邻项逐一抵消，故 \(\chi_A(A)=0\)。这一计算保持每个矩阵系数的位置，没有使用矩阵环中任意元素交换的假设。
\end{proof}

例如，对 \(A=\begin{pmatrix}a&b\\c&d\end{pmatrix}\)，有
\[
\chi_A(t)=t^2-(a+d)t+(ad-bc),
\qquad A^2-(a+d)A+(ad-bc)I_2=0.
\]
这些等式在 \(\mathbb Z\)、多项式环以及含幂零元的交换环上都成立。证明中用到的交换性属于标量环，而不是所有 \(2\times2\) 矩阵之间的交换性。

\subsection{Plücker 坐标}
\label{b2:subsec:0904:02}

本小节改设 \(k\) 为域、\(V=k^n\)。对 \(r\) 维子空间 \(W\subseteq V\)，选一组有序基 \(w_1,\ldots,w_r\)，写
\[
w_1\wedge\cdots\wedge w_r=\sum_{|I|=r}p_Ie_I.
\]
这些 \(p_I\) 是以 \(w_j\) 为列的矩阵的 \(r\) 阶子式。若换一组基，整个楔积乘以换基矩阵的非零行列式。因此，把所有 \(p_I\) 同时乘以同一个非零标量视为同一组坐标，就得到只依赖 \(W\) 的\term[Plucker-zuobiao]{Plücker 坐标}[Plücker coordinates]。这种只确定到共同非零倍数的坐标称为齐次坐标；子空间 \(W\) 确定的是 \(\bigwedge^rV\) 中的一条直线 \(k(w_1\wedge\cdots\wedge w_r)\)，而不是这条直线上的某个指定非零向量。

\ProseParagraphBreak
这些坐标确实能够辨认子空间。对非零纯楔积 \(\omega=w_1\wedge\cdots\wedge w_r\)，有
\begin{equation}\label{b2:eq:plucker-recovers-plane}
W=\{\,v\in V\mid v\wedge\omega=0\,\}.
\end{equation}
若 \(v\in W\)，按 \(w_j\) 展开即得楔积为零。若 \(v\notin W\)，则 \(v,w_1,\ldots,w_r\) 线性无关，\cref{b2:prop:wedge-linear-independence}说明它们的楔积非零。因此两个非零纯楔积只相差标量时，恢复出的子空间相同。这里的非零与纯楔积条件都不可省略；对一般外张量，右端不能无条件解释为一个 \(r\) 维子空间。\cref{b2:ex:09-plucker-plane-recovery}给出了由具体坐标恢复平面的计算。

\ProseParagraphBreak
还需判定哪些坐标真的来自一个子空间。例如，若一个 \(4\times2\) 矩阵的两列为 \(u,v\)，它的六个二阶子式同时来自这两列，因而必须满足彼此的相容关系。下面的二次式正是这种约束。

\begin{proposition}[四维空间中二重楔积的判据]\label{b2:prop:plucker-four-dimensional}
设 \(k\) 为任意域，\(e_1,\ldots,e_4\) 为 \(k^4\) 的标准基，并设
\[
\omega=\sum_{1\le i<j\le4}p_{ij}e_i\wedge e_j\in\bigwedge^2k^4.
\]
它可以写为 \(u\wedge v\)，当且仅当
\begin{equation}\label{b2:eq:plucker-quadric}
p_{12}p_{34}-p_{13}p_{24}+p_{14}p_{23}=0.
\end{equation}
此结论对任意特征的域成立。\(\omega\ne0\) 时，\(u,v\) 自动线性无关。
\end{proposition}
\begin{proof}
若 \(u=\sum u_ie_i\)、\(v=\sum v_ie_i\)，则 \(p_{ij}=u_iv_j-u_jv_i\)。代入\eqref{b2:eq:plucker-quadric}，得到
\[
\begin{aligned}
&(u_1v_2-u_2v_1)(u_3v_4-u_4v_3)\\
&\quad-(u_1v_3-u_3v_1)(u_2v_4-u_4v_2)\\
&\quad+(u_1v_4-u_4v_1)(u_2v_3-u_3v_2)=0.
\end{aligned}
\]
展开后的十二项按相同的 \(u_iu_jv_av_b\) 配对，符号相反，所以为零。

反过来，\(\omega=0\) 时取 \(u=0\)。若 \(\omega\ne0\)，至少一个 \(p_{ij}\ne0\)。重排基可使它成为 \(p_{12}\)：对相邻交换 \((1\ 2),(2\ 3),(3\ 4)\)，直接将下标重排并用 \(p_{ji}=-p_{ij}\)，可见关系左边都变为其负元，所以关系在重排后仍成立。令 \(a=p_{12}\ne0\)。我们先把待求的两列 \(u,v\) 的前两行选为 \((1,0)\)、\((0,a)\)，使 \(12\) 子式已经等于 \(a\)：
\[
\begin{pmatrix}
1&0\\
0&a\\
u_3&v_3\\
u_4&v_4
\end{pmatrix}.
\]
再匹配四个包含第 \(1\) 或第 \(2\) 行的子式，就得到
\[
v_3=p_{13},\qquad v_4=p_{14},\qquad
-a u_3=p_{23},\qquad -a u_4=p_{24}.
\]
前两行使这四个子式成为关于未知系数的线性方程。由于 \(a\ne0\)，它们唯一确定 \(u_3,u_4,v_3,v_4\)，即
\[
u=e_1-\frac{p_{23}}a e_3-\frac{p_{24}}a e_4,
\qquad v=a e_2+p_{13}e_3+p_{14}e_4.
\]
展开 \(u\wedge v\)，\(12,13,14,23,24\) 五个坐标分别就是给定的 \(p_{ij}\)，而 \(34\) 坐标为
\[
\frac{p_{13}p_{24}-p_{14}p_{23}}a=p_{34},
\]
最后一步恰用\eqref{b2:eq:plucker-quadric}：五个坐标已由所选两列实现，二次关系保证剩下的第六个坐标也匹配。故 \(u\wedge v=\omega\)。若 \(u,v\) 相关，其楔积为零，所以非零情形下它们独立。
\end{proof}

外幂中能写成一个纯楔积的元素称为\term[kefenjie-zhangliang]{可分解外张量}[decomposable exterior tensor]，也简称可分解张量。这里的“可分解”专指能写成一个楔积，区别于张量积中能写成一个纯张量。上面的关系表明，在 \(\bigwedge^2k^4\) 中，可分解外张量的坐标并不是任意六个标量。例如 \(e_1\wedge e_2+e_3\wedge e_4\) 的关系左端为 \(1\)，所以在任意域上都不可分解。一般 \(r\) 维子空间的 Plücker 坐标也满足一族二次关系；这里的 \(r=2,n=4\) 是只需一个非平凡二次方程的最小例子。

\ProseParagraphBreak
另一方面，简记 \(e_{ij}=e_i\wedge e_j\)、\(e_{1234}=e_1\wedge e_2\wedge e_3\wedge e_4\)。展开 \(\omega\wedge\omega\) 时，只有指标互不重复的交叉项可能非零。两个次数二的因子交换时产生的符号为 \((-1)^{2\cdot2}=1\)，所以 \(e_{12}\wedge e_{34}\) 与 \(e_{34}\wedge e_{12}\) 同号，都等于 \(e_{1234}\)。每一对这样的交叉项因而贡献两倍，得到
\[
\omega\wedge\omega
=2(p_{12}p_{34}-p_{13}p_{24}+p_{14}p_{23})
e_1\wedge e_2\wedge e_3\wedge e_4.
\]
特征不为二时，\(\omega\wedge\omega=0\) 与 Plücker 关系等价；特征二时，每个这样的平方都是零，包括刚才的不可分解张量。因而在特征二中仍应使用坐标关系本身来判定，不能把平方为零当作充分条件。

\subsection{Koszul 复形的外代数背景}
\label{b2:subsec:0904:03}

重新回到任意交换环 \(R\)。给定整数 \(n\ge1\) 与 \(a_1,\ldots,a_n\in R\)，以 \(e_1,\ldots,e_n\) 为 \(R^n\) 的标准基。生成元之间的关系由映射
\[
\lambda:R^n\longrightarrow R,\qquad e_i\longmapsto a_i
\]
的核记录：\(\sum_i r_ie_i\) 在核中，恰好表示 \(\sum_i r_i a_i=0\)。交换性立即给出关系 \(a_ie_j-a_je_i\)，因为它的像是 \(a_i a_j-a_j a_i=0\)。二次外幂中的 \(e_i\wedge e_j\) 正好记录这一对指标，更高外幂则用来组织这些关系之间的相容性。

\ProseParagraphBreak
对任意 \(R\) 模 \(M\)，为了把线性映射 \(\lambda:M\to R\) 延伸到整个外代数，我们要求它在一次分量上等于 \(\lambda\)，在标量上为零，并遵循分次 Leibniz 规则：作用于一个楔积时，可以作用于前一因子，也可以越过前一因子后作用于后一因子；越过次数 \(p\) 的因子时带上符号 \((-1)^p\)。次数为 \(-1\) 表示把 \(\bigwedge^pM\) 映到 \(\bigwedge^{p-1}M\)。若这样的算子连续作用两次总为零，即平方为零，就称为微分。本小节将所有负次外幂约定为零，使零次上的微分也有明确的目标。

\begin{proposition}[收缩微分]\label{b2:prop:exterior-contraction}
设 \(R\) 为交换含幺环，\(M\) 为幺性 \(R\) 模，\(\lambda:M\to R\) 为 \(R\) 线性映射，并约定 \(\bigwedge^pM=0\) 对 \(p<0\) 成立。存在唯一次数为 \(-1\) 的 \(R\) 线性分次导子
\[
\iota_\lambda:\bigwedge^pM\longrightarrow\bigwedge^{p-1}M,
\]
在 \(R=\bigwedge^0M\) 上取零，在 \(M=\bigwedge^1M\) 上等于 \(\lambda\)，且对 \(\xi\in\bigwedge^pM\)、\(\eta\in\bigwedge^qM\) 满足分次 Leibniz 规则
\[
\iota_\lambda(\xi\wedge\eta)
=\iota_\lambda(\xi)\wedge\eta+(-1)^p\xi\wedge\iota_\lambda(\eta).
\]
它在纯楔积上的公式为
\begin{equation}\label{b2:eq:contraction-formula}
\iota_\lambda(m_1\wedge\cdots\wedge m_p)
=\sum_{j=1}^p(-1)^{j-1}\lambda(m_j)
m_1\wedge\cdots\wedge\widehat{m_j}\wedge\cdots\wedge m_p.
\end{equation}
这里 \(p\ge1\)，帽号表示省去该因子。这个算子称为由 \(\lambda\) 确定的\term[shousuo]{收缩}[contraction]，并满足 \(\iota_\lambda^2=0\)。
\end{proposition}
\begin{proof}
若 \(D\) 是满足陈述中条件的分次导子，先把第一个因子分出，Leibniz 规则给出
\[
D(m_1\wedge\cdots\wedge m_p)
=\lambda(m_1)m_2\wedge\cdots\wedge m_p
-m_1\wedge D(m_2\wedge\cdots\wedge m_p).
\]
从 \(p=1\) 开始归纳，便强迫得到\eqref{b2:eq:contraction-formula}。纯楔积生成各次外幂，而零次上的作用已经指定，所以至多存在一个这样的算子。

反过来，式\eqref{b2:eq:contraction-formula}的右端对输入多线性。若 \(m_r=m_s\)、\(r<s\)，除删除第 \(r\) 或第 \(s\) 个因子的项以外，其余项仍有重复因子，均为零。删去第 \(r\) 个因子后，把原第 \(s\) 个因子移回第 \(r\) 个位置，共交换 \(s-r-1\) 次，所以该项的总符号是 \((-1)^{r-1+s-r-1}=(-1)^{s-2}\)，与删去第 \(s\) 个因子的项符号相反。二者系数中的 \(\lambda(m_r),\lambda(m_s)\) 相同，故抵消。这证明右端交替，外幂泛性质给出各次上的线性算子。把零次上的作用取为零，就得到次数为 \(-1\) 的 \(\iota_\lambda\)，其在一次分量上的作用是 \(\lambda\)。

当 \(p\) 或 \(q\) 为零时，所示乘法公式由 \(R\) 线性与标量上的零作用直接得到。对两个正次纯楔积的乘积，删除位置或者在前 \(p\) 项，或者在后 \(q\) 项。前者给出第一项，后者相对从 \(\eta\) 开始计数多出 \(p\)，产生 \((-1)^p\)，得到分次 Leibniz 规则。由双线性推广到所有齐次元素。

最后，两次收缩就是选两个不同位置，各施加一次 \(\lambda\) 并删去相应因子。每对位置有两种删除顺序，它们留下相同的楔积与相同的标量，只在符号上相反。具体说，固定 \(i<j\)，先删 \(i\) 再删原第 \(j\) 个因子的符号为 \((-1)^{i+j-3}\)，先删 \(j\) 再删 \(i\) 的符号为 \((-1)^{i+j-2}\)。两项的剩余因子次序相同，标量系数分别为 \(\lambda(m_i)\lambda(m_j)\) 及其反序，由交换性相等，故相互抵消。所有项成对消去，得到 \(\iota_\lambda^2=0\)。
\end{proof}
\symidx{\(\iota_\lambda\)：外代数上的收缩算子}

\begin{corollary}\label{b2:cor:contractions-anticommute}
设 \(R\) 为交换含幺环，\(M\) 为幺性 \(R\) 模，\(\lambda,\mu:M\rightarrow R\) 为 \(R\) 线性映射。由它们确定的外代数收缩算子满足
\[
\iota_\lambda\iota_\mu+\iota_\mu\iota_\lambda=0.
\]
这个反交换关系对任意特征成立。
\end{corollary}
\begin{proof}
由\eqref{b2:eq:contraction-formula}，\(\iota_{\lambda+\mu}\) 与 \(\iota_\lambda+\iota_\mu\) 在每个纯楔积上相同，在零次部分上也都为零，所以
\(\iota_{\lambda+\mu}=\iota_\lambda+\iota_\mu\)。对 \(\lambda+\mu\) 应用\cref{b2:prop:exterior-contraction}，得到
\[
0=\iota_{\lambda+\mu}^2
=\iota_\lambda^2+\iota_\lambda\iota_\mu
+\iota_\mu\iota_\lambda+\iota_\mu^2.
\]
两个平方项已为零，剩下的就是所求关系。这里只展开算子乘积并消去已知零项，没有除以二，因此特征二时也成立。
\end{proof}

令 \(n\ge1\)、\(M=R^n\)，选基 \(e_1,\ldots,e_n\)，并指定 \(\lambda(e_i)=a_i\in R\)。取 \(K_p=\bigwedge^pR^n\)、\(\mathrm{d}_p=\iota_\lambda\)，得到
\[
0\longrightarrow K_n\xrightarrow{\mathrm{d}_n}K_{n-1}\longrightarrow\cdots
\longrightarrow K_1\xrightarrow{\mathrm{d}_1}K_0\longrightarrow0.
\]
一列 \(R\) 模与 \(R\) 线性箭头若满足每两个相邻箭头的复合为零，就称为一个\term[lianfuxing]{链复形}[chain complex]；此处的条件 \(\mathrm{d}_{p-1}\mathrm{d}_p=0\) 已由\cref{b2:prop:exterior-contraction}证明。上面由 \(a_1,\ldots,a_n\in R\) 构造的链复形称为序列 \(a_1,\ldots,a_n\) 的\term[Koszul-fuxing]{Koszul 复形}[Koszul complex]，记为 \(K_\bullet(a_1,\ldots,a_n;R)\)。\symidx{\(K_\bullet(a_1,\ldots,a_n;R)\)：Koszul 复形}它在低次上的公式是
\[
\mathrm{d}_1(e_i)=a_i,\qquad \mathrm{d}_2(e_i\wedge e_j)=a_ie_j-a_je_i.
\]
\(K_1=R^n\) 记录这组理想生成元，\(K_2\) 记录由交换性自动产生的两两关系。因此 \(\operatorname{im}\mathrm{d}_1=(a_1,\ldots,a_n)\)，它的余核就是 \(R/(a_1,\ldots,a_n)\)；把这个商接到末端，是由 \(\mathrm{d}_1\) 的像决定的。得到增广复形
\[
0\longrightarrow K_n\longrightarrow\cdots\longrightarrow K_1
\xrightarrow{\mathrm{d}_1}R\longrightarrow R/(a_1,\ldots,a_n)\longrightarrow0.
\]
例如，当 \(n=2\) 时，未增广的复形为
\[
0\longrightarrow R\xrightarrow{\binom{-a_2}{a_1}}R^2
\xrightarrow{\begin{psmallmatrix}a_1&a_2\end{psmallmatrix}}R\longrightarrow0.
\]
相邻矩阵的乘积为 \(-a_1a_2+a_2a_1=0\)，\(\mathrm{d}_2\) 记录了最直接的两两关系。

\ProseParagraphBreak
当 \(n=3\)，写 \(a_1=a,a_2=b,a_3=c\)，并简记 \(e_{ij}=e_i\wedge e_j\)、\(e_{123}=e_1\wedge e_2\wedge e_3\)。三条两两关系为 \(ae_2-be_1\)、\(ae_3-ce_1\)、\(be_3-ce_2\)，而三次外幂给出
\[
\mathrm{d}_3(e_{123})=a e_{23}-b e_{13}+c e_{12}.
\]
它说明这三条关系本身也有相容关系，因为
\[
\begin{aligned}
\mathrm{d}_2\mathrm{d}_3(e_{123})
&=a(be_3-ce_2)-b(ae_3-ce_1)+c(ae_2-be_1)\\
&=0.
\end{aligned}
\]
于是三次外幂记录的是二次外幂所给关系之间的关系，更高次外幂继续以相同方式组织这些自然关系的相容性。

\ProseParagraphBreak
这些自然关系未必生成全部关系。取任意域 \(k\)，令 \(R=k[x]\)、\(a_1=a_2=x\)。此时 \(\mathrm{d}_1(u,v)=x(u+v)\)，\(\mathrm{d}_2(t)=(-xt,xt)\)。因为 \(R\) 是整环且 \(x\ne0\)，有
\[
\ker\mathrm{d}_1=R(-1,1),\qquad
\operatorname{im}\mathrm{d}_2=xR(-1,1)\subsetneq R(-1,1).
\]
严格包含来自 \(1\notin xR\)：关系 \((-1,1)\) 在核中，却不是 \(\mathrm{d}_2\) 的像。因此序列 \(x,x\) 的 Koszul 复形在 \(R^2\) 处不正合。

\ProseParagraphBreak
复合为零给出 \(\operatorname{im}\mathrm{d}_p\subseteq\ker\mathrm{d}_{p-1}\)；在相应位置正合则要求二者相等。增广复形在 \(R\) 及商模处正合，但在其他位置未必正合。一条增广链复形若在每个位置正合，且末端模之前的各项均为自由模，就称为这个末端模的\term[ziyou-xiaojie]{自由消解}[free resolution]。这里各 \(K_p=\bigwedge^pR^n\) 都是自由模，因此是否给出自由消解，关键在于尚待检验的正合性。

\ProseParagraphBreak
最短的例子是 \(n=1\)，复形就是 \(0\to R\xrightarrow{\times a_1}R\to0\)。若 \(k\) 为域，\(R=k[\varepsilon]/(\varepsilon^2)\)、\(a_1=\varepsilon\)，则左侧映射的核为非零理想 \((\varepsilon)\)，所以增广后也不是自由消解。哪些序列使增广后的 Koszul 复形成为 \(R/(a_1,\ldots,a_n)\) 的自由消解，将在第三卷第15.2—15.3节结合正则序列系统讨论。若一个生成元已经是单位，这里就能直接构造每个核元素的原像。

\begin{proposition}\label{b2:prop:koszul-unit-contractible}
设 \(R\) 为交换含幺环，\(n\ge1\)，\(a_1,\ldots,a_n\in R\)，且 \(a_1\) 是单位。令 \(e_1,\ldots,e_n\) 为 \(R^n\) 的标准基，\(K_\bullet=K_\bullet(a_1,\ldots,a_n;R)\) 为 Koszul 复形，其微分由 \(R\) 线性映射 \(\lambda:R^n\rightarrow R\)、\(\lambda(e_i)=a_i\) 的收缩给出。把 \(K_p=\bigwedge^pR^n\) 在 \(p<0\) 或 \(p>n\) 时取为零。定义次数为 \(+1\) 的 \(R\) 线性映射
\[
h_p:K_p\longrightarrow K_{p+1},
\qquad h_p(\xi)=a_1^{-1}e_1\wedge\xi.
\]
则在每个 \(0\le p\le n\) 上有
\[
\mathrm{d}_{p+1}h_p+h_{p-1}\mathrm{d}_p=\operatorname{id}_{K_p},
\]
其中边界处的映射取零。因此 \(K_\bullet\) 正合。
\end{proposition}
\begin{proof}
由\cref{b2:prop:exterior-contraction}的分次 Leibniz 规则及 \(\mathrm{d}_1(e_1)=a_1\)，对 \(\xi\in K_p\) 有
\[
\begin{aligned}
\mathrm{d}_{p+1}h_p(\xi)
&=a_1^{-1}\cdot\mathrm{d}_{p+1}(e_1\wedge\xi)\\
&=a_1^{-1}\bigl(a_1\xi-e_1\wedge\mathrm{d}_p(\xi)\bigr)\\
&=\xi-h_{p-1}\mathrm{d}_p(\xi).
\end{aligned}
\]
当 \(p=0\) 时，\(\mathrm{d}_0=0\)；当 \(p=n\) 时，\(e_1\wedge\xi\in\bigwedge^{n+1}R^n=0\)，上述等式在两个端点仍成立。这就证明所述恒等式，通常简写为 \(\mathrm{d}h+h\mathrm{d}=\operatorname{id}\)。

若 \(\mathrm{d}_p(\xi)=0\)，恒等式给出
\(\xi=\mathrm{d}_{p+1}(h_p(\xi))\)。这为每个闭元素（微分为零的元素）明确给出了成为边界（下一次微分的像）的原像。故 \(\ker\mathrm{d}_p\subseteq\operatorname{im}\mathrm{d}_{p+1}\)，反向包含由微分平方为零成立，复形在每处正合。
\end{proof}


\begin{chaptersummary}
交替性要求重复输入时取值为零，比单纯的交换变号更能适应一般系数环。外幂通过商掉重复因子的关系表达这一性质；有限自由模的递增指标楔积构成基，证明只用交替求和和系数检测，因而不要求阶乘可逆。

\ProseParagraphBreak
楔积把所有外幂连成分次代数，分次交换律记录因子换位的次数。直和的外代数分解须在张量乘法中保留这一符号。最高外幂只有一个基向量，使线性自同态表现为乘以行列式；较低次外幂则同时记录各阶子式，函子复合给出 Cauchy–Binet 公式。

\ProseParagraphBreak
伴随矩阵恒等式不仅给出可逆判据，也通过系数比较证明一般交换环上的 Cayley–Hamilton 定理。在域上，中间次数的非零纯楔积记录子空间；四维空间中的一条 Plücker 二次关系判定二向量是否具有这种形式，特征二时不能改用平方为零来判定。整个外代数上的收缩则把生成元的两两关系及其相容性组成 Koszul 复形。收缩平方为零保证相邻箭头复合为零，单位条件给出一个可直接证明正合的特例，一般序列的正合性仍需另行检查。
\end{chaptersummary}

\BookExercises

\begin{exercise}\label{b2:ex:09-alternating-characteristic-two}
在 \(\mathbb F_2^2\) 上分别写出全部反对称双线性形式和全部交替双线性形式的矩阵。两类形式各有多少个？
\end{exercise}
\begin{solution}
特征二使反对称条件成为 \(A^{\mathsf t}=A\)，所以矩阵为 \(\begin{pmatrix}a&b\\b&c\end{pmatrix}\)，三个参数独立，有 \(2^3=8\) 个。交替还要求对角线为零，所以矩阵为 \(\begin{pmatrix}0&b\\b&0\end{pmatrix}\)，只有两个。充分性可直接展开 \(x^{\mathsf t}Ax\)：交替情形下两个交叉项之和为零，对角项已经为零。\(A=I_2\) 是反对称而不交替的例子，因为在 \(e_1\) 上的自配对为一。
\end{solution}

\begin{exercise}\label{b2:ex:09-skew-not-exterior}
将 \(M=\mathbb Z/6\mathbb Z\)、\(L=\mathbb Z/4\mathbb Z\) 视为 \(\mathbb Z\) 模。求出全部双线性映射 \(M\times M\to L\)，判断哪些反对称、哪些交替，并说明哪些能够通过 \(\bigwedge^2M\) 分解。
\end{exercise}
\begin{solution}
双线性映射由 \(\ell=b(\bar1,\bar1)\in L\) 决定，公式为 \(b(\bar a,\bar c)=ac\ell\)。代表元改变六的倍数时，该式良定义恰要求 \(6\ell=0\)。在 \(\mathbb Z/4\mathbb Z\) 中，这等价于 \(\ell=\bar0\) 或 \(\bar2\)，所以共有两个映射。二者都满足 \(2\ell=0\)，从而 \(b(\bar c,\bar a)=-b(\bar a,\bar c)\)，均反对称；交替条件却要求 \(b(\bar1,\bar1)=0\)，故只有零映射交替。因为 \(M\) 由一个元素生成，\(\bigwedge^2M=0\)，也只有零映射能经过它分解。非零映射再次显示：目标中的二挠可以使交换变号成立，而重复输入仍不为零。
\end{solution}

\begin{exercise}\label{b2:ex:09-alternating-top-form}
设 \(R\) 为交换环，\(L\) 为 \(R\) 模，\(b:(R^3)^3\to L\) 为交替三线性映射，且 \(b(e_1,e_2,e_3)=\ell\)。证明 \(b\) 由 \(\ell\) 唯一决定，并计算 \(b(e_1+e_2,e_2+e_3,e_3+e_1)\)。
\end{exercise}
\begin{solution}
由\cref{b2:prop:exterior-alternating-universal}，\(b\) 对应线性映射 \(\bigwedge^3R^3\to L\)。该外幂以 \(e_1\wedge e_2\wedge e_3\) 为单个基向量，因此由其像 \(\ell\) 唯一决定。题给三向量楔积展开后，只有 \(e_1\wedge e_2\wedge e_3\) 与 \(e_2\wedge e_3\wedge e_1\) 两项没有重复因子；后三循环符号为正，所以结果为 \(2\ell\)。这个答案在有二挠时也成立，不能预先把它当作非零。
\end{solution}

\begin{exercise}\label{b2:ex:09-cyclic-exterior}
设 \(I\) 为交换环 \(R\) 的理想，\(M=R/I\)。计算 \(\bigwedge^pM\) 的全部次数，并证明
\[
\bigwedge M\cong R[x]/(x^2,Ix)
\]
为分次 \(R\) 代数同构，其中 \(x\) 次数为一。
\end{exercise}
\begin{solution}
\(M\) 由 \(\bar1\) 生成，任意至少二重的纯楔积都是 \(\bar1\wedge\bar1\wedge\cdots\) 的标量倍数，所以为零。因此零次分量为 \(R\)，一次分量为 \(R/I\)，更高分量为零。右侧商代数的元素唯一写成常数 \(r\) 加一次项 \(\bar a x\)，乘法由 \(x^2=0\) 与 \(Ix=0\) 决定。将 \(\bar1\in M\) 送到 \(x\) 给出双向保持这些关系的映射，或逐次数比较上述唯一表达，即得同构。
\end{solution}

\begin{exercise}\label{b2:ex:09-wedge-signs}
在 \(\bigwedge\mathbb Z^4\) 中计算
\[
(e_1+e_3)\wedge(e_2-e_4),\qquad
(e_1\wedge e_3)\wedge(e_2\wedge e_4),\qquad
(e_1\wedge e_2+e_3\wedge e_4)^2.
\]
说明为什么“正次数元素都平方为零”是错误的。
\end{exercise}
\begin{solution}
第一项为 \(e_1\wedge e_2-e_1\wedge e_4-e_2\wedge e_3-e_3\wedge e_4\)。第二项把 \(e_3\) 与 \(e_2\) 交换一次，得到 \(-e_1\wedge e_2\wedge e_3\wedge e_4\)。第三项的两个同侧平方为零，两个交叉项因次数均为二而同号，故结果为 \(2e_1\wedge e_2\wedge e_3\wedge e_4\)，在整数系数下非零。\cref{b2:prop:exterior-graded-commutativity}保证奇数次齐次元素平方为零，并不对所有偶数次元素作同样断言。
\end{solution}

\begin{exercise}\label{b2:ex:09-wedge-independence}
设 \(A:k^n\to k^m\) 为域上的线性映射，\(p\ge1\)。证明 \(\bigwedge^pA\ne0\) 当且仅当 \(\operatorname{rank}A\ge p\)，并用 \(p\) 阶子式表述这个条件。对于
\[
A=\begin{pmatrix}1&0&1\\0&1&1\\1&1&2\end{pmatrix},
\]
在任意域上求最大的 \(p\)，使 \(\bigwedge^pA\ne0\)。
\end{exercise}
\begin{solution}
若 \(\operatorname{rank}A<p\)，任意 \(p\) 个像向量都相关，\cref{b2:prop:wedge-linear-independence}说明每个纯楔积都被 \(\bigwedge^pA\) 送到零；纯楔积生成整个外幂，所以该映射为零。若秩至少为 \(p\)，可选 \(p\) 个向量，使它们的像独立，其楔积非零，故 \(\bigwedge^pA\ne0\)。当 \(p\le\min(m,n)\) 时，由\eqref{b2:eq:exterior-minor-action}，该条件等价于至少一个 \(p\) 阶子式非零；超过这两个维数之一时，外幂映射为零，也没有这样的子式。

题给矩阵的第三列是前两列之和，而左上二阶子式为 \(1\)。因此它在任意域上秩均为二，最大的 \(p\) 是二。这里既用一个非零子式给出秩的下界，也用列关系给出上界，不能只由某个二阶子式非零就断言秩恰为二。
\end{solution}

\begin{exercise}\label{b2:ex:09-exterior-square-map}
将整数矩阵 \(A=\begin{pmatrix}1&1&0\\0&2&1\\0&0&3\end{pmatrix}\) 视为 \(\mathbb Z^3\) 上的线性映射。在基 \(e_1\wedge e_2,e_1\wedge e_3,e_2\wedge e_3\) 下计算 \(\bigwedge^2 A\) 的矩阵，并比较其行列式与 \((\det A)^2\)。
\end{exercise}
\begin{solution}
三列的像分别为
\(2e_1\wedge e_2\)、\(e_1\wedge e_2+3e_1\wedge e_3\)、\(e_1\wedge e_2+3e_1\wedge e_3+6e_2\wedge e_3\)。所以矩阵为
\[
\begin{pmatrix}2&1&1\\0&3&3\\0&0&6\end{pmatrix}.
\]
其行列式为 \(36\)，而 \(\det A=6\)，在本例中确有 \(\det(\bigwedge^2A)=(\det A)^2\)。这些数值是整数恒等式，向任意交换环映射后仍然成立；这里只核验本例，不以它代替一般公式的证明。
\end{solution}

\begin{exercise}\label{b2:ex:09-odd-square}
设 \(R\) 为交换环，\(M\) 为具有 \(n\) 个基向量的自由 \(R\) 模，\(N=\bigoplus_{p\ge1}\bigwedge^pM\) 为外代数中的正次理想。证明 \(N^{n+1}=0\)，并证明元素 \(a_0+\eta\)、\(a_0\in R,\eta\in N\)，可逆当且仅当 \(a_0\) 是 \(R\) 的单位。写出逆元的有限求和公式。
\end{exercise}
\begin{solution}
每个正次齐次分量的次数至少为一，所以 \(n+1\) 个正次元素相乘的所有分量都在次数 \(n+1\) 以上，由\cref{b2:thm:exterior-power-basis}全部为零。因此 \(N^{n+1}=0\)。投影到零次分量是保持单位元的代数同态 \(\bigwedge M\to R\)，故可逆元素的常数项也必须可逆。

反过来，若 \(a_0\) 可逆，令 \(z=a_0^{-1}\eta\)。此时 \(z^{n+1}=0\)，有限几何级数给出
\[
(a_0+\eta)^{-1}=a_0^{-1}\sum_{j=0}^n(-z)^j.
\]
左乘或右乘 \(a_0+\eta=a_0(1+z)\)，相邻项均抵消，只剩 \(1+(-1)^n z^{n+1}=1\)。这里每项都是同一个元素 \(z\) 的幂，且标量 \(a_0\) 在中心，所以即使外代数并非普通交换代数，仍得到双侧逆。
\end{solution}

\begin{exercise}\label{b2:ex:09-direct-sum-degree-three}
设 \(R\) 为交换环，\(M,N\) 为 \(R\) 模，分别有基 \(e_1,e_2\) 与 \(f_1,f_2\)。用直和外幂分解写出 \(\bigwedge^3(M\oplus N)\) 的一组基，并说明 \((0,f_1)\wedge(e_1,0)\wedge(0,f_2)\) 对应哪个张量及其符号。
\end{exercise}
\begin{solution}
因为两侧的三次外幂均为零，只剩
\[
\left(\bigwedge^2M\otimes_RN\right)
\oplus\left(M\otimes_R\bigwedge^2N\right).
\]
一组基为 \((e_1\wedge e_2)\otimes f_1\)、\((e_1\wedge e_2)\otimes f_2\)、\(e_1\otimes(f_1\wedge f_2)\)、\(e_2\otimes(f_1\wedge f_2)\)。题给楔积要把 \(e_1\) 移到最前，交换一次，故对应 \(-e_1\otimes(f_1\wedge f_2)\)。若系数环特征二，负号等于正号，但符号产生的规则不变。
\end{solution}

\begin{exercise}\label{b2:ex:09-exterior-not-injective}
设 \(k\) 为域，\(R=k[x,y]\)、\(I=(x,y)\)。进一步计算正文反例中的 \(\bigwedge^2I\)，证明
\[
R/I\longrightarrow\bigwedge^2I,\qquad \bar r\longmapsto r(x\wedge y)
\]
是 \(R\) 模同构，并求包含 \(I\hookrightarrow R\) 所诱导的二次外幂映射的核。
\end{exercise}
\begin{solution}
因为 \(I\) 由 \(x,y\) 生成，双线性展开说明 \(\bigwedge^2I\) 由 \(w=x\wedge y\) 生成。又
\[
xw=x\wedge(xy)=x\wedge(yx)=y(x\wedge x)=0,
\qquad yw=(xy)\wedge y=x(y\wedge y)=0.
\]
所以 \(Iw=0\)，题给映射良定义且满射。\cref{b2:prop:exterior-injection-failure}中由一次项构造的交替映射诱导 \(\widetilde\beta:\bigwedge^2I\to k=R/I\)，并满足 \(\widetilde\beta(w)=1\)。它与题给映射的复合是 \(R/I\) 的恒等，故题给映射也单射。由于 \(\bigwedge^2R=0\)，包含所诱导的映射是零映射，其核为整个 \(\bigwedge^2I\cong R/I\)。这个例子不仅检测到一个非零核元，还算出了全部核。
\end{solution}

\begin{exercise}\label{b2:ex:09-determinant-basis-change}
设 \(A=(a_{ij})\) 为交换环上的 \(n\times n\) 上三角矩阵，\(1\le p\le n\)。把递增指标组 \(I=(i_1,\ldots,i_p)\) 按字典序排列。证明 \(\bigwedge^pA\) 在相应外幂基下仍为上三角矩阵，求出其对角元，并证明
\[
\det\left(\bigwedge^pA\right)=(\det A)^{\binom{n-1}{p-1}}.
\]
\end{exercise}
\begin{solution}
\(Ae_j\) 只涉及 \(e_1,\ldots,e_j\)。若 \(e_I\) 在 \(Ae_{j_1}\wedge\cdots\wedge Ae_{j_p}\) 中出现非零系数，则必有 \(i_s\le j_s\) 对所有 \(s\) 成立：否则前 \(s\) 个输入所选的不同基指标都须不超过 \(j_s\)，但目标指标组中这样的指标至多有 \(s-1\) 个，不可能。于是非零系数只能出现于字典序不大于 \(J\) 的行 \(I\)，证明矩阵上三角。

对应 \(I=J\) 的对角元是子矩阵 \(A_{J,J}\) 的行列式，而它仍上三角，故为 \(\prod_{j\in J}a_{jj}\)。于是
\[
\det\left(\bigwedge^pA\right)
=\prod_{|J|=p}\prod_{j\in J}a_{jj}
=\prod_{j=1}^n a_{jj}^{\binom{n-1}{p-1}}
=(\det A)^{\binom{n-1}{p-1}}.
\]
中间的指数来自固定 \(j\) 后，从其余 \(n-1\) 个指标中选取 \(p-1\) 个的方式数；没有使用可逆性或约去任何对角元。
\end{solution}

\begin{exercise}\label{b2:ex:09-gram-cauchy-binet}
设 \(A\) 为实 \(2\times n\) 矩阵。证明
\[
\det(AA^{\mathsf t})=\sum_{1\le i<j\le n}\det\bigl(A_{\{1,2\},\{i,j\}}\bigr)^2.
\]
由此证明该行列式非负，并说明它何时为正。
\end{exercise}
\begin{solution}
将\cref{b2:thm:cauchy-binet}用于 \(A\) 与 \(A^{\mathsf t}\)。每个对应的行子式是原列子式的转置，行列式相同，因此每项是平方。和非负，且为正当且仅当至少一个二阶子式非零。若有两个独立列，对应二阶子式非零；若没有，全部列至多生成一维空间。故为正恰好等价于 \(A\) 的秩为二。\(n<2\) 时为空和，行列式为零，结论仍成立。
\end{solution}

\begin{exercise}\label{b2:ex:09-block-triangular}
对交换环上的分块矩阵 \(C=\begin{pmatrix}A&B\\0&D\end{pmatrix}\)，其中 \(A,D\) 分别为 \(m,n\) 阶方阵，证明 \(\det C=\det A\det D\)，不要求两个块可逆。
\end{exercise}
\begin{solution}
将 \(R^{m+n}\) 的基分为前 \(m\) 个 \(e_i\) 和后 \(n\) 个 \(f_j\)。前 \(m\) 列的楔积为 \((\det A)e_1\wedge\cdots\wedge e_m\)。在其后乘上后 \(n\) 列的楔积时，只要从任何一列选到前 \(m\) 个坐标的分量，就与前面的全部 \(e_i\) 重复而为零。因此只有后块 \(D\) 的分量贡献，其系数为 \(\det D\)，得到结论。这个推理只用展开和重复因子为零，不需要消去 \(\det A\) 或取逆。
\end{solution}

\begin{exercise}\label{b2:ex:09-adjugate-mod-eight}
在 \(R=\mathbb Z/8\mathbb Z\) 上判断 \(A=\begin{pmatrix}1&2\\3&1\end{pmatrix}\) 是否可逆；若可逆，用伴随矩阵求逆。再判断 \(B=\begin{pmatrix}2&0\\0&1\end{pmatrix}\) 是否可逆。
\end{exercise}
\begin{solution}
\(\det A=1-6=\bar3\)，它的逆为 \(\bar3\)，因此
\[
A^{-1}=3\begin{pmatrix}1&-2\\-3&1\end{pmatrix}
=\begin{pmatrix}3&2\\7&3\end{pmatrix}\pmod8.
\]
直接相乘也得到单位矩阵。\(\det B=\bar2\) 非零但不是单位，所以 \(B\) 不可逆；事实上 \(B(\bar4,\bar0)=(\bar0,\bar0)\)，连单射都不满足。
\end{solution}

\begin{exercise}\label{b2:ex:09-cayley-recurrence}
令 \(A=\begin{pmatrix}1&1\\1&0\end{pmatrix}\)，\(F_0=0,F_1=1,F_{n+1}=F_n+F_{n-1}\)。用 Cayley–Hamilton 定理证明 \(A^n=F_nA+F_{n-1}I_2\) 对所有 \(n\ge1\) 成立，并求 \(A^5\)。
\end{exercise}
\begin{solution}
\(\chi_A(t)=t^2-t-1\)，故 \(A^2=A+I_2\)。\(n=1\) 的公式成立；若 \(n\) 时成立，则
\(A^{n+1}=F_n(A+I_2)+F_{n-1}A=F_{n+1}A+F_nI_2\)，归纳完成。\(F_5=5,F_4=3\)，所以 \(A^5=5A+3I_2=\begin{pmatrix}8&5\\5&3\end{pmatrix}\)。这是整数矩阵恒等式，在任意交换环中解释相应整数系数后仍成立。
\end{solution}

\begin{exercise}\label{b2:ex:09-characteristic-coefficients}
设 \(A\in M_n(R)\)、\(n\ge1\)，其中 \(R\) 交换。证明 \(\chi_A(t)\) 的 \(t^{n-1}\) 系数为 \(-\operatorname{tr}A\)，常数项为 \((-1)^n\det A\)。说明该论证为什么不需要特征值。
\end{exercise}
\begin{solution}
常数项为 \(\det(-A)=(-1)^n\det A\)，由各列提出负号得到。为得到 \(t^{n-1}\)，Leibniz 展开中须选取 \(n-1\) 个对角线上的 \(t\)；相应置换已固定 \(n-1\) 个指标，因而也固定最后一个，所以只能来自 \(\prod_i(t-a_{ii})\)。从第 \(i\) 项选 \(-a_{ii}\)、其余选 \(t\)，求和得到 \(-\sum_i a_{ii}\)。全程是交换环上的多项式系数计算，不要求特征多项式分裂或存在特征值。
\end{solution}

\begin{exercise}\label{b2:ex:09-plucker-explicit}
在任意域上，判断坐标按 \(12,13,14,23,24,34\) 排列为 \((1,2,0,3,1,2)\) 的二张量是否可分解；若是，给出两个因子。这里整数按基域的特征解释。
\end{exercise}
\begin{solution}
关系左端为 \(1\cdot2-2\cdot1+0\cdot3=0\)，所以\cref{b2:prop:plucker-four-dimensional}适用。可取
\[
u=e_1-3e_3-e_4,\qquad v=e_2+2e_3.
\]
展开后五个前项坐标分别为 \(1,2,0,3,1\)，\(34\) 坐标为 \((-3)\cdot0-(-1)\cdot2=2\)。因为 \(12\) 坐标始终为一，楔积在任何特征下都非零，两个因子均线性无关。
\end{solution}

\begin{exercise}\label{b2:ex:09-plucker-plane-recovery}
设 \(k\) 为域，\(e_1,\ldots,e_4\) 为 \(k^4\) 的标准基。令 \(\omega=(e_1+e_3)\wedge(e_2+e_4)\in\bigwedge^2k^4\)。求所有满足 \(v\wedge\omega=0\) 的向量 \(v\)，并说明它与 \(\omega\) 的两条生成向量之间的关系。
\end{exercise}
\begin{solution}
写 \(v=\sum_i a_ie_i\)。因为 \(e_1+e_3,e_2+e_4\) 独立，\eqref{b2:eq:plucker-recovers-plane}说明解空间正是它们的线性生成空间，所以
\[
v=a(e_1+e_3)+b(e_2+e_4),\qquad a,b\in k.
\]
等价条件为 \(a_3=a_1\)、\(a_4=a_2\)。也可把 \(\omega=e_1\wedge e_2+e_1\wedge e_4-e_2\wedge e_3+e_3\wedge e_4\) 代入，再比较四个递增三重楔积的系数，得到同样的两个独立方程。
\end{solution}

\begin{exercise}\label{b2:ex:09-contractions-anticommute}
设 \(R\) 为交换环，\(M\) 为 \(R\) 模，\(m\in M\)、\(\lambda:M\to R\) 线性。令 \(L_m(\xi)=m\wedge\xi\)。证明
\[
\iota_\lambda L_m+L_m\iota_\lambda=\lambda(m)\operatorname{id}_{\bigwedge M}.
\]
当 \(M=R^2\)、\(m=e_1\)、\(\lambda(e_1)=1\)、\(\lambda(e_2)=0\) 时，在基 \(1,e_1,e_2,e_1\wedge e_2\) 下写出两个算子的矩阵并核对等式。
\end{exercise}
\begin{solution}
收缩的分次乘法公式给出
\[
\iota_\lambda(m\wedge\xi)=\lambda(m)\xi-m\wedge\iota_\lambda(\xi),
\]
移项即得等式，无须除以二。在指定基下，两个矩阵分别为
\[
[L_{e_1}]=\begin{pmatrix}0&0&0&0\\1&0&0&0\\0&0&0&0\\0&0&1&0\end{pmatrix},\qquad
[\iota_\lambda]=\begin{pmatrix}0&1&0&0\\0&0&0&0\\0&0&0&1\\0&0&0&0\end{pmatrix}.
\]
前者将 \(1,e_2\) 分别送到 \(e_1,e_1\wedge e_2\)，后者将 \(e_1,e_1\wedge e_2\) 分别送到 \(1,e_2\)，其余基向量都送到零。因此两种复合的和在四个基向量上均为恒等。这条混合关系与\cref{b2:cor:contractions-anticommute}中的两个收缩反交换关系不同：它的右端可以非零。
\end{solution}

\begin{exercise}\label{b2:ex:09-koszul-unit}
设 \(a_1,\ldots,a_n\) 生成交换环 \(R\) 的单位理想，取 \(b_i\in R\) 使 \(\sum_i b_i a_i=1\)。在 Koszul 复形中令 \(u=\sum_i b_ie_i\)、\(h(\xi)=u\wedge\xi\)，证明 \(\mathrm{d}h+h\mathrm{d}=\operatorname{id}\)。对 \(R=\mathbb Z\)、序列 \(2,3\)，显式写出各次 \(h\) 的矩阵，说明正合性不要求序列中一定含有单位。
\end{exercise}
\begin{solution}
因为 \(\mathrm{d}(u)=\sum_i b_i a_i=1\)，收缩的分次乘法公式给出 \(\mathrm{d}(u\wedge\xi)=\xi-u\wedge\mathrm{d}\xi\)，故所求等式成立。若 \(\mathrm{d}\xi=0\)，便有 \(\xi=\mathrm{d}(h\xi)\)；配合 \(\mathrm{d}^2=0\)，得到每个位置正合。

在整数例中，取 \(u=-e_1+e_2\)，因为 \(-2+3=1\)。在 \(K_1\) 的基 \(e_1,e_2\) 和 \(K_2\) 的基 \(e_1\wedge e_2\) 下，
\[
\mathrm{d}_1=\begin{pmatrix}2&3\end{pmatrix},\quad
\mathrm{d}_2=\binom{-3}{2},\quad
h_0=\binom{-1}{1},\quad h_1=\begin{pmatrix}-1&-1\end{pmatrix},\quad h_2=0.
\]
直接相乘得 \(\mathrm{d}_1h_0=1\)、\(\mathrm{d}_2h_1+h_0\mathrm{d}_1=I_2\)、\(h_1\mathrm{d}_2=1\)。所以该复形正合，尽管 \(2,3\) 在 \(\mathbb Z\) 中都不是单位。\cref{b2:prop:koszul-unit-contractible}的单位元条件是充分条件；本题用 Bézout 组合将它推广到生成单位理想的序列。
\end{solution}

\begin{exercise}\label{b2:ex:09-koszul-three}
设 \(k\) 为任意域，令 \(R=k[t]\)，考虑序列 \(t,t,t\) 的 Koszul 复形。在 \(K_1\) 的基 \(e_1,e_2,e_3\) 与 \(K_2\) 的基 \(e_1\wedge e_2,e_1\wedge e_3,e_2\wedge e_3\) 下写出全部微分矩阵，核对相邻乘积为零，并计算
\[
H_p\coloneqq\ker\mathrm{d}_p/\operatorname{im}\mathrm{d}_{p+1}\qquad(0\le p\le3),
\]
其中 \(\mathrm{d}_0=0\)、\(K_4=0\)。这些商模称为复形的同调模；它们如何判断各位置的正合性？
\end{exercise}
\begin{solution}
微分矩阵为
\[
\mathrm{d}_1=t\begin{pmatrix}1&1&1\end{pmatrix},\qquad
\mathrm{d}_2=t\begin{pmatrix}-1&-1&0\\1&0&-1\\0&1&1\end{pmatrix},\qquad
\mathrm{d}_3=t\begin{pmatrix}1\\-1\\1\end{pmatrix}.
\]
两个相邻乘积都为零。由于 \(R\) 是整环，求核时可消去非零的 \(t\)。令 \(u=(-1,1,0)\)、\(v=(-1,0,1)\)，则
\[
\ker\mathrm{d}_1=Ru\oplus Rv,\qquad
\operatorname{im}\mathrm{d}_2=tRu\oplus tRv.
\]
后一等式来自中间矩阵的三列 \(u,v,v-u\)。同样，\(\ker\mathrm{d}_2=R(1,-1,1)\)，而 \(\operatorname{im}\mathrm{d}_3=tR(1,-1,1)\)；\(\mathrm{d}_3\) 单射。因此
\[
H_0\cong R/(t),\qquad H_1\cong(R/(t))^2,\qquad
H_2\cong R/(t),\qquad H_3=0.
\]
在第 \(p\) 项正合恰好等价于 \(H_p=0\)。这个复形在 \(K_1,K_2\) 处不正合，尽管两个相邻乘积都为零；接上末端 \(R/(t)\) 只能修复零次余核，不能消去这两个非零同调模。
\end{solution}

% !TeX root = ../../Book2.tex
% 纲要标题：## 第10章　双线性型、二次型与经典群
% 纲要位置：计划纲要.md，第 1372 行。

\chapter{双线性型、二次型与经典群}
\label{b2:ch:10}
\BookChapterNumber{b2:ch:10}

\begin{chapterabstract}
一个型不仅给向量赋值，还规定哪些方向相互正交。双线性型的换基产生合同，而线性算子的换基产生相似；本章先分清这两种分类问题，再建立交替型标准形、特征不为二的二次型对角化及实惯性定律。一般域上的二次型则通过双曲平面的分离、反射与 Witt 消去整理为双曲部分和各向异性部分。最后讨论保持各种型的经典群，并以 Clifford 代数的商构造说明二次型与结合代数的联系。
\end{chapterabstract}

\begin{learninggoals}
\begin{itemize}
\item 辨认非退化配对、根空间与正交补，说明何时正交补给出直和分解。
\item 对具体矩阵作合同化简，比较不同域及不同特征下的二次型。
\item 掌握双曲平面的分离和反射，能够逐步说明 Witt 分解的存在性、唯一性与消去。
\item 从保持型的条件写出正交群、辛群与酉群，区分双线性与 Hermitian 约定，并理解 Clifford 代数的泛性质。
\end{itemize}
\end{learninggoals}

在实平面上，\(q_+(x,y)=x^2+y^2\) 与 \(q_-(x,y)=x^2-y^2\) 的矩阵都可逆，但两者的几何行为不同：前者只在原点为零，后者在直线 \(x=y\) 和 \(x=-y\) 上也为零。若把底域改为复数，坐标代换 \(y'=iy\) 就有 \(x^2-y^2=x^2+(y')^2\)，两型因而等距。一个型的分类因此不仅要问矩阵是否可逆，还要问在什么域上、允许怎样换基，以及哪些向量与子空间被它判为正交或零值。

\ProseParagraphBreak

对称型可以通过对角化逐步拆解；交替型和一般双线性型却需要不同的操作。特别在特征二，\(q(v+w)-q(v)-q(w)\) 仍给出极化信息，却不能靠除以 \(2\) 从中恢复 \(q\)。把这些适用条件说清，再看保持型的变换，正交群及由二次关系构造的代数才有准确的起点。

\ProseParagraphBreak

本章需要\chapref{b2:ch:01}的有限维向量空间、对偶、矩阵换基及行列式的基本性质。\secref{b2:sec:1001}对任意域讨论双线性型；\secref{b2:sec:1002}明确区分特征二；\secref{b2:sec:1003}始终假设特征不为二，并在有限维非退化空间中证明 Witt 理论。\subsecref{b2:subsec:1004:03}还使用\cref{b2:thm:tensor-algebra-universal}的张量代数泛性质。

\ProseParagraphBreak
双线性型及其几何可参阅 Roman 的《Advanced Linear Algebra》\cite[Chapter 11]{Roman2008AdvancedLinearAlgebra}；实惯性与 Witt 消去分别见同书\cite[Theorems 11.25--11.26, 11.34]{Roman2008AdvancedLinearAlgebra}。比较文献时应先检查极化型是否除以二。本章在特征不为二时统一写 \(q(v)=b(v,v)\)，因而双曲平面的表达式为 \(2xy\)。

% !TeX root = ../../Book2.tex
% 纲要标题：### 10.1 双线性型
% 纲要位置：计划纲要.md，第 1374 行。

\section{双线性型}
\label{b2:sec:1001}

线性映射接受一个向量，双线性型则同时比较两个向量。选定基以后，两者都能写成矩阵，但换基时的规则不同。本节先把配对两端的空间分开；只有在说明对称性或交替性以后，才把左右正交条件合并。

% 纲要要点（供目录核对）：
%
% * 配对、有限维非退化性及正交补；用对偶限制映射判定维数与直和
% * 对称型与交替型；交替型的二维分离在任意特征成立
% * 型的合同与算子的相似
%

\subsection{配对、非退化性与正交补}
\label{b2:subsec:1001:01}

\begin{definition}\label{b2:def:bilinear-pairing}
设 \(k\) 为域，\(U,V\) 为 \(k\) 向量空间。如果映射 \(b:U\times V\to k\) 在固定任意一个变量后对另一个变量都是 \(k\) 线性的，就称 \(b\) 为\term[shuangxianxing-peidui]{双线性配对}[bilinear pairing]。它的左根空间与右根空间分别为
\[
\{\,u\in U\mid b(u,v)=0\;\text{对所有}\;v\in V\,\},\qquad
\{\,v\in V\mid b(u,v)=0\;\text{对所有}\;u\in U\,\}.
\]
两者均为零时，称配对\term[peidui-feituihua]{非退化}[nondegenerate]。当 \(U=V\) 时，称 \(b\) 为 \(V\) 上的\term[shuangxianxing-xing]{双线性型}[bilinear form]。
\end{definition}

以下空间均有限维。固定 \(u\in U\) 后，它与第二侧所有向量的配对值组成 \(V\) 上的线性泛函 \(b(u,-)\)；固定 \(v\in V\) 则得到 \(U\) 上的线性泛函 \(b(-,v)\)。因此配对把一侧的向量变成另一侧的线性泛函，给出两个映射
\[
L_b:U\longrightarrow V^*,\quad u\longmapsto b(u,-),
\qquad
R_b:V\longrightarrow U^*,\quad v\longmapsto b(-,v).
\]
它们的核就是上述两个根空间。对这样的有限维 \(k\) 向量空间上的配对 \(b:U\times V\to k\)，若 \(L_b:U\to V^*\) 是同构，就称 \(b\) 为\term[wanmei-peidui]{完美配对}[perfect pairing]；这表示 \(V\) 上的每个线性泛函都唯一地由某个 \(u\in U\) 与 \(V\) 配对得到。

\begin{proposition}\label{b2:prop:finite-perfect-pairing}
设 \(k\) 为域，\(U,V\) 为有限维 \(k\) 向量空间，\(b:U\times V\to k\) 为双线性配对。令 \(L_b:U\to V^*\)、\(u\mapsto b(u,-)\) 及 \(R_b:V\to U^*\)、\(v\mapsto b(-,v)\)，则 \(b\) 非退化当且仅当 \(L_b\) 与 \(R_b\) 都是同构；这也等价于其中任一个是同构。此时 \(\dim U=\dim V\)。
\end{proposition}
\begin{proof}
非退化使两个映射都单射，故 \(\dim U\le\dim V\le\dim U\)，它们因而同为同构。若只知 \(L_b\) 是同构，且 \(R_b(v)=0\)，则全部 \(U\) 中元素都与 \(v\) 配为零。\(L_b\) 满射说明每个 \(V\) 上的线性泛函都在 \(v\) 处为零，由\cref{b2:lem:functional-extension}的分离结论，\(v=0\)。又两空间等维，所以 \(R_b\) 也是同构。交换两个变量证明另一方向。
\end{proof}

例如，求值 \(V^*\times V\to k\) 在有限维时完美。若取 \(b:k^2\times k^3\to k\)、\(b(u,v)=u_1v_1+u_2v_2\)，左根空间为零，右根空间却是一维；只检验一端不够。

\ProseParagraphBreak
上面从两个单射推出两个同构，使用了有限维的维数比较。若取无限维空间 \(V=k^{(\mathbb N)}\) 及配对 \(b(u,v)=\sum_i u_iv_i\)，每个向量支集有限使求和有限，且用标准基向量 \(e_i\) 检测便知两端根空间都为零。可是线性泛函 \(\ell(v)=\sum_i v_i\) 满足 \(\ell(e_i)=1\) 对每个 \(i\) 都成立，不可能等于任何 \(b(u,-)\)，因为 \(u\) 的支集有限；因此这个非退化配对的 \(L_b\) 并不满射。

\ProseParagraphBreak
对于同一空间上的一般型，条件 \(b(u,v)=0\) 与 \(b(v,u)=0\) 未必等价。以下两种型允许我们合并这些条件。

\begin{definition}\label{b2:def:symmetric-alternating-form}
设 \(k\) 为域，\(V\) 为 \(k\) 向量空间，\(b:V\times V\to k\) 为双线性型。如果 \(b(u,v)=b(v,u)\) 对所有 \(u,v\in V\) 成立，就称 \(b\) 为\term[duichen-xing]{对称型}[symmetric bilinear form]；如果 \(b(v,v)=0\) 对所有 \(v\in V\) 成立，就称它为\term[jiaoti-xing]{交替型}[alternating bilinear form]。如果 \(b(u,v)=-b(v,u)\) 对所有 \(u,v\in V\) 成立，就称它为\term[fanduichen-xing]{反对称型}[skew-symmetric bilinear form]。
\end{definition}

将 \(b(u+v,u+v)=0\) 展开，便知交替型总是反对称的。对于子空间 \(W\)，先分别记
\[
W^{\perp_r}=\{\,v\in V\mid b(w,v)=0\ (w\in W)\,\},\qquad
{}^{\perp_l}W=\{\,v\in V\mid b(v,w)=0\ (w\in W)\,\}.
\]
若 \(b\) 对称，或满足 \(b(v,u)=-b(u,v)\)，两者相同，统记为 \(W^\perp\)，称为\term[zhengjiao-bu]{正交补}[orthogonal complement]；此时 \(\operatorname{rad}b=V^\perp\) 称为型的\term[xing-de-gen]{根空间}[radical of a form]。这里的“补”首先是正交条件的名称，并不预先断言 \(V=W\oplus W^\perp\)。
\symidx{\(W^\perp\)：对称型或交替型下的正交补}
\symidx{\(\operatorname{rad}b\)：双线性型的根空间}

\begin{theorem}\label{b2:thm:orthogonal-complement}
设 \(k\) 为域，\(V\) 为有限维 \(k\) 向量空间，\(b:V\times V\to k\) 为对称型或交替型，\(W\le V\)。记 \(W^\perp=\{v\in V\mid b(v,w)=0\;\text{对所有}\;w\in W\}\)。若 \(b\) 在 \(V\) 上非退化，则
\[
\dim W^\perp=\dim V-\dim W,\qquad (W^\perp)^\perp=W.
\]
无论 \(b\) 在 \(V\) 上是否退化，下述正交直和分解成立当且仅当 \(b|_{W\times W}\) 非退化：
\[
V=W\mathbin{\perp}W^\perp,
\]
其中符号 \(\perp\) 表示两个分量互相正交的直和。在上述分解成立的前提下，若 \(b\) 在 \(V\) 上也非退化，则 \(W^\perp\) 上的限制型非退化；若 \(b\) 在 \(V\) 上退化，则不保证这一点。
\end{theorem}
\begin{proof}
正交补是满足全部配对方程的向量组成的核。具体地，考虑线性映射
\[
\Phi:V\longrightarrow W^*,\qquad
\Phi(v)(w)=b(w,v)\quad(w\in W).
\]
由型的对称性或反对称性，\(\ker\Phi=W^\perp\)。当 \(V\) 上非退化时，\(\Phi\) 满射：给定 \(W\) 上的泛函，先由\cref{b2:lem:functional-extension}延拓到 \(V\)，再由\cref{b2:prop:finite-perfect-pairing}写成 \(b(-,v)\)。于是秩与核的维数公式给出第一个等式。正交条件的对称性给出 \(W\subseteq(W^\perp)^\perp\)，再由维数相等得到相等。

限制型的根空间满足
\[
\operatorname{rad}(b|_{W\times W})=W\cap W^\perp,
\]
因为两边都由 \(W\) 中与全部 \(W\) 配为零的向量组成。因此正交直和分解若成立，限制型就非退化。反过来，现在只假设 \(W\) 上非退化。对任意 \(v\in V\)，\cref{b2:prop:finite-perfect-pairing}在 \(W\) 上给出唯一 \(w\in W\)，使 \(b(x,w)=b(x,v)\) 对所有 \(x\in W\) 成立。因此 \(v-w\in W^\perp\)，而上述根空间等式给出 \(W\cap W^\perp=0\)，这就证明直和分解。若 \(V\) 上也非退化，而 \(z\in W^\perp\) 与全部 \(W^\perp\) 正交，则它还与 \(W\) 正交，故与 \(V\) 正交，只能为零。
\end{proof}

若 \(b\) 在 \(k^2\) 上的矩阵为 \(\operatorname{diag}(1,0)\)，取 \(W=ke_1\)，则 \(W\) 上非退化，且 \(W^\perp=ke_2\)；后者上的限制型却为零。

\ProseParagraphBreak
例如，实平面上 \(b(x,y)=x_1y_1-x_2y_2\) 非退化，但直线 \(W=\mathbb R(1,1)\) 满足 \(W^\perp=W\)。维数公式正确，直和分解却不成立，因为限制型为零。非退化性不能从整个空间自动传给每个子空间。

\subsection{对称型与交替型}
\label{b2:subsec:1001:02}

交替性已经蕴含反对称性。特征不为二时，反对称性反过来给出 \(2b(v,v)=0\)，所以二者等价。特征二时不能这样约去二；对称型 \(b(x,y)=xy\) 就不是交替型。交替型的定义在任意特征下都要求对角取值为零。

\begin{theorem}[交替型的标准形]\label{b2:thm:alternating-form-normal}
设 \(k\) 为任意特征的域，\(V\) 为有限维 \(k\) 向量空间，\(b:V\times V\to k\) 为交替双线性型，则 \(V\) 有一组基
\[
e_1,f_1,\ldots,e_m,f_m,z_1,\ldots,z_r,
\]
使不同的两维分量及最后的根空间互相正交，并有
\[
b(e_i,f_i)=1,\quad b(f_i,e_i)=-1,\quad b(e_i,e_i)=b(f_i,f_i)=0.
\]
\(z_1,\ldots,z_r\) 是根空间的基。因此型的秩为 \(2m\)，并完全决定其在给定维数下的合同类型。特别地，非退化交替型只存在于偶数维空间。
\end{theorem}
\begin{proof}
若型为零，任取基即可。否则要分离一个非退化的两维平面，再在其正交补中继续同样的构造；交替型的每条直线上的限制型都为零，所以这里不能只分离一条直线。选 \(e,f_0\) 使 \(b(e,f_0)\ne0\)，将 \(f_0\) 除以该值，得到 \(b(e,f)=1\)。\(e,f\) 线性无关：若 \(ae+cf=0\)，分别与 \(e,f\) 配对便得 \(c=a=0\)。它们的线性生成空间 \(H\) 上的矩阵为
\[
\begin{pmatrix}0&1\\-1&0\end{pmatrix},
\]
行列式为一，故限制型非退化。由\cref{b2:thm:orthogonal-complement}中不要求整体非退化的分解，\(V=H\perp H^\perp\)。在较低维的 \(H^\perp\) 上归纳，直到剩余限制型为零。每个已分离的两维块非退化，故最后的零型分量恰是整个型的根空间。

上述基给出的矩阵由 \(m\) 个同样的可逆两维块及一个 \(r\) 维零块组成，所以秩为 \(2m\)。秩不随换基改变，\(r=\dim V-2m\) 也被确定；同秩的两个型把上述基逐一对应便得到保持配对的同构。
\end{proof}

非退化交替型的一组 \(e_i,f_i\) 称为\term[xin-ji]{辛基}[symplectic basis]。若先排列 \(e_1,\ldots,e_m\)，再排列 \(f_1,\ldots,f_m\)，其矩阵就是
\[
J_m=\begin{pmatrix}0&I_m\\-I_m&0\end{pmatrix}.
\]
这些论证没有除以二，因而在特征二中仍然有效；此时 \(-I_m=I_m\)，单个辛块变成
\[
\begin{pmatrix}0&1\\1&0\end{pmatrix}.
\]
它既对称又交替：对称性来自 \(-1=1\)，交替性仍由对角线为零保证。

\subsection{型的合同与算子的相似}
\label{b2:subsec:1001:03}

选基 \(\mathcal E=(e_1,\ldots,e_n)\)，令 \(B=\bigl(b(e_i,e_j)\bigr)\)。若 \(x,y\) 是两个向量的坐标列，则
\[
b(u,v)=x^{\mathsf T}By.
\]
这既说明矩阵由型唯一确定，也说明每个矩阵都确定一个双线性型。非退化恰好等于 \(B\) 可逆；对称性等于 \(B^{\mathsf T}=B\)，交替性等于 \(B^{\mathsf T}=-B\) 且对角元全零。

\begin{proposition}\label{b2:prop:bilinear-congruence}
设 \(k\) 为域，\(V\) 为 \(n\) 维 \(k\) 向量空间，\(b:V\times V\to k\) 为双线性型，其旧基矩阵为 \(B\)。若新基的列在旧基中的坐标组成 \(P\in\operatorname{GL}_n(k)\)，则型的新矩阵为
\[
B'=P^{\mathsf T}BP.
\]
存在这样的 \(P\) 时，称两矩阵\term[hetong]{合同}[congruent]。两个型之间的\term[dengju-tonggou]{等距同构}[isometry]是保持双线性取值的线性双射；它们等距，当且仅当相应矩阵合同。
\end{proposition}
\begin{proof}
新坐标为 \(x',y'\) 时，旧坐标为 \(Px',Py'\)，故取值是 \((x')^{\mathsf T}P^{\mathsf T}BPy'\)。逐个代入标准基向量便得到矩阵公式。一般线性双射也由可逆矩阵给出，保持两变量取值的等式正是同一个合同等式。
\end{proof}

算子换基的公式则为 \(A'=P^{-1}AP\)：输入的新坐标先乘 \(P\) 变为旧坐标，作用 \(A\) 后，输出的旧坐标再乘 \(P^{-1}\) 变回新坐标。双线性型的两个位置都是输入变量，两侧的坐标都乘 \(P\)，所以给出 \(P^{\mathsf T}BP\)。实一维型的矩阵 \((1)\) 与 \((4)\) 通过 \(P=(2)\) 合同，但一维矩阵相似只能使其不变，故二者不相似。

\ProseParagraphBreak
合同保持秩。对于非退化型，等式
\[
\det B'=(\det P)^2\det B
\]
还说明行列式在 \(k^\times/(k^\times)^2\) 中的平方类不变。不过秩与行列式平方类合在一起仍不足以分类。例如实数域上 \(I_2\) 与 \(-I_2\) 都有秩二和行列式一，但在非零向量上，一个二次取值恒正，另一个恒负，所以不可能合同。下一节将从二次取值出发说明这里还缺少什么信息。

% !TeX root = ../../Book2.tex
% 纲要标题：### 10.2 二次型
% 纲要位置：计划纲要.md，第 1380 行。

\section{二次型}
\label{b2:sec:1002}

双线性型在两个位置取同一个向量，便得到二次表达式。这个过程能否逆转，取决于系数域的特征。本节始终令 \(V\) 有限维，并先在任意域上定义二次型；讨论对角化和惯性时，再明确加入特征与次序条件。

% 纲要要点（供目录核对）：
%
% * 特征不为2时的对角化及一维型的平方类判据
% * 实数域上的惯性定律；正负指数由定号子空间的最大维数确定
% * 特征2时二次型与极化型的区别；相同非退化极化型不足以判定等距
%

\subsection{特征不为二时的对角化}
\label{b2:subsec:1002:01}

\begin{definition}\label{b2:def:quadratic-form}
设 \(k\) 为域，\(V\) 为 \(k\) 向量空间。如果映射 \(q:V\to k\) 满足 \(q(av)=a^2\cdot q(v)\) 对所有 \(a\in k,v\in V\) 成立，且映射
\begin{equation}\label{b2:eq:quadratic-polar}
\beta_q(u,v)=q(u+v)-q(u)-q(v)
\end{equation}
是 \(k\) 双线性的，就称 \(q\) 为 \(V\) 上的\term[erci-xing]{二次型}[quadratic form]，称 \(\beta_q\) 为它的\term[jihua-xing]{极化型}[polar form]。带有指定二次型 \(q\) 的向量空间 \((V,q)\) 称为\term[erci-kongjian]{二次空间}[quadratic space]；这一名称允许 \(q\) 退化。
\end{definition}

第一个条件说明沿同一方向伸缩时，取值按标量的平方变化；极化型的双线性则控制不同方向之间的混合项。若 \(e_1,\ldots,e_n\) 是一组基，反复展开 \(q(u+v)=q(u)+q(v)+\beta_q(u,v)\)，再用缩放条件，就得到任意特征下都成立的坐标公式
\[
q\left(\sum_i x_ie_i\right)
=\sum_i x_i^2\cdot q(e_i)
+\sum_{i<j}x_ix_j\cdot\beta_q(e_i,e_j).
\]
因此这两个公理恰好使坐标表达式由平方项与两个不同坐标的乘积组成。

\ProseParagraphBreak
极化型总对称。当 \(\operatorname{char}k\ne2\) 时，本章采用归一化记号
\begin{equation}\label{b2:eq:normalized-polar}
b_q(u,v)=\frac12\beta_q(u,v),\qquad q(v)=b_q(v,v).
\end{equation}
第二个等式来自 \(\beta_q(v,v)=q(2v)-2q(v)=2q(v)\)：除以二正是为了使两个变量取同一向量时恢复 \(q\) 的值。反过来，对任意对称双线性型 \(b\)，令 \(q(v)=b(v,v)\)，则 \(\beta_q=2b\)。所以在特征不为二时，两类对象一一对应。以下在这一特征假设下，未加下标的 \(b\) 表示 \(b_q\)，未归一化的极化型仍记为 \(\beta_q\)。
\symidx{\(\beta_q\)、\(b_q=\beta_q/2\)：极化型及特征不为二时的归一化型}

在这项特征假设下，称 \(q\) 非退化，是指 \(b_q\) 非退化；两个二次型等距，是指线性双射保持其二次取值。由极化公式，这等价于保持 \(b_q\)。

\begin{theorem}[二次型的对角化]\label{b2:thm:quadratic-diagonalization}
设 \(k\) 为特征不为二的域，\(V\) 为有限维 \(k\) 向量空间，\(q:V\to k\) 为二次型，并令 \(b_q=\beta_q/2\)，则 \(V\) 有一组基，使
\[
q(x_1e_1+\cdots+x_ne_n)=a_1x_1^2+\cdots+a_rx_r^2,
\qquad a_i\ne0.
\]
最后 \(n-r\) 个基向量组成 \(\operatorname{rad}b_q\) 的基。特别地，每个对称矩阵都合同于对角矩阵。记 \(\langle a\rangle\) 为一维二次型 \(t\mapsto at^2\)。对任意 \(a,b\in k^\times\)，在上述 \(\operatorname{char}k\ne2\) 假设下，有
\[
\langle a\rangle\cong\langle b\rangle
\quad\Longleftrightarrow\quad
a/b\in(k^\times)^2,
\]
其中同构指保持二次取值的线性同构。
\end{theorem}
\begin{proof}
若 \(q\) 为零，极化公式使 \(b_q=0\)，任意基都合适。否则取 \(e\) 使 \(a=q(e)\ne0\)。直线 \(ke\) 上的限制型非退化，\cref{b2:thm:orthogonal-complement}给出
\[
V=ke\perp(ke)^\perp.
\]
为消去 \(e\) 方向与其余方向之间的混合项，要从 \(v\) 中减去一个 \(e\) 的倍数，使剩余向量与 \(e\) 正交。由 \(b(e,e)=a\)，这一倍数只能是 \(b(e,v)/a\)，因而任意 \(v\) 分解为
\[
v=\frac{b(e,v)}a e+\left(v-\frac{b(e,v)}a e\right),
\]
第二项与 \(e\) 正交。在较小维的正交补上归纳，得到全部非零一维分量及最后的零型分量。不同分量之间的混合项为零，所以二次表达式具有所述形式；从 \(b_q\) 的对角矩阵可见，最后这些基向量恰好生成 \(\operatorname{rad}b_q\)。

一维线性双射必为 \(t\mapsto ct\)，其中 \(c\in k^\times\)。
从 \(\langle a\rangle\) 到 \(\langle b\rangle\) 的这个映射保持二次取值，当且仅当 \(a=bc^2\)。
故两型等距当且仅当 \(a/b\) 是非零平方。
\end{proof}

常将非退化对角型写为 \(\langle a_1,\ldots,a_r\rangle\)。对角化并不表示这列系数已经唯一：由一维型的平方类判据，\(\langle1\rangle\) 与 \(\langle4\rangle\) 在 \(\mathbb Q\) 上等距，但与 \(\langle2\rangle\) 不等距。
\symidx{\(\langle a_1,\ldots,a_r\rangle\)：对角二次型}

\ProseParagraphBreak
例如，在 \(\mathbb Q\) 上取 \(q(x,y)=x^2+2xy+3y^2\)。令 \(u=x+y,v=y\)，便得 \(q=u^2+2v^2\)。新基在旧坐标中是 \((1,0),(-1,1)\)，两者对于矩阵 \(\begin{psmallmatrix}1&1\\1&3\end{psmallmatrix}\) 正交；新矩阵为 \(\operatorname{diag}(1,2)\)。这是同时变换两个变量的合同计算。

\begin{corollary}\label{b2:cor:quadratic-algebraically-closed}
设 \(k\) 为特征不为二的代数闭域，\(V\) 为有限维 \(k\) 向量空间，则 \(V\) 上的二次型在等距同构意义下由其秩完全分类。
\end{corollary}
\begin{proof}
由\cref{b2:thm:quadratic-diagonalization}对角化。每个非零 \(a_i\) 都有平方根，缩放相应基向量后将它改为一，得到 \(\operatorname{diag}(I_r,0)\)。反过来合同保持秩，故不同的 \(r\) 不等距。
\end{proof}

在代数闭域上，所有非零系数属于同一个平方类，所以对角化之后只剩秩；在 \(\mathbb Q\) 上，非零平方类已经有很多种，例如一与二不能彼此缩放得到。实数域上的非零系数则可缩为正一或负一，但这两个平方类不能合并。

\subsection{实数域上的惯性定律}
\label{b2:subsec:1002:02}

在 \(k=\mathbb R\) 上，非零平方不能改变符号。对角化后可以把每个非零系数的绝对值缩放为一，正负号却必须保留。对于实向量空间 \(V\) 上的二次型 \(q\)，若每个非零 \(v\in V\) 都满足 \(q(v)>0\)，就称 \(q\)\term[zhengding]{正定}[positive definite]；若 \(-q\) 正定，就称 \(q\) 负定。

\begin{theorem}[Sylvester 惯性定律]\label{b2:thm:sylvester-inertia}
设 \(V\) 为有限维实向量空间，\(q:V\to\mathbb R\) 为二次型，则 \(V\) 有一组基，使 \(q\) 在相应坐标下有标准表达式
\[
q(x)=x_1^2+\cdots+x_{p_+}^2-x_{p_++1}^2-\cdots-x_{p_++p_-}^2,
\qquad p_++p_-+r=\dim V.
\]
三元组 \((p_+,p_-,r)\) 唯一，且
\[
\begin{aligned}
p_+&=\max\{\,\dim W\mid W\le V,\ q|_W\text{ 正定}\,\},\\
p_-&=\max\{\,\dim W\mid W\le V,\ -q|_W\text{ 正定}\,\}.
\end{aligned}
\]
两个实二次型等距，当且仅当这三个数相同。
\end{theorem}
\begin{proof}
存在性由\cref{b2:thm:quadratic-diagonalization}及基向量乘 \(1/\sqrt{|a_i|}\) 得到。为证唯一性，记正项、负项及零项的坐标子空间为 \(V_+,V_-,V_0\)。\(V_+\) 上的型正定，维数为 \(p_+\)。若 \(W\) 是任意一个正定子空间，投影 \(W\to V_+\) 必单射：投影为零的 \(w\) 属于 \(V_-\oplus V_0\)，因而 \(q(w)\le0\)，正定性使 \(w=0\)。所以 \(\dim W\le p_+\)。这说明 \(p_+\) 是正定子空间的最大维数，与所选标准基无关。

对 \(-q\) 作同样论证，负项数 \(p_-\) 也由型确定。\(r=\dim V-p_+-p_-\) 因而确定；它也等于根空间的维数。相同三元组的型把各类标准基逐一对应，便得到等距同构，证明分类断言。
\end{proof}

这一结果称为\term[Sylvester-guanxing-dinglv]{Sylvester 惯性定律}[Sylvester's law of inertia]。\footnote{西尔维斯特（James Joseph Sylvester，1814-1897），英国数学家，研究矩阵、行列式与不变量理论。}正项数、负项数分别称\term[zheng-guanxing-zhishu]{正惯性指数}[positive index of inertia]、\term[fu-guanxing-zhishu]{负惯性指数}[negative index of inertia]，差 \(p_+-p_-\) 称为\term[fuhaocha]{符号差}[signature]。正定比非退化强。\(q(v)=0\) 只要求一个二次取值为零，而 \(v\in\operatorname{rad}b_q\) 要求 \(b_q(v,w)=0\) 对所有 \(w\) 成立；在特征不为二时，后者取 \(w=v\) 就推出前者，反向却不成立。例如实型 \(q(x,y)=x^2-y^2\) 非退化，\(v=(1,1)\) 满足 \(q(v)=0\)，但 \(b_q(v,(1,0))=1\)，所以不在根空间中。

\ProseParagraphBreak
例如 \(q(x,y,z)=2xy+z^2\)。令 \(u=(x+y)/\sqrt2,v=(x-y)/\sqrt2\)，得到 \(q=u^2-v^2+z^2\)，惯性为 \((2,1,0)\)。

\subsection{特征二的二次型与极化型}
\label{b2:subsec:1002:03}

现在假设 \(\operatorname{char}k=2\)。仍用\eqref{b2:eq:quadratic-polar}定义 \(\beta_q\)，但不能定义 \(\beta_q/2\)。又因 \(q(0)=0\)，有 \(\beta_q(v,v)=0\)，所以极化型必为交替型。

\begin{proposition}\label{b2:prop:quadratic-char-two-coordinates}
设 \(k\) 为特征二的域，\(V\) 为以 \(e_1,\ldots,e_n\) 为基的有限维 \(k\) 向量空间，\(q:V\to k\) 为二次型，\(\beta_q\) 为其极化型。在这组基的二次表达式中，混合项 \(x_ix_j\) 的系数为 \(\beta_q(e_i,e_j)\)（\(i<j\)）；各平方项 \(x_i^2\) 的系数 \(q(e_i)\) 可以任意改变而不改变极化型。因此，当 \(n\ge1\) 时，单凭 \(\beta_q\) 不能恢复 \(q\)。
\end{proposition}
\begin{proof}
坐标公式已由二次型的两个公理展开得到，混合项系数正是 \(\beta_q(e_i,e_j)\)。在特征二中，每个平方项的极化为零，因为 \((x+y)^2=x^2+y^2\)；改变它的系数仍得到二次型，且不改变极化型。若 \(n\ge1\)，改变一个平方项系数就改变相应 \(e_i\) 处的取值，所以极化型不能确定 \(q\)。
\end{proof}

例如一维上的零型与 \(q(x)=x^2\) 具有相同的零极化型，但不是同一个二次型。在 \(\mathbb F_2^2\) 上，\(q(x,y)=xy\) 的极化型为
\[
\beta_q\bigl((x,y),(u,v)\bigr)=xv+yu,
\]
其矩阵 \(\begin{psmallmatrix}0&1\\1&0\end{psmallmatrix}\) 可逆。特征二中，线性组合的平方仍是各平方项的线性组合，例如 \((ax+by)^2=a^2x^2+b^2y^2\)，所以换基也不能让平方项产生所需的混合项。任何仅含平方项的表达式 \(\sum a_ix_i^2\) 在任意基下都有零极化型，与这个 \(q\) 的非退化极化型矛盾。因此对称双线性型的矩阵分类不能替代特征二二次型的分类。

\ProseParagraphBreak
特征二时，映射 \(q\mapsto\beta_q\) 已经丢失平方项的信息；而特征不为二时，\(q(v)=\dfrac12\beta_q(v,v)\) 能恢复全部取值。“极化型非退化”“二次型有非零零点”与“根空间非零”需要分别说明，尤其不能由极化型为零就断言二次型为零。即使极化型非退化，也可能有两个不等距的二次型与它对应。

\begin{proposition}\label{b2:prop:quadratic-char-two-polar-failure}
设 \(V=\mathbb F_2^2\)，定义二次型
\[
q_0(x,y)=xy,\qquad q_1(x,y)=x^2+xy+y^2.
\]
两者有相同的非退化极化型
\[
\beta_{q_0}\bigl((x,y),(u,v)\bigr)
=\beta_{q_1}\bigl((x,y),(u,v)\bigr)=xv+yu,
\]
但二次空间 \((V,q_0)\) 与 \((V,q_1)\) 不等距。
\end{proposition}
\begin{proof}
两个表达式均按标量的平方缩放；直接展开极化公式，平方项在特征二中消去，两者都得到 \(xv+yu\)，所以它们确为二次型。这个极化型的矩阵为 \(\begin{psmallmatrix}0&1\\1&0\end{psmallmatrix}\)，可逆，因而非退化。

空间恰有三个非零向量 \((1,0),(0,1),(1,1)\)。逐个计算得
\[
\begin{aligned}
q_0(1,0)&=0,&q_0(0,1)&=0,&q_0(1,1)&=1,\\
q_1(1,0)&=1,&q_1(0,1)&=1,&q_1(1,1)&=1.
\end{aligned}
\]
因此 \(q_0\) 有两个非零零点，\(q_1\) 没有非零零点。线性双射会置换这三个非零向量，而等距同构还必须保持二次取值，所以它不可能把 \(q_0\) 变成 \(q_1\)。
\end{proof}

下一节讨论的 Witt 分解重新假设特征不为二。

% !TeX root = ../../Book2.tex
% 纲要标题：### 10.3 Witt 理论初步
% 纲要位置：计划纲要.md，第 1386 行。

\section{Witt 理论初步}
\label{b2:sec:1003}

从本节起固定 \(\operatorname{char}k\ne2\)，所讨论的二次空间均有限维，且以 \(b=\beta_q/2\) 为相应对称型。实数域上的正负号分类不能直接搬到一般域上，但非零零点仍能帮助我们分离出一种固定的两维型。Witt 理论先取出这些两维部分，再研究剩下没有非零零点的部分。

% 纲要要点（供目录核对）：
%
% * 双曲平面、各向同性与各向异性；各向同性不等于属于根空间
% * Witt 分解与消去定理；反射用于对齐共同非退化分量
% * Witt 群的概念；忽略双曲部分后由相反型给出加法逆元
%

\subsection{双曲平面、各向同性与各向异性}
\label{b2:subsec:1003:01}

\begin{definition}\label{b2:def:isotropic-hyperbolic}
设 \(k\) 为特征不为二的域，\((V,q)\) 为 \(k\) 上的二次空间。若 \(v\in V\) 满足 \(v\ne0\) 且 \(q(v)=0\)，就称 \(v\) 为\term[gexiang-tongxing-xiangliang]{各向同性向量}[isotropic vector]。若不存在这样的向量，就称 \((V,q)\)\term[gexiang-yixing]{各向异性}[anisotropic]；零空间按此定义也是各向异性的。若子空间 \(T\subseteq V\) 满足 \(q|_T=0\)，就称 \(T\) 为\term[quan-gexiang-tongxing]{全各向同性子空间}[totally isotropic subspace]。
\end{definition}

各向同性只要求一个非零向量的自身取值 \(q(v)=b(v,v)\) 为零；属于根空间则要求 \(b(v,w)=0\) 对所有 \(w\in V\) 成立。因此根空间中的非零向量必各向同性，反过来却不成立。全各向同性要求子空间中每个向量的二次取值为零，由极化公式，这等价于 \(b|_{T\times T}=0\)，比仅指定若干基向量各自取值为零更强。非退化空间的根空间为零，但仍可以有非零全各向同性子空间，例如上一节的 \(x^2-y^2\)。

\begin{definition}\label{b2:def:hyperbolic-plane}
设 \(k\) 为特征不为二的域，\((V,q)\) 为二维 \(k\) 二次空间，\(b=\beta_q/2\)。如果 \(V\) 有基 \(e,f\)，满足 \(q(e)=q(f)=0\)、\(b(e,f)=1\)，就称 \((V,q)\) 为\term[shuangqu-pingmian]{双曲平面}[hyperbolic plane]，记为 \(H\)。它的表达式及矩阵为
\[
q(xe+yf)=2xy,\qquad B_H=\begin{pmatrix}0&1\\1&0\end{pmatrix}.
\]
若一个 \(k\) 二次空间等距于若干个 \(H\) 的正交直和，就称其为\term[shuangqu-kongjian]{双曲空间}[hyperbolic space]，并记 \(H^{\perp r}\) 为 \(r\) 个双曲平面的正交直和。
\end{definition}
\symidx{\(H\)、\(H^{\perp r}\)：双曲平面及其正交直和}

表达式中的因子二来自本书的约定 \(q(v)=b(v,v)\)：两个混合项分别为 \(xyb(e,f)\) 与 \(xyb(f,e)\)，各等于 \(xy\)。若把未归一化的极化型 \(\beta_q\) 在基向量间的值取为一，才会得到 \(xy\) 的表达式。

\ProseParagraphBreak
所有这样给出的平面都等距，因为把指定的 \(e,f\) 相互对应便保持矩阵。\(\det B_H=-1\ne0\)，所以它非退化，却包含全各向同性直线 \(ke\) 与 \(kf\)。这两条直线的和不是全各向同性子空间，因为 \(q(e+f)=2\ne0\)；各向同性向量 \(e\) 也不属于根空间，因为 \(b(e,f)=1\)。基 \(e+f/2,e-f/2\) 相互正交，二次取值为一和负一，因而
\[
H\cong\langle1,-1\rangle.
\]
若 \(k=\mathbb R\)，这正是一个正方向与一个负方向组成的部分。

\begin{lemma}[双曲平面的分离]\label{b2:lem:hyperbolic-plane-split}
设 \(k\) 为特征不为二的域，\((V,q)\) 为有限维非退化 \(k\) 二次空间，\(b=\beta_q/2\)。若 \(e\in V\) 满足 \(e\ne0\) 且 \(q(e)=0\)，则存在 \(f\in V\)，使 \(e,f\) 生成一个双曲平面 \(H\)，并有 \(V=H\perp H^\perp\)，其中正交补非退化。
\end{lemma}
\begin{proof}
由于 \(q(e)=0\)，直线 \(ke\) 上的限制型是零型，不能把它单独取作非退化正交分量，也不能像分离非零对角项时那样除以 \(q(e)\)。我们要给它补上另一个各向同性方向，使两个方向之间的配对非零。

非退化性保证某个 \(f_0\) 满足 \(b(e,f_0)\ne0\)。缩放以后令此值为一。先尝试 \(f=f_0+te\)，其中 \(t\in k\)。由 \(q(v)=b(v,v)\)、\(q(e)=0\) 及 \(b(e,f_0)=b(f_0,e)=1\)，有
\[
b(e,f)=1,\qquad
q(f)=q(f_0)+2t\,b(f_0,e)+t^2q(e)=q(f_0)+2t.
\]
所以保持 \(b(e,f)=1\) 的同时，让 \(q(f)=0\) 只需求解线性方程 \(q(f_0)+2t=0\)。因为特征不为二，取 \(t=-q(f_0)/2\)，即
\[
f=f_0-\frac{q(f_0)}2e.
\]
两向量的配对矩阵为 \(B_H\)，所以它们线性无关且限制型非退化。现在应用\cref{b2:thm:orthogonal-complement}，即得所需正交分解与正交补的非退化性。
\end{proof}

\begin{proposition}\label{b2:prop:isotropic-subspace-hyperbolic}
设 \(k\) 为特征不为二的域，\((V,q)\) 为有限维非退化 \(k\) 二次空间。若 \(T\subseteq V\) 为 \(d\) 维全各向同性子空间，则 \(T\) 包含于某个等距于 \(H^{\perp d}\) 的非退化子空间。
\end{proposition}
\begin{proof}
对 \(d\) 归纳，\(d=0\) 时取零空间。\(d>0\) 时，取 \(0\ne e\in T\)，由\cref{b2:lem:hyperbolic-plane-split}补上 \(f\)，令 \(H=ke+kf\)。因为 \(T\) 全各向同性，其中每个 \(t\) 已与 \(e\) 正交；要让它落入 \(H^\perp\)，还须消去它与 \(f\) 的配对。利用 \(b(e,f)=1\)，沿 \(e\) 调整，置
\[
t'=t-b(t,f)e.
\]
此时 \(b(t',e)=0\)，且 \(b(t',f)=b(t,f)-b(t,f)b(e,f)=0\)。又 \(t'\in T\)，所以每个 \(t\) 都有表示 \(t=b(t,f)e+t'\)。由于 \(e\notin H^\perp\)，这个表示给出直和
\[
T=ke\oplus\bigl(T\cap H^\perp\bigr).
\]
交的维数为 \(d-1\)，仍全各向同性；在非退化的 \(H^\perp\) 内应用归纳假设，再加上 \(H\)，便得到所求子空间。
\end{proof}

\subsection{Witt 分解与消去定理}
\label{b2:subsec:1003:02}

不断分离双曲平面，每次维数减少二，便能得到分解的存在性。要说明所得各向异性部分不依赖选择，需要先证明正交直和可以消去共同的非退化部分。这里的困难是，\(U\perp V\) 到 \(U\perp W\) 的等距同构未必把指定的 \(U\) 分量送到另一边的 \(U\) 分量。将 \(U\) 对角化后，只需逐条对齐共同的非退化直线，再把映射限制到正交补；下面的反射正是用来调整这些直线的像。

\begin{lemma}[反射与等长向量]\label{b2:lem:orthogonal-reflection}
设 \(k\) 为特征不为二的域，\((V,q)\) 为有限维 \(k\) 二次空间，\(b=\beta_q/2\)。若 \(z\in V\) 满足 \(q(z)\ne0\)，则 \(V=kz\perp(kz)^\perp\)。在 \(kz\) 上取负恒等、在 \((kz)^\perp\) 上取恒等所定义的映射 \(\tau_z\) 是等距自同构，称为沿 \(z\) 的\term[zhengjiao-fanshe]{反射}[reflection]。它的公式是
\[
\tau_z(v)=v-\frac{2b(v,z)}{q(z)}z.
\]
若 \(u,v\in V\) 满足 \(q(u)=q(v)=a\ne0\)，则存在至多两个这样的反射的复合，将 \(u\) 送到 \(v\)。
\end{lemma}
\begin{proof}
因为 \(q(z)\ne0\)，直线 \(kz\) 的限制型非退化，\cref{b2:thm:orthogonal-complement}给出 \(V=kz\perp(kz)^\perp\)，这里不要求整个 \(V\) 非退化。对任意 \(v\in V\)，这一分解具体写成
\[
v=\frac{b(v,z)}{q(z)}z+v_\perp,
\qquad v_\perp\in(kz)^\perp.
\]
只把第一项取负，得到 \(-\dfrac{b(v,z)}{q(z)}z+v_\perp=v-\dfrac{2b(v,z)}{q(z)}z\)，这就导出了所给公式。两个正交分量上的作用均保持二次取值，且平方为恒等，所以 \(\tau_z\) 是等距自同构。

对等长向量，计算得
\[
q(u-v)=2\bigl(a-b(u,v)\bigr),\qquad
q(u+v)=2\bigl(a+b(u,v)\bigr).
\]
两者之和为 \(4a\ne0\)，故至少一个非零。若 \(q(u-v)\ne0\)，则
\[
\frac{2b(u,u-v)}{q(u-v)}=1,
\qquad \tau_{u-v}(u)=v.
\]
若 \(q(u-v)=0\)，则 \(q(u+v)\ne0\)，同样计算得 \(\tau_{u+v}(u)=-v\)。再施加 \(\tau_v\)，便把 \(-v\) 送到 \(v\)。这里 \(q(v)=a\ne0\)，所以第二个反射确有定义。\(u=v\) 时也可直接取恒等映射。
\end{proof}

这项反射构造可与 Roman\cite[Chapter 11, Theorem 11.32]{Roman2008AdvancedLinearAlgebra}对照。非零长度的条件十分具体：它既保证至少一个 \(u\pm v\) 能定义反射，也保证必要时能再沿 \(v\) 反射。各向同性向量不能直接使用这个论证。

\begin{theorem}[Witt 消去]\label{b2:thm:witt-cancellation}
设 \(k\) 为特征不为二的域，\(U,V,W\) 为有限维非退化 \(k\) 二次空间。若
\[
U\perp V\cong U\perp W,
\]
则 \(V\cong W\)。
\end{theorem}
\begin{proof}
先设共同部分一维。在两边分别取基向量 \(e,f\)，使 \(q(e)=q(f)=a\ne0\)，并给定等距同构
\[
\phi:ke\perp V\longrightarrow kf\perp W.
\]
目标空间中的 \(\phi(e)\) 与 \(f\) 等长且长度非零。由\cref{b2:lem:orthogonal-reflection}，有目标空间的等距自同构 \(\rho\)，使 \(\rho\phi(e)=f\)。因此 \(\rho\phi\) 将 \((ke)^\perp=V\) 送入 \((kf)^\perp=W\)：对 \(v\in V\)，保持配对给出 \(b\bigl(\rho\phi(v),f\bigr)=b(v,e)=0\)。反向应用它的逆同构，得到像恰为 \(W\)，故限制映射是 \(V\to W\) 的等距同构。

一般情形中，由\cref{b2:thm:quadratic-diagonalization}，共同非退化部分可写成
\[
U=\langle a_1\rangle\perp\cdots\perp\langle a_d\rangle,
\qquad a_i\ne0.
\]
依次对两边共同的一维部分应用已证情形。每次余下的空间仍是非退化正交直和，故下一次消去仍满足假设。经过 \(d\) 次即得 \(V\cong W\)；\(d=0\) 时本来就是该结论。
\end{proof}

上述结果称为\term[Witt-xiaoqu-dingli]{Witt 消去定理}[Witt cancellation theorem]。

\begin{theorem}[Witt 分解]\label{b2:thm:witt-decomposition}
设 \(k\) 为特征不为二的域，\((V,q)\) 为有限维非退化 \(k\) 二次空间，则 \(V\) 有分解
\[
V\cong H^{\perp r}\perp A,
\]
其中 \(A\) 各向异性。整数 \(r\) 与 \(A\) 的等距类型唯一。\(r\) 等于全各向同性子空间的最大维数，称为\term[Witt-zhishu]{Witt 指数}[Witt index]。
\end{theorem}
\begin{proof}
若有非零各向同性向量，便用\cref{b2:lem:hyperbolic-plane-split}取出一个双曲平面。在非退化正交补上重复，每次维数减少二，因此有限步后剩余部分 \(A\) 没有非零各向同性向量。若原空间已经各向异性，则取 \(r=0\)。

设还有 \(V\cong H^{\perp s}\perp B\)，其中 \(B\) 各向异性。不妨 \(r\ge s\)。由\cref{b2:thm:witt-cancellation}消去 \(s\) 个双曲平面，得到 \(H^{\perp(r-s)}\perp A\cong B\)。若 \(r>s\)，左侧含非零各向同性向量，右侧却没有，矛盾。因此 \(r=s\)，继续用所得等距即得 \(A\cong B\)。

所列双曲平面中的全部 \(e_i\) 线性无关，且 \(q(e_i)=0\)、不同平面的向量之间配对为零，所以它们生成一个 \(r\) 维全各向同性子空间。反过来，任意 \(d\) 维全各向同性子空间由\cref{b2:prop:isotropic-subspace-hyperbolic}包含在某个 \(H^{\perp d}\) 中。将其非退化正交补再作已经证明存在的 Witt 分解，得到整个 \(V\) 的一个分解，其中双曲平面的数目至少为 \(d\)。由唯一性，它恰为 \(r\)，故 \(d\le r\)。
\end{proof}

这一分解称为\term[Witt-fenjie]{Witt 分解}[Witt decomposition]。若 \(q\) 退化，记 \(R=\operatorname{rad}b\)。对 \(z\in R\)，有 \(q(z)=b(z,z)=0\) 及 \(b(v,z)=0\)，所以 \(q(v+z)=q(v)\)。因此商空间上有良定义的二次型
\[
\bar q(v+R)=q(v).
\]
它非退化：若 \(v+R\) 与全部商向量正交，就有 \(b(v,w)=0\) 对所有 \(w\in V\) 成立，从而 \(v\in R\)、\(v+R=0\)。根空间 \(R\)、商空间 \(V/R\) 及这个商型都由原型直接确定，不需要选取补空间。

\ProseParagraphBreak
要在 \(V\) 内写出直和分解，可取 \(R\) 的一个普通线性补空间 \(V_0\)。自然映射 \(V_0\to V/R\) 是线性同构，且保持二次取值，所以 \(V_0\) 的限制型非退化。对它应用 Witt 分解，得到
\[
V=R\perp V_0
\cong R\perp H^{\perp r}\perp A.
\]
第一项是零型。与第一章的商空间和补空间之别相同，\(V_0\) 在 \(V\) 中的具体位置一般不唯一，但每个这样的补空间都等距于同一个商型；商空间的 Witt 指数与各向异性部分的等距类型因而由原二次型确定。

\subsection{Witt 群的概念}
\label{b2:subsec:1003:03}

正交直和给出的等距类可以相加，但对非零 \(V\)，\((V,q)\perp(V,-q)\) 的底层空间是 \(V\oplus V\)，维数仍为 \(2\dim V\)，不能等距于零空间。一维情形已经揭示了应忽略什么：对 \(a\in k^\times\)，在 \(\langle a,-a\rangle\) 中取
\[
e=(1,1),\qquad f=\left(\frac1{2a},-\frac1{2a}\right).
\]
直接计算得 \(q(e)=q(f)=0\)、\(b(e,f)=1\)。这两个向量线性无关，因为把 \(ce+df=0\) 分别与 \(e,f\) 配对便得 \(d=c=0\)。所以 \(\langle a,-a\rangle\cong H\)。若允许增删双曲平面而不改变类，这对互为负号的一维型之和便为零；对角化后同样的配对将使 \((V,-q)\) 成为 \((V,q)\) 的加法逆元。

\begin{definition}\label{b2:def:witt-equivalence}
设 \(k\) 为特征不为二的域，\(V,W\) 为有限维非退化 \(k\) 二次空间。如果存在整数 \(r,s\ge0\)，使
\[
V\perp H^{\perp r}\cong W\perp H^{\perp s},
\]
就称 \(V\) 与 \(W\)\term[Witt-dengjia]{Witt 等价}[Witt equivalent]。所有等价类组成的集合记为 \(W(k)\)。在类上定义 \([V]+[W]=[V\perp W]\)。
\end{definition}

\begin{theorem}\label{b2:thm:witt-group}
设 \(k\) 为特征不为二的域，\(W(k)\) 为有限维非退化 \(k\) 二次空间按添加双曲平面所得的 Witt 等价类集合，则 Witt 等价是等价关系，正交直和给出 \(W(k)\) 上良定义的交换群运算。零元为零空间的类，\((V,q)\) 的负元由 \((V,-q)\) 表示。每个类恰有一个各向异性代表的等距类型。
\end{theorem}
\begin{proof}
自反性和对称性直接成立。若 \(V\perp H^{\perp r}\cong W\perp H^{\perp s}\)，又 \(W\perp H^{\perp u}\cong X\perp H^{\perp v}\)，则给第一个等距加上 \(H^{\perp u}\)，给第二个加上 \(H^{\perp s}\)，得到
\[
V\perp H^{\perp(r+u)}\cong X\perp H^{\perp(v+s)}.
\]
故关系传递。对两个等距分别取正交直和，便证明加法与代表选择无关。直和的结合、交换及零空间的单位性质给出相应的群公理。

还须构造负元。将 \(q\) 对角化为 \(\langle a_1,\ldots,a_n\rangle\)，其中 \(a_i\ne0\)。\(q\perp(-q)\) 重排分量后为若干个 \(\langle a_i,-a_i\rangle\)，由定义前的一维计算，每个分量都等距于 \(H\)。因此 \((V,q)\perp(V,-q)\cong H^{\perp n}\)，其类为零，给出负元。

\cref{b2:thm:witt-decomposition}说明每个类都由各向异性部分表示。若各向异性 \(A,B\) Witt 等价，则 \(A\perp H^{\perp r}\cong B\perp H^{\perp s}\) 的两边已是 Witt 分解，唯一性迫使 \(r=s\) 且 \(A\cong B\)。
\end{proof}

这个群称为域 \(k\) 的\term[Witt-qun]{Witt 群}[Witt group]。记号中的集合没有大小困难：每个有限维型都可以由某个有限对称矩阵表示，全部这些矩阵组成集合，再取等价类即可。
\symidx{\(W(k)\)：特征不为二的域的 Witt 群}

\begin{proposition}\label{b2:prop:witt-real-complex}
对有限维非退化实二次空间，正惯性指数与负惯性指数之差给出群同构 \(W(\mathbb R)\cong\mathbb Z\)。对有限维非退化复二次空间，维数的奇偶性给出群同构 \(W(\mathbb C)\cong\mathbb Z/2\mathbb Z\)。
\end{proposition}
\begin{proof}
实型由\cref{b2:thm:sylvester-inertia}分为 \(p_+\) 个 \(\langle1\rangle\) 与 \(p_-\) 个 \(\langle-1\rangle\)。一正一负合成双曲平面；配对后剩余系数全为一或全为负一，是各向异性型。它的维数与符号由差 \(p_+-p_-\) 决定，所以其 Witt 类只由这个差决定。符号差在正交直和下相加，加双曲平面不改变它；差为零时全部分量都能配成双曲平面，而任意整数都能由全正或全负型实现，故所得同态既单射又满射。

复数域上，先由\cref{b2:thm:quadratic-diagonalization}对角化，再利用每个非零复数都有平方根，把全部非零系数缩放为一。因此非退化型只有维数这一分类数据。\(\langle1,1\rangle\cong\langle1,-1\rangle\cong H\)，第一个等距通过一个坐标乘以 \(i\) 得到。因此偶数维型的类为零，奇数维型的类为 \(\bigl[\langle1\rangle\bigr]\)。加双曲平面不改变维数奇偶，故这个非零类不可能等于零，得到所述二阶群。
\end{proof}

张量积还能在这个群上定义乘法。\appref{b2:app:B}的\cref{b2:thm:B-witt-ring}证明该乘法对Witt类良定义，\cref{b2:prop:B-discriminant-witt}进一步讨论维数奇偶、判别式和基本理想；这些是本节加法理论之后的补充。

% !TeX root = ../../Book2.tex
% 纲要标题：### 10.4 经典群及复数结构
% 纲要位置：计划纲要.md，第 1392 行。

\section{经典群及复数结构}
\label{b2:sec:1004}

一个型给空间增加了结构，保持这种结构的可逆线性变换便组成一个群。不同的型给出不同的经典群；复数域上还须区别双线性配对与带共轭的配对。本节最后的 Clifford 代数说明二次型也能直接参与一个结合代数的定义。

% 纲要要点（供目录核对）：
%
% * 正交群、辛群及酉群
% * Hermitian 型与复数结构；用实部和乘以i恢复完整配对
% * 完整证明非退化型的伴随算子存在唯一性、反乘法性质及像与核关系；型的相似变换与乘子的实际像
% * 区别型的相似变换、矩阵共轭与一般线性群；正交、辛、Hermitian 条件分别写清
% * Clifford 代数的构造作为选读延伸；零型回到外代数，证明双曲平面及一维模型
% #### 10.4.3* Clifford 代数的构造
%
% ---
%

\subsection{正交群、辛群与酉群}
\label{b2:subsec:1004:01}

设 \(\operatorname{char}k\ne2\)，\((V,q)\) 为有限维非退化二次空间。保持 \(q\) 的线性自同构组成\term[zhengjiao-qun]{正交群}[orthogonal group]
\[
\operatorname O(V,q)=\{\,g\in\operatorname{GL}(V)\mid q(gv)=q(v)\text{ 对所有 }v\in V\,\}.
\]
由\eqref{b2:eq:normalized-polar}，保持 \(q\) 等价于保持 \(b_q\)。在 Gram 矩阵为 \(B\) 的基下，\footnote{格拉姆（Jørgen Pedersen Gram，1850-1916），丹麦数学家与精算师，研究向量内积和正交化。}这个群写为
\[
\operatorname O(B)=\{\,G\in\operatorname{GL}_n(k)\mid G^{\mathsf T}BG=B\,\}.
\]
保持非退化交替型 \(b\) 的线性自同构则组成\term[xin-qun]{辛群}[symplectic group]。在\cref{b2:thm:alternating-form-normal}给出的辛基下，记为
\[
\operatorname{Sp}_{2m}(k)=\{\,G\in\operatorname{GL}_{2m}(k)\mid G^{\mathsf T}J_mG=J_m\,\}.
\]
辛群的这一定义允许任意特征。
\symidx{\(\operatorname O(V,q)\)、\(\operatorname{Sp}_{2m}(k)\)：正交群与辛群}

这些集合确为群：恒等变换保持型；保持型的两个映射的复合也保持型；在保持型的等式中代入逆映射的像，就得到逆映射也保持型。改变基只改变群在一般线性群中的矩阵实现，不改变其抽象含义。

\ProseParagraphBreak
对 \(G^{\mathsf T}BG=B\) 取行列式，因 \(\det B\ne0\)，得到 \((\det G)^2=1\)。于是 \((\det G-1)(\det G+1)=0\)；域中没有零因子，故 \(\det G\in\{1,-1\}\)，而特征不为二保证这两个值不同。行列式为一的子群称为\term[teshu-zhengjiao-qun]{特殊正交群}[special orthogonal group]，记为 \(\operatorname{SO}(V,q)\)。当 \(V\ne0\) 时，\cref{b2:thm:quadratic-diagonalization}给出非零长度向量，沿它的反射在该直线上为负恒等、在正交补上为恒等，所以行列式为负一。故此时特殊正交群在正交群中指数为二。

\ProseParagraphBreak
复向量空间 \(\mathbb C^n\) 上的配对
\[
h_0(z,w)=\overline z^{\mathsf T}w=\sum_i\overline{z_i}w_i
\]
在第二变量线性，在第一变量共轭线性。保持它的复线性变换组成\term[you-qun]{酉群}[unitary group]
\[
\operatorname U(n)=\{\,G\in\operatorname{GL}_n(\mathbb C)\mid G^*G=I_n\,\},
\qquad G^*=\overline G^{\mathsf T}.
\]
于是 \(|\det G|^2=1\)。复正交型 \(z^{\mathsf T}w\) 对两个变量都线性，\(h_0\) 却在第一变量带有共轭。这个差别已经在乘 \(i\) 时显现：
\[
(iz)^{\mathsf T}(iw)=-z^{\mathsf T}w,\qquad
\overline{iz}^{\mathsf T}(iw)=\overline z^{\mathsf T}w.
\]
所以 \(iI_n\) 保持 \(h_0\)，但当 \(n>0\) 时不属于 \(\operatorname O(I_n)\)。
\symidx{\(\operatorname U(n)\)、\(G^*=\overline G^{\mathsf T}\)：标准酉群与共轭转置}

\ProseParagraphBreak
二维辛群还有一个直接可算的描述。若 \(G=\begin{psmallmatrix}a&b\\c&d\end{psmallmatrix}\)，相乘得到
\[
G^{\mathsf T}\begin{pmatrix}0&1\\-1&0\end{pmatrix}G
=(ad-bc)\begin{pmatrix}0&1\\-1&0\end{pmatrix}.
\]
所以 \(\operatorname{Sp}_2(k)=\operatorname{SL}_2(k)\)，且这个等式在特征二中仍成立。它来自这个二维计算，不能把正交群与辛群的一般定义混为一谈。

\subsection{Hermitian 型、伴随与型的相似变换}
\label{b2:subsec:1004:02}

\begin{definition}\label{b2:def:hermitian-form}
设域 \(K\) 配有一个\term[yu-duihe]{对合}[involution]，即自同构 \(\sigma\) 满足 \(\sigma^2=\operatorname{id}\)，并记 \(\overline a=\sigma(a)\)。给定 \(K\) 向量空间 \(V\)，如果双可加映射 \(h:V\times V\to K\) 满足
\[
h(au,bv)=\overline a\,b\cdot h(u,v)
\]
且 \(h(v,u)=\overline{h(u,v)}\) 对所有 \(a,b\in K\)、\(u,v\in V\) 成立，就称 \(h\) 为 \(V\) 上的\term[Hermitian-xing]{Hermitian 型}[Hermitian form]（也称厄米型）。
\termidx[emi-xing]{厄米型}
若 \(h(v,w)=0\) 对全部 \(w\) 成立只能有 \(v=0\)，称 \(h\) 非退化。保持 \(h\) 的 \(K\) 线性自同构组成的群记为 \(\operatorname U(V,h)\)。
\end{definition}

恒等对合给出对称双线性型；复共轭给出通常的 Hermitian 型。本书固定第一变量共轭线性、第二变量线性的约定。若 \(H=\bigl(h(e_i,e_j)\bigr)\)，则
\[
h(u,v)=\overline x^{\mathsf T}Hy,\quad H^*=H,
\qquad H'=P^*HP.
\]
最后一个公式由 \(x=Px',y=Py'\) 直接代入得到。对角元都属于固定域 \(K^\sigma=\{\,a\in K\mid\sigma(a)=a\,\}\)：恒等对合的固定域为 \(K\)，复共轭的固定域则为 \(\mathbb R\)。保持型的矩阵满足 \(G^*HG=H\)。

\ProseParagraphBreak
分离一维正交分量需要先找到 \(h(v,v)\ne0\) 的向量。若对合恒等且特征为二，非零型却可能是交替型；例如 \(K^2\) 上的 \(h(u,v)=u_1v_2+u_2v_1\) 非零，但每个 \(h(v,v)\) 都为零。这样的型须使用辛型的二维块，不能从非零长度的直线开始分离。这正是下面假设所排除的情形。

\begin{proposition}\label{b2:prop:hermitian-diagonalization}
设 \(K\) 为配有对合 \(\sigma\) 的域，\(K^\sigma=\{a\in K\mid\sigma(a)=a\}\)，\(V\) 为有限维 \(K\) 向量空间，\(h:V\times V\to K\) 为 Hermitian 型。若 \(\sigma\ne\operatorname{id}\)，或 \(\operatorname{char}K\ne2\)，则 \(h\) 有正交基，且对角系数属于 \(K^\sigma\)。在 \(K=\mathbb C\)、\(\sigma\) 为复共轭的情形，对角系数可以归一化为一、负一与零，其个数唯一。
\end{proposition}
\begin{proof}
先证明非零的型必有 \(h(v,v)\ne0\) 的向量。反设全部对角取值为零。令 \(c=h(u,v)\)，从 \(h(u+v,u+v)=0\) 得 \(c+\overline c=0\)。再对 \(u+av\) 作同样展开，得到
\[
ac+\overline a\,\overline c=(a-\overline a)c=0.
\]
若对合非恒等，选择 \(a\ne\overline a\) 便得 \(c=0\)。若对合恒等而特征不为二，第一式就是 \(2c=0\)，也得 \(c=0\)。任意 \(u,v\) 的配对均为零，与型非零矛盾。

若型为零，任意基均正交。否则取 \(e\) 使 \(d=h(e,e)\ne0\)。对每个 \(v\)，向量 \(v-\dfrac{h(e,v)}d\,e\) 与 \(e\) 正交，因为第二变量线性给出 \(h(e,v)-h(e,v)d/d=0\)；由 Hermitian 对称性，两种次序的配对都为零。又 \(Ke\) 与其正交补的交为零，故分离出这个非退化一维分量，再按维数归纳。对角元被 \(\sigma\) 固定，已经由 Hermitian 对称性保证。

在复数情形，这些系数均为实数，缩放基向量可把非零系数改为正负一。若正项空间维数为 \(p_+\)，任意正定复子空间投影到它都单射，因为投影为零的向量的二次取值非正。故 \(p_+\) 是正定子空间最大复维数。负项的个数对 \(-h\) 作同样论证确定，零项数由总维数减去二者得到。这与\cref{b2:thm:sylvester-inertia}中的论证相同，只把实坐标平方换成复坐标模平方。
\end{proof}

非退化复 Hermitian 型的惯性为 \((p_+,p_-,0)\) 时，对应群常记为 \(\operatorname U(p_+,p_-)\)，其标准矩阵为 \(\operatorname{diag}(I_{p_+},-I_{p_-})\)。标准酉群是 \(p_-=0\) 的情形。这里的“酉”只表示保持指定 Hermitian 型，不自动表示该型正定。标准酉群还可以完全用实线性变换描述：忘去 \(\mathbb C^n\) 的复标量后，保留实配对 \(\operatorname{Re}h_0\) 与乘 \(i\) 的实线性算子，就能恢复原来的复结构和配对。

\begin{proposition}\label{b2:prop:unitary-real-complex-structure}
设 \(n\ge0\)，\(h_0(z,w)=\sum_{j=1}^n\overline{z_j}w_j\) 为 \(\mathbb C^n\) 上的标准 Hermitian 型。将 \(\mathbb C^n\) 看成实空间，记 \(Jz=iz\)、\(g(z,w)=\operatorname{Re}h_0(z,w)\)，则 \(J^2=-I\)，而保持 \(h_0\) 的复线性自同构群 \(\operatorname U(n)\) 恰为保持 \(g\) 并与 \(J\) 交换的实线性自同构组成的群。
\end{proposition}
\begin{proof}
实线性映射与 \(J\) 交换，当且仅当它还保持乘 \(i\)，也就当且仅当它复线性。又由第一变量的共轭线性，\(h_0(Jz,w)=-i h_0(z,w)\)，所以
\[
g(Jz,w)=\operatorname{Im}h_0(z,w),\qquad
h_0(z,w)=g(z,w)+i\,g(Jz,w).
\]
若映射保持 \(h_0\)，它显然保持实部，且作为复线性映射与 \(J\) 交换。反过来，若 \(A\) 保持 \(g\) 并与 \(J\) 交换，则
\[
g(JAz,Aw)=g(AJz,Aw)=g(Jz,w),
\]
所以它也保持 \(h_0\) 的虚部，连同实部便保持 \(h_0\)。
\end{proof}

若实向量空间 \(W\) 上的实线性算子 \(J:W\to W\) 满足 \(J^2=-I_W\)，就称 \(J\) 为 \(W\) 上的一个\term[fushu-jiegou]{复数结构}[complex structure]：规定 \((a+bi)v=av+bJv\)，就使原实空间成为复向量空间。结合律由 \(J^2=-I_W\) 验证。若 \(W\) 的实维数有限，复维数也有限；一组复基 \(w_1,\ldots,w_r\) 对应实基 \(w_1,Jw_1,\ldots,w_r,Jw_r\)，所以 \(\dim_{\mathbb R}W=2\dim_{\mathbb C}W\) 必为偶数。上面的命题因而把酉群同时解释为一种实正交群的子群和一种复线性群。

\ProseParagraphBreak
型还可以用来给线性算子定义伴随。下面把对称型、交替型和 Hermitian 型一并处理；符号中的共轭在前两种情形取恒等，在第三种情形取指定的域对合。

\begin{proposition}\label{b2:prop:form-adjoint-operator}
设 \(K\) 为配有对合 \(\sigma:a\mapsto\bar a\) 的域，\(V\) 为有限维 \(K\) 向量空间，\(\varepsilon\in\{1,-1\}\)。设 \(b:V\times V\to K\) 为非退化的 \(\varepsilon\)-Hermitian 型，即 \(b\) 双可加，第一变量共轭线性、第二变量线性，并满足
\[
b(au,cv)=\bar a\,c\cdot b(u,v),\qquad
b(v,u)=\varepsilon\overline{b(u,v)}
\]
对所有 \(a,c\in K\)、\(u,v\in V\) 成立，且 \(b(u,v)=0\) 对所有 \(v\) 成立只能有 \(u=0\)。对每个 \(K\) 线性算子 \(T:V\rightarrow V\)，存在唯一的 \(K\) 线性算子 \(T^\dagger:V\rightarrow V\) 满足
\[
b(Tu,v)=b(u,T^\dagger v)\qquad(u,v\in V).
\]
称 \(T^\dagger\) 为 \(T\) 关于 \(b\) 的\term[bansui-suanzi]{伴随算子}[adjoint operator]。选定基，令 \(H=(b(e_i,e_j))\)，并以相同字母表示算子的矩阵，则 \(H\) 可逆、\(H^*=\varepsilon H\)，其中 \(X^*=\overline X^{\mathsf T}\)。对任意另一 \(K\) 线性算子 \(S:V\rightarrow V\)，有
\[
T^\dagger=H^{-1}T^*H,\qquad
(ST)^\dagger=T^\dagger S^\dagger,\qquad
(T^\dagger)^\dagger=T.
\]
伴随运算可加，并满足 \((aT)^\dagger=\bar aT^\dagger\)。对任意子空间 \(W\le V\)，记 \(W^\perp=\{\,v\in V\mid b(w,v)=0\text{ 对所有 }w\in W\,\}\)，则
\[
(\operatorname{im}T)^\perp=\ker T^\dagger,\qquad
\operatorname{rank}T^\dagger=\operatorname{rank}T.
\]
\end{proposition}
\begin{proof}
固定第二变量时，\(u\mapsto b(u,v)\) 是共轭线性泛函，不能在非恒等对合下把它当作普通 \(V^*\) 中的元素。令
\[
C_\sigma(V,K)=
\{\,\ell:V\to K\mid \ell\text{ 可加，且 }\ell(au)=\bar a\,\ell(u)
\ (a\in K,\ u\in V)\,\}
\]
按逐点加法与标量乘法组成 \(K\) 向量空间。选任意一组 \(V\) 的基，共轭线性泛函由它在基上的值唯一确定，这些值又可任意指定，所以 \(\dim_K C_\sigma(V,K)=\dim_KV\)。映射
\[
\Psi:V\longrightarrow C_\sigma(V,K),\qquad
w\longmapsto \bigl(u\longmapsto b(u,w)\bigr)
\]
由第二变量线性而 \(K\) 线性。若 \(\Psi(w)=0\)，则 \(b(u,w)=0\) 对所有 \(u\) 成立，由 \(\varepsilon\)-Hermitian 对称性也有 \(b(w,u)=0\)，非退化性给出 \(w=0\)。有限维且等维使 \(\Psi\) 成为同构。

对每个 \(v\in V\)，\(\ell_v(u)=b(Tu,v)\) 属于 \(C_\sigma(V,K)\)，故存在唯一 \(w\) 使 \(\Psi(w)=\ell_v\)。定义 \(T^\dagger v=w\)，便得到要求的配对等式。又 \(v\mapsto\ell_v\) 是 \(K\) 线性，后接 \(\Psi^{-1}\) 说明 \(T^\dagger\) 线性；每个 \(v\) 所对应的 \(w\) 唯一也给出算子唯一性。这与\cref{b2:prop:finite-perfect-pairing}的有限维表示机制相同，只是第一变量的泛函改为共轭线性泛函。

选定基以后，要求的配对等式成为
\[
\overline u^{\mathsf T}T^*Hv
=\overline u^{\mathsf T}HT^\dagger v.
\]
对全部基向量 \(u,v\) 成立，当且仅当 \(T^*H=HT^\dagger\)。非退化性使 \(H\) 可逆，故得到矩阵公式。复合反序也可直接从配对看出：
\[
b(STu,v)=b(Tu,S^\dagger v)
=b(u,T^\dagger S^\dagger v).
\]
由唯一性，\((ST)^\dagger=T^\dagger S^\dagger\)。矩阵中的共轭转置也反转乘法，与这个内在公式一致。再由 \(H^*=\varepsilon H\) 以及 \(\varepsilon^2=1\)，有
\[
(T^\dagger)^\dagger
=H^{-1}H^*T(H^*)^{-1}H=T.
\]
可加性和标量公式同样直接来自共轭转置。

向量 \(v\) 属于 \((\operatorname{im}T)^\perp\)，当且仅当 \(b(Tu,v)=0\) 对所有 \(u\) 成立；伴随等式把这个条件改写为 \(b(u,T^\dagger v)=0\) 对所有 \(u\) 成立。由第二变量的非退化性，这恰好等价于 \(T^\dagger v=0\)，证明正交补与核的等式。

对任意 \(W\le V\)，限制映射 \(C_\sigma(V,K)\to C_\sigma(W,K)\) 满射：把 \(W\) 的一组基扩充为 \(V\) 的基，并在新增基向量上任意指定值，即可延伸共轭线性泛函。因此 \(\Psi\) 后接限制的映射 \(V\to C_\sigma(W,K)\) 满射，核为 \(W^\perp\)，故
\(\dim W^\perp=\dim V-\dim W\)。取 \(W=\operatorname{im}T\)，并对 \(T^\dagger\) 使用秩与核的维数公式，得到
\[
\operatorname{rank}T^\dagger
=\dim V-\dim\ker T^\dagger
=\dim\operatorname{im}T=\operatorname{rank}T.
\]
\end{proof}
\symidx{\(T^\dagger\)：相对于指定非退化型的伴随算子}

例如，实标准内积给出 \(T^\dagger=T^{\mathsf T}\)，复标准 Hermitian 型给出 \(T^\dagger=T^*\)。对辛矩阵 \(J_m\)，则有 \(T^\dagger=J_m^{-1}T^{\mathsf T}J_m\)。这都是整个线性算子的伴随，不是由代数余子式组成的伴随矩阵。这个反转乘法的运算将在\secref{b2:sec:B03}作为代数上的对合进一步讨论。

\ProseParagraphBreak
算子换基用矩阵相似 \(A\mapsto P^{-1}AP\)，型的换基则用合同 \(H\mapsto P^*HP\)；二者都是重新描述同一个对象。型的换基不要求 \(P\) 保持原矩阵 \(H\)。若把可逆矩阵 \(g\) 看作固定空间上的变换，等距条件才是 \(g^*Hg=H\)；以下的型相似变换则只要求 \(g^*Hg=cH\)，即把全部配对值同时乘以一个非零标量。它与算子的矩阵相似是不同概念。

\begin{definition}\label{b2:def:form-similitude}
设 \(K\) 为带域对合 \(a\mapsto\bar a\) 的域，固定域为 \(F=\{\,a\in K\mid\bar a=a\,\}\)，\(V=K^n\)、\(n\ge1\)。取可逆矩阵 \(H\) 满足 \(H^*=\varepsilon H\)，其中 \(X^*=\overline X^{\mathsf T}\)、\(\varepsilon\in\{1,-1\}\)，并令 \(b(u,v)=\overline u^{\mathsf T}Hv\)。若 \(g\in\operatorname{GL}_K(V)\) 满足
\[
b(gu,gv)=c\cdot b(u,v)\qquad(u,v\in V)
\]
对某个 \(c\in F^\times\) 成立，就称 \(g\) 为一个\term[xing-xiangsi-bianhuan]{型的相似变换}[similitude]，称 \(c\) 为其\term[xiangsi-chengzi]{乘子}[multiplier]，记为 \(\mu(g)\)。全部这类变换组成相似变换群。共轭为恒等且特征不为二时，对称型的相似变换群记为 \(\operatorname{GO}(b)\)；共轭为恒等时，交替型的相似变换群记为 \(\operatorname{GSp}(b)\)；对给定非平凡域对合，Hermitian 型的相似变换群记为 \(\operatorname{GU}(b)\)。特征二中的二次型相似条件须直接用二次型定义，不能只由其极化型判断。
\end{definition}

型非退化且 \(n\ge1\)，保证可以取 \(u,v\) 使 \(z=b(u,v)\ne0\)。若 \(c,c'\) 都满足相似条件，则 \(cz=c'z\)，在域中消去非零的 \(z\)，得到 \(c=c'\)，所以乘子唯一。

\ProseParagraphBreak
即使先只要求 \(c\in K^\times\)，交换这两个变量也会迫使它属于固定域。一方面，
\[
b(gv,gu)=c\,b(v,u)=c\varepsilon\bar z;
\]
另一方面，由型的对称性，
\[
b(gv,gu)=\varepsilon\overline{b(gu,gv)}
=\varepsilon\bar c\,\bar z.
\]
因 \(\varepsilon\bar z\ne0\)，这两个等式给出 \(c=\bar c\)。

\begin{proposition}\label{b2:prop:similitude-multiplier}
设 \(K\) 为带域对合 \(a\mapsto\bar a\) 的域，\(F\) 为其固定域，\(V=K^n\)、\(n\ge1\)，且 \(b(u,v)=\overline u^{\mathsf T}Hv\)，其中 \(H\) 可逆、\(H^*=\varepsilon H\)、\(\varepsilon\in\{1,-1\}\)，星号表示共轭转置。型 \(b\) 的相似变换组成群，乘子为到 \(F^\times\) 的群同态，其核恰为保持 \(b\) 的等距群。记 \(g^\dagger=H^{-1}g^*H\)，则对 \(g\in\operatorname{GL}_n(K)\)、\(c\in F^\times\)，有
\[
g\text{ 是乘子为 }c\text{ 的相似变换}
\quad\Longleftrightarrow\quad
g^*Hg=cH
\quad\Longleftrightarrow\quad
g^\dagger g=cI_n.
\]
\end{proposition}
\begin{proof}
两个相似变换依次作用时，配对值乘以两个乘子的积；恒等变换的乘子为一。把 \(u,v\) 换成 \(g^{-1}u,g^{-1}v\)，可知逆变换的乘子为 \(\mu(g)^{-1}\)。这证明群结构及同态性。乘子为一恰好是等距条件。把配对写成矩阵，再使用 \(g^\dagger=H^{-1}g^*H\)，便得两个等价式。
\end{proof}

乘子的核已经确定，接下来还需计算它的实际像，而这个像受型的性质限制。在非零空间上的正定实型 \(b(u,v)=u^{\mathsf T}v\) 上，非零 \(v\) 给出 \(\mu(g)=b(gv,gv)/b(v,v)>0\)，而 \(\sqrt c I_n\) 实现每个 \(c>0\)，故像恰为 \(\mathbb R_{>0}\)。非零空间上的复正定 Hermitian 型也有相同的乘子像。另一方面，对标准辛型取
\[
g_t=\operatorname{diag}(tI_m,I_m),\qquad t\in K^\times.
\]
直接按块相乘，得到 \(g_t^{\mathsf T}J_mg_t=tJ_m\)，所以当 \(m\ge1\) 时，辛相似变换的乘子取遍 \(K^\times\)。因此从核得到的正合序列应以乘子的实际像结束，不能无条件把它写成到全部非零标量的满射。

\OptionalSubsection{Clifford 代数的构造}{b2:subsec:1004:03}

下面恢复任意域 \(k\)，并令 \(q\) 按\cref{b2:def:quadratic-form}定义。张量代数保留向量空间的线性关系而不另加乘法关系；现在要求向量 \(v\) 的像满足平方等式 \(v^2=q(v)1\)，就是把二次型的取值直接写进乘法。将这些等式的差生成双边理想，便使它们在左右乘法下始终成立。构造允许型退化，也允许特征二；此时仍使用未经除以二的 \(\beta_q\)。

\begin{definition}\label{b2:def:clifford-algebra}
设 \(k\) 为任意域，\(V\) 为 \(k\) 向量空间，\(q:V\to k\) 为二次型。取张量代数 \(T_k(V)\)，令 \(I_q\) 为全部 \(v\otimes v-q(v)1\)、\(v\in V\) 生成的双边理想。商代数
\[
\operatorname{Cl}(V,q)=T_k(V)/I_q
\]
称为 \(q\) 的\term[Clifford-daishu]{Clifford 代数}[Clifford algebra]。从 \(V\) 到这个商代数的自然线性映射记为 \(\iota\)。
\end{definition}
\symidx{\(\operatorname{Cl}(V,q)\)：二次型的 Clifford 代数}

若 \(q=0\)，关系理想恰由所有 \(v\otimes v\) 生成。由\cref{b2:prop:exterior-algebra-universal}的张量商描述，得到自然代数同构
\[
\operatorname{Cl}(V,0)\cong\bigwedge_k V.
\]
因此外代数就是所有一次元素平方为零的 Clifford 代数；一般二次型把这些平方从零改为相应的标量。

\begin{theorem}[Clifford 代数的泛性质]\label{b2:thm:clifford-universal}
设 \(k\) 为任意域，\(V\) 为 \(k\) 向量空间，\(q:V\to k\) 为二次型，\(\iota:V\to\operatorname{Cl}(V,q)\) 为自然映射。对任意结合含幺 \(k\) 代数 \(A\)，若 \(k\) 线性映射 \(f:V\to A\) 满足 \(f(v)^2=q(v)1_A\) 对每个 \(v\in V\) 成立，则存在唯一含幺 \(k\) 代数同态
\[
\widetilde f:\operatorname{Cl}(V,q)\longrightarrow A,
\qquad \widetilde f\iota=f.
\]
此外在 Clifford 代数中有
\begin{equation}\label{b2:eq:clifford-polar-relation}
\iota(u)\iota(v)+\iota(v)\iota(u)=\beta_q(u,v)1.
\end{equation}
\end{theorem}
\begin{proof}
\cref{b2:thm:tensor-algebra-universal}将 \(f\) 唯一延伸为 \(T_k(V)\to A\) 的含幺代数同态。题设等式使 \(I_q\) 的每个生成元映到零，双边理想的全部元素也就映到零，故它通过商代数分解。反过来，任何这样的分解都给出原来的张量代数同态，所以唯一。

将 \(\iota(u+v)^2=q(u+v)1\) 展开，减去 \(\iota(u)^2=q(u)1\) 与 \(\iota(v)^2=q(v)1\)，便得到\eqref{b2:eq:clifford-polar-relation}。特征不为二时，右侧为 \(2b_q(u,v)1\)。
\end{proof}

关系 \(\iota(v)^2=q(v)1\) 记录向量的二次取值；特征不为二时，\eqref{b2:eq:clifford-polar-relation}还说明，若 \(b_q(u,v)=0\)，则 \(\iota(u)\iota(v)=-\iota(v)\iota(u)\)。乘法同时记录了二次取值与正交条件。

\ProseParagraphBreak
上述商构造给出了自然线性映射 \(\iota\) 及泛性质。\(\iota\) 的单射性和代数的维数还需由基的结构证明；在特征不为二的有限维情形，\cref{b2:thm:B-clifford-basis}给出其单射性及维数 \(2^{\dim V}\)。下面直接计算一个低维模型。

\begin{example}\label{b2:exm:clifford-hyperbolic-plane}
设 \(k\) 为特征不为二的域，\(H=ke\oplus kf\) 上的二次型为 \(q(xe+yf)=2xy\)，则 \(\operatorname{Cl}(H,q)\cong M_2(k)\)。
\end{example}
\begin{proof}
以双曲基 \(e,f\) 的像仍记为 \(e,f\)，关系为
\[
e^2=f^2=0,\qquad ef+fe=2.
\]
先说明 \(1,e,f,ef\) 生成整个代数的向量空间。任何字若出现相邻相同字母便为零；否则字母交替。长度至少三时含 \(efe\) 或 \(fef\)，而
\[
efe=e(2-ef)=2e,\qquad fef=(2-ef)f=2f.
\]
递减字长后只剩长度至多二的字，再用 \(fe=2-ef\) 即得四个所列元素。

为找到满足这些关系的矩阵，先取矩阵单位 \(E_{12},E_{21}\)。它们的平方均为零，且 \(E_{12}E_{21}+E_{21}E_{12}=I_2\)；把第一个矩阵乘二，就能实现右端的 \(2I_2\)。因此取
\[
E=\begin{pmatrix}0&2\\0&0\end{pmatrix},\qquad
F=\begin{pmatrix}0&0\\1&0\end{pmatrix}.
\]
它们满足 \(E^2=F^2=0\)、\(EF+FE=2I_2\)，因此线性映射 \(xe+yf\mapsto xE+yF\) 的平方为 \(2xyI_2\)。由\cref{b2:thm:clifford-universal}得到代数同态到 \(M_2(k)\)。\(I_2,E,F,EF\) 线性无关：两个非对角位置先确定 \(E,F\) 的系数为零，两个对角位置再确定余下系数为零，所用的二可逆。它们构成四维矩阵空间的基。来源已由四个元素生成，故同态既单射又满射。
\end{proof}

在这个例子里，二次关系产生了一个非交换矩阵代数。关系 \(v\otimes v-q(v)1\) 混合张量次数二与零，所以商代数一般不继承张量次数的自然数分次；将次数模二以后，两项却同属偶次。因而 \(I_q\) 对奇偶分次仍是齐次理想，商代数继承直和分解
\[
\operatorname{Cl}(V,q)=\operatorname{Cl}^0(V,q)\oplus
\operatorname{Cl}^1(V,q),
\qquad
\operatorname{Cl}^i\operatorname{Cl}^j\subseteq
\operatorname{Cl}^{\,i+j\ (\mathrm{mod}\ 2)}.
\]
在特征不为二的非零有限维非退化二次空间上，\cref{b2:prop:B-spin-action}再用偶部分的可逆元构造自旋群到特殊正交群的作用，并算出其核。

\begin{proposition}\label{b2:prop:clifford-line-models}
设 \(k\) 为任意域，\(V=ke\) 为一维 \(k\) 向量空间，\(a\in k\)，二次型为 \(q(xe)=ax^2\)。则存在含幺 \(k\) 代数同构
\[
\operatorname{Cl}(V,q)\cong k[t]/(t^2-a),
\]
将 \(\iota(e)\) 送到 \(t\) 的剩余类。特别地，在 \(k=\mathbb R\) 时，
\[
\operatorname{Cl}(\mathbb R e,q)\cong
\begin{cases}
\mathbb R\times\mathbb R,&a>0,\\
\mathbb C,&a<0,\\
\mathbb R[\varepsilon]/(\varepsilon^2),&a=0.
\end{cases}
\]
\end{proposition}
\begin{proof}
令 \(A=k[t]/(t^2-a)\)，\(\tau=t+(t^2-a)\)。线性映射 \(f:V\to A\)、\(xe\mapsto x\tau\) 满足
\[
f(xe)^2=x^2\tau^2=ax^2\,1_A=q(xe)1_A.
\]
由\cref{b2:thm:clifford-universal}，得到含幺代数同态
\(\alpha:\operatorname{Cl}(V,q)\to A\)，使 \(\alpha(\iota(e))=\tau\)。

反过来，一生成元的自由结合代数是 \(k[t]\)，故赋值 \(t\mapsto\iota(e)\) 给出含幺代数同态 \(k[t]\to\operatorname{Cl}(V,q)\)。关系 \(\iota(e)^2=a1\) 使 \(t^2-a\) 映到零，因而同态通过商，得到
\(\beta:A\to\operatorname{Cl}(V,q)\)。复合 \(\alpha\beta\) 固定 \(\tau\) 与标量，故在由它们生成的 \(A\) 上为恒等；复合 \(\beta\alpha\) 固定 \(\iota(e)\) 与标量，而它们生成整个 Clifford 代数，故这个复合也为恒等。这证明所述同构，对特征二同样适用。

首一二次多项式的带余除法说明，\(A\) 中每个元素唯一写成 \(x+y\tau\)。若 \(a>0\)，令 \(r=\sqrt a\)，赋值 \(\tau\mapsto(r,-r)\) 给出 \(\mathbb R\) 代数同态
\[
A\longrightarrow\mathbb R\times\mathbb R,\qquad
x+y\tau\longmapsto(x+yr,x-yr).
\]
它的逆将 \((u,v)\) 送到 \(\dfrac{u+v}2+\dfrac{u-v}{2r}\tau\)，所以是同构。若 \(a<0\)，令 \(r=\sqrt{-a}\)，赋值 \(\tau\mapsto ir\) 给出
\[
A\longrightarrow\mathbb C,\qquad x+y\tau\longmapsto x+iyr.
\]
因 \(r\ne0\)，每个复数唯一有这种表达，故也是同构。若 \(a=0\)，只需把 \(\tau\) 政名为 \(\varepsilon\)，关系就是 \(\varepsilon^2=0\)，得到双数代数。
\end{proof}

本章的 Witt 群若要连同张量积一起计算，可在读完\secref{b2:sec:1003}后转到\secref{b2:sec:B02}；型的伴随算子同代数对合的联系则见\secref{b2:sec:B03}，那一节还需要\chapref{b2:ch:17}的中心单代数与四元数。


\begin{chaptersummary}
配对使向量与泛函相联系；有限维非退化性把这种联系变成同构，继而给出正交补的维数公式。只有当子空间上的限制型非退化时，正交补才直接给出直和。型的矩阵按合同变化，算子的矩阵按相似变化，二者的分类不变量不同。

\ProseParagraphBreak
交替型由空间维数与秩共同分类。特征不为二的二次型可以对角化，但进一步的分类取决于域：实数域上，正负惯性指数与根空间维数共同分类；复数域上，空间维数与秩共同分类。一般域上，非零各向同性向量使双曲平面得以分离；反射保证共同的非退化分量可以消去，从而确定唯一的各向异性部分。Witt 群进一步忽略双曲部分，保留二次型在稳定等距下的信息。

\ProseParagraphBreak
正交群、辛群与酉群分别保持指定的型。伴随算子把保持型的条件写成 \(g^\dagger g=I\)，型的相似变换则允许右端为非零标量；这个标量的实际取值范围依赖底域和型。Hermitian 型的共轭作用不能从公式中省略，复数结构则把酉群与实正交群联系起来。Clifford 代数将二次取值转为乘法关系；它的泛性质来自张量代数及双边理想商，具体维数结论仍须具体证明。
\end{chaptersummary}

\BookExercises

\begin{exercise}\label{b2:ex:10-left-right-orthogonals}
设 \(k\) 为域，在 \(k^2\) 上取 \(b(x,y)=x_1y_1+x_1y_2+x_2y_2\)。记 \(e_1,e_2\) 为标准基。证明型非退化，并对 \(W=ke_1\) 分别求 \(W^{\perp_r}\) 与 \({}^{\perp_l}W\)。说明为什么一般型不能只写一个没有侧别的正交补。
\end{exercise}
\begin{solution}
矩阵为 \(\begin{psmallmatrix}1&1\\0&1\end{psmallmatrix}\)，行列式为一，故非退化。\(b(e_1,y)=y_1+y_2\)，所以 \(W^{\perp_r}=k(1,-1)\)；而 \(b(x,e_1)=x_1\)，所以 \({}^{\perp_l}W=ke_2\)。两条直线不同，在特征二中也不同。非退化保证两边的维数正确，却不把不同变量中的正交条件变成同一个条件。
\end{solution}

\begin{exercise}\label{b2:ex:10-degenerate-orthogonals}
设 \(k\) 为域，\(V\) 为有限维 \(k\) 向量空间，\(b\) 为 \(V\) 上的对称型，允许退化，令 \(R=\operatorname{rad}b\)。证明对任意 \(W\le V\)，
\[
\dim W^\perp=\dim V-\dim W+\dim(W\cap R),
\qquad (W^\perp)^\perp=W+R.
\]
\end{exercise}
\begin{solution}
型在商空间 \(\overline V=V/R\) 上良定义：改变任一代表元所添的项都来自根空间，配对不变。所得型非退化，因为与全部商类正交的代表元正属于 \(R\)。记 \(\overline W=(W+R)/R\)，则 \(W^\perp/R=\overline W^\perp\)。由\cref{b2:thm:orthogonal-complement}，后者的维数为 \(\dim\overline V-\dim\overline W\)。代入 \(\dim\overline W=\dim W-\dim(W\cap R)\)，得到第一式。再次取正交补，商空间中得到 \(\overline W\)，取原像便是 \(W+R\)，证明第二式。
\end{solution}

\begin{exercise}\label{b2:ex:10-symplectic-basis}
在任意域 \(k\) 上，四维交替型的矩阵为
\[
B=\begin{pmatrix}0&1&1&0\\-1&0&0&0\\-1&0&0&1\\0&0&-1&0\end{pmatrix}.
\]
找出两个互相正交的辛平面，并给出辛基。说明该计算在特征二中是否仍有效。
\end{exercise}
\begin{solution}
第一对取 \(e_1,e_2\)。向量 \(e_3-e_2\) 与这两者的配对均为零，\(e_4\) 也与这两者正交，而 \(b(e_3-e_2,e_4)=1\)。所以 \(e_1,e_2,e_3-e_2,e_4\) 是按两维块排列的辛基；这四个向量线性无关，由原标准基到它们的变换可逆。只用了加减和单位系数，没有除以二，故在特征二中同样成立。
\end{solution}

\begin{exercise}\label{b2:ex:10-one-dimensional-square-class}
分类 \(\mathbb F_5\) 上全部非退化一维二次型，并从 \(\langle1\rangle,\langle2\rangle,\langle3\rangle,\langle4\rangle\) 中选出每个等距类的代表。再比较 \(\langle1\rangle\) 与 \(\langle2\rangle\) 在 \(\mathbb Q\) 和 \(\mathbb Q(\sqrt2)\) 上的等距性，说明扩大底域怎样影响平方类判据。
\end{exercise}
\begin{solution}
由\cref{b2:thm:quadratic-diagonalization}中的一维判据，只须分类非零系数的平方类。\(\mathbb F_5^\times\) 中的平方为 \(1,4\)，另一个平方类为 \(2,3\)，所以可取 \(\langle1\rangle\) 与 \(\langle2\rangle\) 作代表；\(\langle4\rangle\cong\langle1\rangle\)、\(\langle3\rangle\cong\langle2\rangle\)。

在有理数域上 \(2\) 不是平方：若互素整数满足 \(m^2=2n^2\)，则先得 \(m\) 偶，再得 \(n\) 偶，矛盾。因此两型不等距。扩大到 \(\mathbb Q(\sqrt2)\) 后，\(2=(\sqrt2)^2\)，两型等距。原域中不同的平方类可以在扩域后合并，原有等距关系则仍被保留。
\end{solution}

\begin{exercise}\label{b2:ex:10-diagonalization-radical}
在 \(\mathbb R^3\) 上，对
\[
q(x,y,z)=x^2+2xy+2y^2+2yz+z^2
\]
给出可逆坐标代换，使它对角化；求惯性、根空间，以及使限制型非退化的一个补空间。
\end{exercise}
\begin{solution}
令 \(u=x+y,v=y+z,w=y\)，逆变换为 \(x=u-w,y=w,z=v-w\)。于是 \(q=u^2+v^2\)，惯性为 \((2,0,1)\)。根空间为 \(u=v=0\)，即 \(\mathbb R(-1,1,-1)\)。取 \(w=0\)，得到 \(y=0\) 的平面，其限制型为 \(x^2+z^2\)，非退化；它与上述根空间构成直和。
\end{solution}

\begin{exercise}\label{b2:ex:10-inertia-parameter}
对实参数 \(t\)，分类 \(q_t(x,y,z)=x^2+2xy+ty^2-z^2\) 的惯性，并求出它何时非退化、何时具有非零各向同性向量。
\end{exercise}
\begin{solution}
写为 \((x+y)^2+(t-1)y^2-z^2\)。当 \(t>1,t=1,t<1\) 时，惯性依次为 \((2,1,0),(1,1,1),(1,2,0)\)。所以非退化恰在 \(t\ne1\) 时成立。无论 \(t\) 如何，\((1,0,1)\) 都是非零零点，故总有非零各向同性向量。
\end{solution}

\begin{exercise}\label{b2:ex:10-char-two-same-polar}
在 \(\mathbb F_2^2\) 上取 \(q_0(x,y)=xy\)、\(q_1(x,y)=x^2+xy+y^2\)。分别求保持 \(q_0\)、\(q_1\) 的线性自同构群，并与保持它们共同极化型的群比较。
\end{exercise}
\begin{solution}
\cref{b2:prop:quadratic-char-two-polar-failure}已经给出三条非零向量上的取值。保持 \(q_0\) 的变换必须保持零点集合 \(\{e_1,e_2\}\)，因而只能固定或交换这两个基向量；两者都保持 \(q_0\)，所得群为 \(\{I_2,\begin{psmallmatrix}0&1\\1&0\end{psmallmatrix}\}\cong C_2\)。\(q_1\) 在每个非零向量上都取一，故它由全部 \(\operatorname{GL}_2(\mathbb F_2)\) 保持。选择第一列有三种，第二列须与第一列线性无关，有两种，因此该群有六个元素；它忠实地置换三条非零向量，所以同构于 \(S_3\)。

共同的极化型矩阵为 \(J=\begin{psmallmatrix}0&1\\1&0\end{psmallmatrix}\)。二维计算给出 \(G^{\mathsf T}JG=(\det G)J\)，而 \(\mathbb F_2\) 中每个可逆矩阵的行列式都是一，故其保持群也是 \(\operatorname{GL}_2(\mathbb F_2)\)。保持极化型的变换未必保持原二次型，\(q_0\) 的群在这里是严格较小的子群。
\end{solution}

\begin{exercise}\label{b2:ex:10-binary-isotropy}
设 \(k\) 为特征不为二的域，\(d\in k^\times\)。证明 \(\langle1,-d\rangle\) 各向异性，当且仅当 \(d\) 不是平方；若 \(d\) 是平方，写出一组双曲基。
\end{exercise}
\begin{solution}
非零零点满足 \(x^2=dy^2\)。若 \(y=0\)，则 \(x=0\)，故必须 \(y\ne0\)，得到 \(d=(x/y)^2\)。反过来，若 \(d=c^2\)、\(c\ne0\)，取
\[
e=(c,1),\qquad f=(1/(2c),-1/(2c^2)).
\]
直接得到 \(q(e)=q(f)=0\) 及 \(b(e,f)=1\)，故为双曲基。
\end{solution}

\begin{exercise}\label{b2:ex:10-hyperbolic-split-explicit}
设 \(k\) 为特征不为二的域，在 \(k^3\) 上取 \(q(x,y,z)=2xy+z^2\)。以各向同性向量 \(e=(1,-1/2,1)\) 开始，取 \(f_0=(0,0,1)\)，再将它调整为与 \(e\) 配对为一的各向同性向量 \(f\)。实际完成一次双曲平面的分离，写出正交补的基和限制型。
\end{exercise}
\begin{solution}
相应对称型为 \(b\bigl((x,y,z),(u,v,w)\bigr)=xv+yu+zw\)。先有 \(b(e,f_0)=1\)、\(q(f_0)=1\)，于是按\cref{b2:lem:hyperbolic-plane-split}的修正方法取
\[
f=f_0-\frac{q(f_0)}2e=(-1/2,1/4,1/2).
\]
它满足 \(q(f)=0\)、\(b(e,f)=1\)。向量 \((x,y,z)\) 同时与 \(e,f\) 正交的条件为 \(-x+2y+2z=0\)、\(x-2y+2z=0\)。相加先得 \(z=0\)，再得 \(y=x/2\)，故正交补由 \(g=(1,1/2,0)\) 生成，且 \(q(g)=1\)。因此 \(e,f,g\) 给出 \(H\perp\langle1\rangle\) 的具体分解；这里调整 \(f_0\) 的系数由各向同性条件决定。
\end{solution}

\begin{exercise}\label{b2:ex:10-two-reflections}
在实三维型 \(q(x,y,z)=x^2+y^2-z^2\) 中，令 \(u=(1,0,0)\)、\(v=(1,1,1)\)。解释为何沿 \(u-v\) 的反射没有定义，并用两个有定义的反射把 \(u\) 送到 \(v\)。
\end{exercise}
\begin{solution}
\(q(u)=q(v)=1\)，但 \(q(u-v)=q(0,-1,-1)=0\)，反射公式的分母为零。另一方面 \(u+v=(2,1,1)\) 的二次取值为四，且 \(b(u,u+v)=2\)，故 \(\tau_{u+v}(u)=u-(u+v)=-v\)。\(q(v)=1\) 保证 \(\tau_v\) 有定义，并有 \(\tau_v(-v)=v\)。所求映射为 \(\tau_v\tau_{u+v}\)。
\end{solution}

\begin{exercise}\label{b2:ex:10-real-isotropic-dimension}
设有限维实向量空间上的二次型的惯性为 \((p_+,p_-,r)\)。证明其全各向同性子空间的最大维数为 \(\min(p_+,p_-)+r\)，并说明为什么根空间的贡献必须单列。
\end{exercise}
\begin{solution}
先在非退化标准型中讨论。若 \(T\) 全各向同性，它到正坐标空间和负坐标空间的两个投影都单射：其中一个投影为零，则一个非零向量的型值只能严格为正或严格为负，与零值矛盾。因此 \(\dim T\le\min(p_+,p_-)\)。把一条正坐标轴与一条负坐标轴配对，取各对的和向量，就得到达到上界的全各向同性子空间。

一般情形下，先除去根空间。\(T\) 的像全各向同性，故维数至多为 \(\min(p_+,p_-)\)，而核 \(T\cap\operatorname{rad}b\) 维数至多为 \(r\)。把前面构造的子空间与整个根空间直和即可达到总上界。根空间是零型，既不贡献正方向也不贡献负方向，所以它不是双曲分量，应另外加上其维数。
\end{solution}

\begin{exercise}\label{b2:ex:10-witt-binary-relation}
设 \(k\) 为特征不为二的域，\(a,b\in k\)，且 \(a,b,a+b\) 都非零。通过显式正交基证明
\[
\langle a,b\rangle\cong\left\langle a+b,\frac{ab}{a+b}\right\rangle.
\]
据此写出 \(W(k)\) 中的一条关系。
\end{exercise}
\begin{solution}
在原正交基 \(e,f\) 中取 \(u=e+f\)、\(v=(be-af)/(a+b)\)。两者的配对为 \((ab-ab)/(a+b)=0\)，且 \(q(u)=a+b\)、\(q(v)=(ab^2+a^2b)/(a+b)^2=ab/(a+b)\)。换基矩阵的行列式为负一，所以是基。这给出等距，并在 Witt 群中得到
\[
\bigl[\langle a\rangle\bigr]+\bigl[\langle b\rangle\bigr]
=\bigl[\langle a+b\rangle\bigr]+\left[\left\langle\frac{ab}{a+b}\right\rangle\right].
\]
最后两层定界符只用于区分一维型与它的 Witt 类。
\end{solution}

\begin{exercise}\label{b2:ex:10-real-complex-witt-map}
证明将实二次型的系数看成复数，诱导群同态 \(W(\mathbb R)\to W(\mathbb C)\)。在\cref{b2:prop:witt-real-complex}的识别下求出该同态及其核。
\end{exercise}
\begin{solution}
实等距矩阵在复数域仍可逆并满足同一个合同等式，双曲平面扩域后仍是双曲平面，所以 Witt 等价及类的加法都被保持。实惯性为 \((p_+,p_-,0)\) 的型，复维数为 \(p_++p_-\)，其奇偶性等于 \(p_+-p_-\) 的奇偶性。因此该同态就是整数模二，核为 \(2\mathbb Z\)。
\end{solution}

\begin{exercise}\label{b2:ex:10-orthogonal-one-dimension}
设 \(k\) 为特征不为二的域。求 \(k\) 上一维非退化二次型的正交群和特殊正交群。再设二维双曲平面的矩阵为 \(B_H=\begin{psmallmatrix}0&1\\1&0\end{psmallmatrix}\)，证明 \(\operatorname{diag}(t,t^{-1})\)、\(t\in k^\times\)，都在其特殊正交群中。
\end{exercise}
\begin{solution}
一维线性自同构为乘 \(c\ne0\)，保持型要求 \(c^2=1\)，故正交群为 \(\{1,-1\}\)，特殊正交群为 \(\{1\}\)。双曲平面上，变换 \((x,y)\mapsto(tx,t^{-1}y)\) 保持 \(2xy\)，且行列式为一，故满足所述条件。这个一参数族表明不定型的保持群允许沿两条各向同性轴作相反的伸缩。
\end{solution}

\begin{exercise}\label{b2:ex:10-first-order-isometries}
设 \(k\) 为域，\(n\ge1\)，\(R=k[\varepsilon]/(\varepsilon^2)\)，\(B\in M_n(k)\) 可逆，给定 \(X\in M_n(k)\)。证明 \(G=I_n+\varepsilon X\) 在 \(R\) 上保持 \(B\)，当且仅当 \(X^{\mathsf T}B+BX=0\)。\(\operatorname{char}k\ne2\) 时，分别计算 \(B=I_n\) 与 \(B=J_m\)（此时 \(n=2m\)）的解空间维数。
\end{exercise}
\begin{solution}
双数中的 \(\varepsilon^2=0\) 使保持型的等式只留下相对于 \(\varepsilon\) 的一次变化，因此所得条件将是关于 \(X\) 的线性方程。\(G\) 的逆为 \(I_n-\varepsilon X\)。展开 \(G^{\mathsf T}BG\)，得到 \(B+\varepsilon(X^{\mathsf T}B+BX)\)。由于 \(1,\varepsilon\) 在 \(k\) 上线性无关，保持 \(B\) 恰好等价于所述等式。\(B=I_n\) 时，\(X\) 反对称且对角元为零，独立参数为上三角部分的 \(n(n-1)/2\) 项。

若 \(B=J_m\)，写 \(X=\begin{psmallmatrix}A&C\\D&E\end{psmallmatrix}\)。分块相乘给出 \(E=-A^{\mathsf T}\)、\(C=C^{\mathsf T}\)、\(D=D^{\mathsf T}\)。所以维数为 \(m^2+2\cdot m(m+1)/2=m(2m+1)\)。
\end{solution}

\begin{exercise}\label{b2:ex:10-indefinite-hermitian}
在 \(\mathbb C^2\) 上取 Hermitian 矩阵 \(H=\begin{psmallmatrix}1&1\\1&0\end{psmallmatrix}\)。求其惯性并给出 \(P^*HP\) 的标准形。再证明
\[
A=\begin{pmatrix}5/3&4/3\\4/3&5/3\end{pmatrix}
\]
属于 \(\operatorname U(1,1)\)，却不属于 \(\operatorname U(2)\)。
\end{exercise}
\begin{solution}
对角取值为 \(|z_1+z_2|^2-|z_2|^2\)。取 \(P=\begin{psmallmatrix}1&-1\\0&1\end{psmallmatrix}\)，即 \(z_1=u-v,z_2=v\)，得到 \(P^*HP=\operatorname{diag}(1,-1)\)，惯性为 \((1,1,0)\)。

记 \(c=5/3,s=4/3\)，有 \(c^2-s^2=1\)。相乘得 \(A^{\mathsf T}\operatorname{diag}(1,-1)A=\operatorname{diag}(1,-1)\)，故 \(A\in\operatorname U(1,1)\)。但第一列的普通模平方为 \(c^2+s^2=41/9\ne1\)，所以 \(A^*A\ne I\)，不在标准酉群中。
\end{solution}

\begin{exercise}\label{b2:ex:10-unitary-real-blocks}
设 \(n\ge1\)，在 \(\mathbb R^{2n}\) 上令 \(J=\begin{psmallmatrix}0&-I_n\\I_n&0\end{psmallmatrix}\)。求所有与 \(J\) 交换的实矩阵的分块形式，并证明其中的正交矩阵恰对应复酉矩阵。
\end{exercise}
\begin{solution}
将四个 \(n\) 阶块代入 \(TJ=JT\)，得到且仅得到
\[
T=\begin{pmatrix}A&-B\\B&A\end{pmatrix}.
\]
它表示复矩阵 \(A+iB\) 的实线性作用。\(T^{\mathsf T}T=I\) 等价于 \(A^{\mathsf T}A+B^{\mathsf T}B=I\) 与 \(A^{\mathsf T}B=B^{\mathsf T}A\)，而这两式恰是 \((A+iB)^*(A+iB)=I\) 的实部和虚部。因此得到所述对应。
\end{solution}

\begin{exercise}\label{b2:ex:10-adjoint-image-kernel}
设 \(k\) 为特征不为二的域，在 \(k^2\) 上取矩阵为 \(H=\begin{psmallmatrix}0&1\\1&0\end{psmallmatrix}\) 的对称双线性型。证明 \(T=\begin{psmallmatrix}0&1\\0&0\end{psmallmatrix}\) 关于此型自伴随，且 \(T^2=0\) 而 \(T\ne0\)；求 \(\operatorname{im}T\)、\(\ker T\) 及 \((\operatorname{im}T)^\perp\)。再证明有限维实正定内积空间或复正定 Hermitian 空间上的自伴随算子若满足 \(T^2=0\)，就必为零，并指出两种情形的区别。
\end{exercise}
\begin{solution}
\(H^{-1}=H\)，相乘得到 \(H^{-1}T^{\mathsf T}H=T\)，故由\cref{b2:prop:form-adjoint-operator}，\(T^\dagger=T\)。又 \(T(e_2)=e_1\)、\(T(e_1)=0\)，所以 \(T^2=0\ne T\)，且 \(\operatorname{im}T=\ker T=ke_1\)。条件 \(b(e_1,v)=v_2=0\) 给出 \((\operatorname{im}T)^\perp=ke_1\)，也与正文的像、核恒等式一致。像与核在这里相等，它们并不构成直和。

在正定情形，自伴随性与 \(T^2=0\) 给出
\[
b(Tv,Tv)=b(v,T^\dagger Tv)=b(v,T^2v)=0.
\]
正定性迫使每个 \(Tv\) 为零，故 \(T=0\)。前一例的 \(ke_1\) 是全各向同性直线，非零向量也可有零取值；非退化性本身不能替代正定性。
\end{solution}

\begin{exercise}\label{b2:ex:10-indefinite-similitude-multipliers}
设 \(b\) 是 \(\mathbb R^2\) 上矩阵为 \(B=\operatorname{diag}(1,-1)\) 的双线性型。取 \(P=\begin{psmallmatrix}1&1\\0&1\end{psmallmatrix}\)，计算换基后的矩阵 \(P^{\mathsf T}BP\)，并判断 \(P\) 是否保持 \(b\)，或是否将它变为非零标量倍数。再求 \(b\) 的相似乘子的像，给出一个乘子为负一的变换，并说明正定型为何不可能有负乘子。
\end{exercise}
\begin{solution}
直接相乘得 \(P^{\mathsf T}BP=\begin{psmallmatrix}1&1\\1&0\end{psmallmatrix}\)。它既不等于 \(B\)，也不是 \(B\) 的标量倍数，所以 \(P\) 虽可用于换基，却不是 \(b\) 的等距变换或型的相似变换。交换两个坐标的矩阵 \(S=\begin{psmallmatrix}0&1\\1&0\end{psmallmatrix}\) 满足 \(S^{\mathsf T}BS=-B\)，是乘子为负一的型的相似变换，却不是等距变换。每个正数 \(c\) 由 \(\sqrt c I_2\) 实现，每个负数 \(c\) 由 \(\sqrt{-c}S\) 实现，故乘子的像为 \(\mathbb R^\times\)。若某个变换 \(g\) 对正定型 \(h\) 有负乘子 \(c\)，则非零 \(v\) 满足 \(h(gv,gv)=c\cdot h(v,v)<0\)，与正定性矛盾。
\end{solution}

\begin{optionalexercise}\label{b2:ex:10-clifford-line}
求实一维二次型的三种 Clifford 代数模型 \(\mathbb R\times\mathbb R\)、\(\mathbb C\)、\(D=\mathbb R[\varepsilon]/(\varepsilon^2)\) 的全部含幺 \(\mathbb R\) 代数自同构。分别求出三个代数中的幂等元，说明它们作为实向量空间维数相同却不能同构。
\end{optionalexercise}
\begin{solution}
\cref{b2:prop:clifford-line-models}给出这三个模型。乘积代数的幂等元为 \((0,0),(1,0),(0,1),(1,1)\)。自同构固定零与单位，故必须固定或交换两个非平凡幂等元；它们又线性生成该代数，所以自同构只有恒等与交换两因子。\(\mathbb C\) 中 \(i\) 的像须满足平方为负一，只能为 \(i\) 或 \(-i\)，两者分别给出恒等与复共轭。

在 \(D\) 中，\(1,\varepsilon\) 为实基。若 \((a+b\varepsilon)^2=0\)，常数项给出 \(a=0\)，故 \(\varepsilon\) 在任何自同构下必变为 \(b\varepsilon\)。可逆性等价于 \(b\ne0\)，而每个这样的缩放确实给出自同构，因此 \(\operatorname{Aut}_{\mathbb R}(D)\cong\mathbb R^\times\)。

\(\mathbb C\) 中幂等元只有零与一。\(D\) 中的幂等条件为 \(a^2=a\)、\((2a-1)b=0\)，同样只有零与一。非平凡幂等元将乘积代数与后两者区分；后两者又由非零幂零元的存在性区分：\(\varepsilon\ne0\) 但 \(\varepsilon^2=0\)，而域 \(\mathbb C\) 没有非零幂零元。所以三者虽都是二维实代数，却两两不同构。
\end{solution}

\begin{optionalexercise}\label{b2:ex:10-clifford-char-two}
在 \(\mathbb F_2^2\) 上取 \(q(x,y)=x^2+y^2\)。证明其 Clifford 代数同构于
\[
\mathbb F_2[u,v]/(u^2,v^2).
\]
直接给出该代数的一个向量空间基，不引用一般 Clifford 代数的维数结论。
\end{optionalexercise}
\begin{solution}
令两个标准基的像为 \(e,f\)，则 \(e^2=f^2=1\)。极化型为零，由\eqref{b2:eq:clifford-polar-relation}，\(ef+fe=0\)，在特征二中即 \(ef=fe\)。置 \(u=e+1,v=f+1\)，得到 \(u^2=v^2=0\)，且它们交换，因此从题中商代数有一个满同态到 Clifford 代数。

反过来，将 \((x,y)\) 送到 \(x(1+u)+y(1+v)\)。因交换性与特征二，这个元素的平方为 \(x^2+y^2\)，由泛性质得到反向同态。两复合固定各生成元，故为同构。多项式按次数分别对 \(u^2,v^2\) 取余，唯一得到 \(a+bu+cv+duv\)：理想 \((u^2,v^2)\) 中每个单项式至少有一个指数不小于二，不能消去这四个单项式。因此 \(1,u,v,uv\) 是基，维数为四。
\end{solution}

